% Consolidation des acquis de la dissertation philosophique — Philosophie 1re Tech — Cours signé OnBuch+
% Programme officiel MINESEC. Compiler avec : tectonic N02-consolidation-acquis-dissertation-philosophique.tex
\newcommand{\DOCMATIERE}{Philosophie}
\newcommand{\DOCNIVEAU}{1re Tech}
\newcommand{\DOCMODULE}{Philosophie – classe de Première technique}
\newcommand{\DOCLECON}{N02}
\newcommand{\DOCTITRE}{Consolidation des acquis de la dissertation philosophique}
% OnBuch+ — préambule « cours signés » ENRICHI (compilé avec tectonic / XeLaTeX)
% Système visuel « Pop » d'OnBuch+ : fond crème (jamais blanc), bordures encre
% épaisses, ombres « dures » décalées, titres Archivo Black en capitales.
% Version MATHÉMATIQUES : graphes (pgfplots/tikz), boîtes pédagogiques
% (définition, propriété, méthode, attention, le savais-tu, exercice résolu,
% exercice, TP, bilan), page de garde et signature OnBuch+ ancrée en bas.
%
% Macros attendues dans main.tex AVANT \input{preamble} :
%   \DOCMATIERE  \DOCNIVEAU  \DOCMODULE  \DOCLECON (numéro)  \DOCTITRE
\documentclass[11pt]{article}

\usepackage{fontspec}
\usepackage[french]{babel}
\usepackage[a4paper,margin=2cm,top=3.4cm,bottom=2.8cm,headheight=1.4cm,footskip=1.3cm]{geometry}
\usepackage{amsmath,amssymb,amsfonts}
\usepackage{mathtools} % \xmapsto, \coloneqq... souvent utilisés par les modèles
\usepackage{cancel} % simplifications barrées (bilans de forces)
% \abs{z} (module, valeur absolue) et \norm{u} (norme), utilisés par les modèles
\providecommand{\abs}{}\let\abs\relax\DeclarePairedDelimiter{\abs}{\lvert}{\rvert}
\providecommand{\norm}{}\let\norm\relax\DeclarePairedDelimiter{\norm}{\lVert}{\rVert}
\usepackage[table]{xcolor} % [table] : fournit aussi \rowcolors (alternance de lignes)
\usepackage{pagecolor}
\usepackage{tikz}
\usetikzlibrary{arrows.meta,positioning,calc,shapes.geometric,shapes.misc,shapes.arrows,decorations.pathmorphing,decorations.pathreplacing,decorations.markings,patterns,shadows,mindmap,trees,fit,backgrounds,matrix,intersections,angles,quotes}
\usepackage{pgfplots}
\usepackage[version=4]{mhchem} % équations nucléaires \ce{^{235}_{92}U}
\usepackage[european,siunitx]{circuitikz} % schémas électriques
\pgfplotsset{compat=1.18}
\usepgfplotslibrary{fillbetween,groupplots} % aires entre courbes (calcul intégral)
\usepackage{enumitem}
\usepackage{fancyhdr}
% [most] charge toutes les bibliothèques tcolorbox (listings, minted, raster,
% documentation...) et gonfle inutilement la mémoire TeX sur un document long
% (~300 boîtes) : on ne charge que ce qui est réellement utilisé.
\usepackage[skins,breakable]{tcolorbox}
\usepackage{graphicx}
\usepackage{colortbl}
\usepackage{booktabs}
\usepackage{tabularx}
\usepackage{multirow}
\usepackage{diagbox} % case d'en-tête diagonale des tableaux à double entrée
% Colonnes extensibles centrée / gauche / droite (C, L, R), souvent utilisées par les modèles
\newcolumntype{C}{>{\centering\arraybackslash}X}
\newcolumntype{L}{>{\raggedright\arraybackslash}X}
\newcolumntype{R}{>{\raggedleft\arraybackslash}X}
\usepackage{array}
\usepackage{multicol}
\usepackage{siunitx}
% parse-numbers=false : accepte aussi \SI{2,5 \times 10^{-3}}{...}
\sisetup{locale=FR,output-decimal-marker={,},group-minimum-digits=4,parse-numbers=false}
\DeclareSIUnit\ohm{\ensuremath{\symup{\Omega}}} % U+2126 absent de la police
\DeclareSIUnit\division{div} % réglages d'oscilloscope (ms/div, V/div)
\DeclareSIUnit\tr{tr} % tours (vitesse de rotation, ex. \SI{1500}{\tr\per\minute})
\DeclareSIUnit\molar{M} % concentration molaire (ex. \SI{3}{\molar}, \SI{1}{\micro\molar})
\DeclareSIUnit\an{an} % année (le modèle confond parfois avec \year, un compteur TeX) — ex. \SI{5}{\kilo\meter\per\an}
\DeclareSIUnit\FCFA{FCFA} % franc CFA (ex. \SI{500}{\FCFA\per\kilo\gram})
\DeclareSIUnit\degreeBrix{\textdegree Brix} % teneur en sucre d'un sirop/jus (réfractométrie)
\DeclareSIUnit\month{mois} % le modèle utilise parfois \month, un compteur TeX, comme unité de durée
\DeclareSIUnit\microliter{\micro\liter} % le modèle écrit parfois \microliter en un seul token
\DeclareSIUnit\million{millions} % dénombrement (ex. \SI{15}{\million\per\mL} spermatozoïdes)
\usepackage{qrcode}
% \degree : commande fréquemment utilisée par les modèles pour un angle ou une
% température en dehors de \SI{}{} (non fournie par les paquets chargés).
\providecommand{\degree}{^{\circ}}
% \qed (fin de démonstration), utilisé par les modèles sans amsthm
\providecommand{\qed}{\hfill$\square$}
% \barre : double barre des valeurs interdites dans un tableau de variations (tkz-tab)
\providecommand{\barre}{\ensuremath{\|}}
% environnement proof (amsthm non chargé) utilisé par les modèles
\ifdefined\proof\else\newenvironment{proof}[1][Démonstration]{\par\noindent\textit{#1.}\ }{\qed\par}\fi
\usepackage{lastpage}
\usepackage{hyperref}
\hypersetup{hidelinks}

% ============================================================
% Palette « Pop » OnBuch+ — mêmes hex que la classe P (plus_ui.dart)
% ============================================================
\definecolor{poporange}{HTML}{F59321}
\definecolor{popdark}{HTML}{E07A0C}
\definecolor{popcream}{HTML}{FFF6E8}
\definecolor{popink}{HTML}{1C1714}
\definecolor{poppurple}{HTML}{7A5AE0}
\definecolor{popgreen}{HTML}{2E9E6A}
\definecolor{poppink}{HTML}{E0457B}
\definecolor{popblue}{HTML}{2F7DE1}
\definecolor{popgold}{HTML}{C98A12}
\definecolor{popmuted}{HTML}{8A7F76}
\definecolor{popline}{HTML}{EADFD2}
\definecolor{popgray}{HTML}{8A7F76}
\colorlet{popgrey}{popgray}
% Alias tolérants (le modèle nomme parfois les couleurs différemment)
\colorlet{popwhite}{white}
\colorlet{popblack}{popink}
\colorlet{popred}{poppink}
\colorlet{popyellow}{popgold}
% Teintes claires (fonds de tableaux/encadrés) — le modèle nomme parfois
% les couleurs avec un suffixe L (ex. popblueL) sans les avoir définies.
\colorlet{poporangeL}{poporange!20}
\colorlet{popdarkL}{popdark!20}
\colorlet{poppurpleL}{poppurple!20}
\colorlet{popgreenL}{popgreen!20}
\colorlet{poppinkL}{poppink!20}
\colorlet{popblueL}{popblue!20}
\colorlet{popgoldL}{popgold!20}
\colorlet{popredL}{popred!20}
\colorlet{popgrayL}{popgray!20}
% Teintes claires pour fonds de tableaux / zones
\colorlet{popblueL}{popblue!10!white}
\colorlet{poporangeL}{poporange!14!white}
\colorlet{popgreenL}{popgreen!12!white}
\colorlet{poppurpleL}{poppurple!12!white}
\colorlet{poppinkL}{poppink!12!white}
\colorlet{popgoldL}{popgold!14!white}

\pagecolor{popcream}

% ============================================================
% Typographie
% ============================================================
\defaultfontfeatures{Ligatures=TeX}
\setmainfont{PlusJakartaSans-Medium.ttf}[Path=../../../fonts/,BoldFont=PlusJakartaSans-Bold.ttf]
% Maths en sans-serif (Fira Math) avec chiffres et lettres droites de Plus
% Jakarta, pour que les formules se fondent dans le texte.
\usepackage[math-style=ISO,bold-style=ISO]{unicode-math}
\setmathfont{FiraMath-Regular.otf}[Path=../../../fonts/]
\setmathfont{PlusJakartaSans-Medium.ttf}[Path=../../../fonts/,range={up/num,up/{Latin,latin},\mathup}]
\usepackage{icomma} % virgule décimale sans espace en mode math
% Caractères mathématiques Unicode absents des polices de texte (Plus Jakarta,
% Archivo Black) : rendus par leur équivalent mathématique, en texte comme en
% formule — sinon ils s'affichent en carré vide (ex. « Arithmétique dans ℤ »).
\usepackage{newunicodechar}
\newunicodechar{ℝ}{\ensuremath{\mathbb{R}}}
\newunicodechar{ℤ}{\ensuremath{\mathbb{Z}}}
\newunicodechar{ℂ}{\ensuremath{\mathbb{C}}}
\newunicodechar{ℕ}{\ensuremath{\mathbb{N}}}
\newunicodechar{ℚ}{\ensuremath{\mathbb{Q}}}
\newunicodechar{𝔻}{\ensuremath{\mathbb{D}}}
\newunicodechar{α}{\ensuremath{\alpha}}
\newunicodechar{β}{\ensuremath{\beta}}
\newunicodechar{γ}{\ensuremath{\gamma}}
\newunicodechar{δ}{\ensuremath{\delta}}
\newunicodechar{ε}{\ensuremath{\varepsilon}}
\newunicodechar{θ}{\ensuremath{\theta}}
\newunicodechar{λ}{\ensuremath{\lambda}}
\newunicodechar{μ}{\ensuremath{\mu}}
\newunicodechar{σ}{\ensuremath{\sigma}}
\newunicodechar{τ}{\ensuremath{\tau}}
\newunicodechar{φ}{\ensuremath{\varphi}}
\newunicodechar{ψ}{\ensuremath{\psi}}
\newunicodechar{ω}{\ensuremath{\omega}}
\newunicodechar{Σ}{\ensuremath{\Sigma}}
% majuscules produites par \MakeUppercase dans les titres (α-aminés -> Α)
\newunicodechar{Α}{\ensuremath{\alpha}}
\newunicodechar{Β}{\ensuremath{\beta}}
\newunicodechar{Γ}{\ensuremath{\Gamma}}
\newunicodechar{Δ}{\ensuremath{\Delta}}
\newunicodechar{Ε}{\ensuremath{\varepsilon}}
\newunicodechar{Θ}{\ensuremath{\Theta}}
\newunicodechar{Λ}{\ensuremath{\Lambda}}
\newunicodechar{Μ}{\ensuremath{\mu}}
\newunicodechar{Τ}{\ensuremath{\tau}}
\newunicodechar{Φ}{\ensuremath{\Phi}}
\newunicodechar{Ψ}{\ensuremath{\Psi}}
\newunicodechar{Ω}{\ensuremath{\Omega}}
\newunicodechar{↦}{\ensuremath{\mapsto}}
\newunicodechar{∈}{\ensuremath{\in}}
\newunicodechar{∉}{\ensuremath{\notin}}
\newunicodechar{∧}{\ensuremath{\wedge}}
\newunicodechar{⇔}{\ensuremath{\Leftrightarrow}}
\newunicodechar{⟺}{\ensuremath{\Longleftrightarrow}}
\newunicodechar{⇒}{\ensuremath{\Rightarrow}}
\newunicodechar{⟹}{\ensuremath{\Longrightarrow}}
\newunicodechar{∘}{\ensuremath{\circ}}
\newunicodechar{∪}{\ensuremath{\cup}}
\newunicodechar{∩}{\ensuremath{\cap}}
\newunicodechar{≡}{\ensuremath{\equiv}}
\newunicodechar{⊂}{\ensuremath{\subset}}
\newunicodechar{⊥}{\ensuremath{\perp}}
\newunicodechar{∃}{\ensuremath{\exists}}
\newunicodechar{∀}{\ensuremath{\forall}}
\newunicodechar{∛}{\ensuremath{\sqrt[3]{\,}}}
\newunicodechar{∜}{\ensuremath{\sqrt[4]{\,}}}
\newunicodechar{⃗}{\ensuremath{{}^{\to}}}
\newunicodechar{ⁿ}{\ifmmode^{n}\else\textsuperscript{\ensuremath{n}}\fi}
\newunicodechar{ˣ}{\ifmmode^{x}\else\textsuperscript{\ensuremath{x}}\fi}
\newunicodechar{ᵇ}{\ifmmode^{b}\else\textsuperscript{\ensuremath{b}}\fi}
\newunicodechar{ʳ}{\ifmmode^{r}\else\textsuperscript{\ensuremath{r}}\fi}
\newunicodechar{ᵘ}{\ifmmode^{u}\else\textsuperscript{\ensuremath{u}}\fi}
\newunicodechar{ʸ}{\ifmmode^{y}\else\textsuperscript{\ensuremath{y}}\fi}
\newunicodechar{ᵃ}{\ifmmode^{a}\else\textsuperscript{\ensuremath{a}}\fi}
\newunicodechar{ᵉ}{\ifmmode^{e}\else\textsuperscript{\ensuremath{e}}\fi}
\newunicodechar{ᵏ}{\ifmmode^{k}\else\textsuperscript{\ensuremath{k}}\fi}
\newunicodechar{ᵖ}{\ifmmode^{p}\else\textsuperscript{\ensuremath{p}}\fi}
\newunicodechar{ᴺ}{\ifmmode^{N}\else\textsuperscript{\ensuremath{N}}\fi}
\newunicodechar{ᶜ}{\ifmmode^{c}\else\textsuperscript{\ensuremath{c}}\fi}
\newunicodechar{ʰ}{\ifmmode^{h}\else\textsuperscript{\ensuremath{h}}\fi}
\newunicodechar{ᵠ}{\ifmmode^{q}\else\textsuperscript{\ensuremath{q}}\fi}
\newunicodechar{ₙ}{\ifmmode_{n}\else\textsubscript{\ensuremath{n}}\fi}
\newunicodechar{ₐ}{\ifmmode_{a}\else\textsubscript{\ensuremath{a}}\fi}
\newunicodechar{ᵢ}{\ifmmode_{i}\else\textsubscript{\ensuremath{i}}\fi}
\newunicodechar{ᵣ}{\ifmmode_{r}\else\textsubscript{\ensuremath{r}}\fi}
\newunicodechar{ₖ}{\ifmmode_{k}\else\textsubscript{\ensuremath{k}}\fi}
\newunicodechar{ᵦ}{\ifmmode_{\beta}\else\textsubscript{\ensuremath{\beta}}\fi}
\newunicodechar{ₓ}{\ifmmode_{x}\else\textsubscript{\ensuremath{x}}\fi}
\newfontfamily\popdisplay{ArchivoBlack-Regular.ttf}[Path=../../../fonts/]
\newfontfamily\poplogo{SpaceGrotesk-Bold.ttf}[Path=../../../fonts/]
\newfontfamily\popsemi{PlusJakartaSans-SemiBold.ttf}[Path=../../../fonts/]
\newfontfamily\popxbold{PlusJakartaSans-ExtraBold.ttf}[Path=../../../fonts/]
\newfontfamily\popxboldls{PlusJakartaSans-ExtraBold.ttf}[Path=../../../fonts/,LetterSpace=3.0]

\setlist[enumerate]{leftmargin=1.7em,itemsep=5pt,topsep=4pt}
\setlist[itemize]{leftmargin=1.7em,itemsep=4pt,topsep=4pt,label={\color{poporange}\textbullet}}
\setlength{\parindent}{0pt}
\setlength{\parskip}{5pt}
\renewcommand{\arraystretch}{1.25}

% pgfplots : style Pop par défaut
\pgfplotsset{
  every axis/.append style={axis line style={popink,line width=1pt},tick style={popink},
    grid style={popline,line width=0.5pt},label style={font=\small},tick label style={font=\footnotesize}},
  popcurve/.style={poporange,line width=1.8pt,no markers,smooth},
}
% Même style disponible dans un \draw TikZ simple (hors pgfplots)
\tikzset{popcurve/.style={poporange,line width=1.8pt}}
\tikzset{popgrid/.style={popline,very thin}}
\colorlet{popcurve}{poporange} % le nom du style est parfois utilisé comme couleur

% ============================================================
% Pill / squiggle
% ============================================================
\newcommand{\poppill}[3]{%
  \tikz[baseline=-0.55ex]\node[draw=popink,line width=1.2pt,rounded corners=2.6pt,%
    fill=#2,inner xsep=6.5pt,inner ysep=3pt]{%
    \popxboldls\fontsize{7}{7}\selectfont\color{#3}\MakeUppercase{#1}};%
}
\newcommand{\popsquiggle}[1]{%
  \tikz[baseline=-0.4ex]{
    \draw[#1,line width=2.6pt,line cap=round]
      plot [smooth] coordinates {(0,0.09) (0.34,0.01) (0.68,0.21) (1.02,0.01) (1.36,0.10)};
  }%
}
% Mot-clé surligné (marqueur orange sous le texte)
\newcommand{\cle}[1]{{\popxbold\color{popdark}#1}}

% ============================================================
% En-tête / pied
% ============================================================
\pagestyle{fancy}
\fancyhf{}
\renewcommand{\headrulewidth}{0pt}
\renewcommand{\footrulewidth}{0pt}
\lhead{\raisebox{-0.15ex}{\poplogo\fontsize{13}{13}\selectfont\color{popink}OnBuch\color{poporange}+}}
\rhead{\poppill{\DOCMATIERE\ \textbullet\ \DOCNIVEAU\ \textbullet\ Leçon \DOCLECON}{white}{popink}}
\lfoot{\popxbold\fontsize{7.6}{7.6}\selectfont\color{popmuted}\MakeUppercase{Cours signé OnBuch+}\ \textbullet\ {\popsemi\DOCTITRE}}
\rfoot{\tikz[baseline=-0.6ex]\node[rounded corners=8.5pt,fill=popink,minimum height=17pt,minimum width=17pt,inner xsep=5pt,inner sep=0]{\popxbold\fontsize{8}{8}\selectfont\color{popcream}\thepage\,/\,\pageref*{LastPage}};}

% ============================================================
% Titres
% ============================================================
\newcommand{\chaptitre}[2]{%
  \begin{tcolorbox}[enhanced,colback=poporange,colframe=popink,boxrule=2.6pt,arc=11pt,
      drop shadow=popink,left=15pt,right=15pt,top=12pt,bottom=11pt,before skip=3pt,after skip=18pt]
    \poppill{#2}{popink}{popcream}\par\vspace{9pt}
    {\raggedright\hyphenpenalty=10000\exhyphenpenalty=10000\popdisplay\fontsize{22}{24}\selectfont\color{popcream}\MakeUppercase{#1}\par}
  \end{tcolorbox}
}
% Section numérotée : pastille orange avec le numéro + titre Archivo
\newcounter{popsec}
\newcommand{\coursec}[1]{%
  \stepcounter{popsec}\setcounter{popsub}{0}%
  \par\vspace{16pt}\noindent
  \begin{minipage}{\linewidth}
  \tikz[baseline=-0.7ex]\node[draw=popink,line width=1.6pt,fill=poporange,rounded corners=4pt,minimum size=22pt,inner sep=0]{\popdisplay\fontsize{12}{12}\selectfont\color{popcream}\thepopsec};%
  \hspace{8pt}{\popdisplay\fontsize{15}{17}\selectfont\color{popink}\MakeUppercase{#1}}\par
  \vspace{4pt}\noindent{\color{poporange}\rule{36pt}{3.6pt}}
  \end{minipage}\par\nopagebreak\vspace{8pt}
}
\newcounter{popsub}
\newcommand{\courssub}[1]{%
  \stepcounter{popsub}%
  \par\vspace{9pt}\noindent{\popxbold\fontsize{11.5}{13}\selectfont\color{popink}{\color{poporange}\thepopsec.\thepopsub}\hspace{6pt}#1}\par\nopagebreak\vspace{4pt}
}

% ============================================================
% Boîtes pédagogiques — même recette : fond blanc, bordure épaisse colorée,
% ombre dure encre, badge plein. Titre optionnel [#1].
% ============================================================
\tcbset{popbase/.style={breakable,enhanced,colback=white,boxrule=2.3pt,arc=9pt,
  shadow={3pt}{-3pt}{0pt}{popink},left=14pt,right=14pt,top=11pt,bottom=11pt,before skip=11pt,after skip=15pt,
  fontupper=\popsemi\fontsize{10.6}{14.5}\selectfont\color{popink},
  fontlower=\popsemi\fontsize{10.6}{14.5}\selectfont\color{popink}}}
% Coche dessinée (le glyphe ✓ n'existe pas dans les polices du document)
\newcommand{\popcheck}[1]{\tikz[baseline=-0.5ex]\draw[#1,line width=1.6pt,line cap=round,line join=round](-0.13,0.0)--(-0.04,-0.09)--(0.13,0.1);}
\providecommand{\coche}{\popcheck{popgreen}}
\providecommand{\resultat}[1]{\par\smallskip\noindent\fcolorbox{popgreen}{popgreenL}{\parbox{\dimexpr\linewidth-2\fboxsep-2\fboxrule\relax}{\textbf{Résultat.} #1}}\par\smallskip}
\newcommand{\popboxhead}[3]{\poppill{#1}{#2}{white}\ifx\relax#3\relax\else\hspace{6pt}{\popxbold\fontsize{10.6}{12}\selectfont\color{popink}#3}\fi\par\vspace{7pt}}

\newtcolorbox{aretenir}[1][]{popbase,colframe=popgreen,before upper={\popboxhead{À retenir}{popgreen}{#1}}}
\newtcolorbox{exemplebox}[1][]{popbase,colframe=poporange,before upper={\popboxhead{Exemple}{poporange}{#1}}}
\newtcolorbox{definition}[1][]{popbase,colframe=popblue,before upper={\popboxhead{Définition}{popblue}{#1}}}
\newtcolorbox{propriete}[1][]{popbase,colframe=poppurple,before upper={\popboxhead{Propriété}{poppurple}{#1}}}
\newtcolorbox{methode}[1][]{popbase,colframe=popgold,before upper={\popboxhead{Méthode}{popgold}{#1}}}
\newtcolorbox{attention}[1][]{popbase,colframe=poppink,before upper={\popboxhead{Attention}{poppink}{#1}}}
\newtcolorbox{experience}[1][]{popbase,colframe=popblue,colback=popblueL,before upper={\popboxhead{Expérience}{popblue}{#1}}}
% « Le savais-tu ? » : carte violette pleine (contexte, culture, Cameroun)
\newtcolorbox{savaistu}[1][]{popbase,colback=poppurple,colframe=popink,
  fontupper=\popsemi\fontsize{10.4}{14.2}\selectfont\color{white},
  before upper={\poppill{Le savais-tu\,?}{white}{poppurple}\ifx\relax#1\relax\else\hspace{6pt}{\popxbold\color{white}#1}\fi\par\vspace{7pt}}}
% Exercice résolu : énoncé en haut, \tcblower puis la solution sur fond crème
\newcounter{popexo}\newcounter{popexr}
\newtcolorbox{exoresolu}[1][]{popbase,colframe=poporange,colbacklower=popcream,
  segmentation style={popink,line width=1.2pt,dashed},
  before upper={\stepcounter{popexr}\popboxhead{Exercice résolu \thepopexr}{poporange}{#1}},
  before lower={\poppill{Solution}{popink}{popcream}\par\vspace{6pt}}}
% Exercice (non résolu en place) — la correction est dans la partie Corrigés
\newtcolorbox{exercice}[1][]{popbase,colframe=popink,
  before upper={\stepcounter{popexo}\popboxhead{Exercice \thepopexo}{popink}{#1}}}
% Corrigé
\newtcolorbox{corrige}[1][]{popbase,colframe=popgreen,colback=popgreenL,
  before upper={\popboxhead{Corrigé}{popgreen}{#1}}}
% Objectifs / prérequis / bilan
\newtcolorbox{objectifs}{popbase,colframe=popink,colback=poporangeL,
  before upper={\popboxhead{Objectifs de la leçon}{popink}{}}}
\newtcolorbox{prerequis}{popbase,colframe=popink,colback=popblueL,
  before upper={\popboxhead{Prérequis}{popblue}{}}}
\newtcolorbox{bilan}{popbase,colframe=popink,colback=white,boxrule=2.8pt,
  before upper={\popboxhead{Fiche bilan}{poporange}{L'essentiel à connaître pour l'examen}}}
% Figure encadrée avec légende
\newtcolorbox{popfigure}[1][]{enhanced,breakable=false,colback=white,colframe=popline,boxrule=1.4pt,arc=8pt,
  left=10pt,right=10pt,top=10pt,bottom=8pt,before skip=10pt,after skip=12pt,halign=center,
  lower separated=false,halign lower=center,
  fontlower=\popsemi\fontsize{9}{11.5}\selectfont\color{popmuted},
  #1}
\newcommand{\legende}[1]{\par\vspace{6pt}{\popsemi\fontsize{9}{11.5}\selectfont\color{popmuted}#1}}

% ============================================================
% Signature OnBuch+
% ============================================================
\newcommand{\poplogoinline}[1]{{\poplogo\fontsize{#1}{#1}\selectfont\color{popink}OnBuch\color{poporange}+}}
\newcommand{\signatureonbuch}{%
  \begin{tcolorbox}[enhanced,colback=popink,colframe=popink,boxrule=2.2pt,arc=10pt,
      drop shadow=poporange,left=14pt,right=14pt,top=10pt,bottom=10pt,before skip=0pt,after skip=0pt]
    \tikz[baseline=-0.7ex]\node[circle,fill=popgreen,draw=popcream,line width=1.3pt,minimum size=22pt,inner sep=0]{\popcheck{white}};%
    \hspace{9pt}\begin{minipage}[c]{0.62\linewidth}
      {\popxbold\fontsize{12}{13}\selectfont\color{popcream}Signé\ \textbullet\ {\poplogo OnBuch\color{poporange}+}}\par\vspace{2pt}
      {\raggedright\popsemi\fontsize{8}{10.5}\selectfont\color{popline}\MakeUppercase{Équipe pédagogique OnBuch+}\\\MakeUppercase{Conforme au programme officiel MINESEC}\par}
    \end{minipage}\hfill
    {\poplogo\fontsize{22}{22}\selectfont\color{popcream}OnBuch\color{poporange}+}
  \end{tcolorbox}%
}

% ============================================================
% Page de garde
% #1 titre · #2 matière · #3 niveau · #4 module · #5 n° leçon · #6 durée ·
% #7 description / objectifs (texte ou itemize)
% ============================================================
\newcommand{\pagegarde}[7]{%
  \thispagestyle{empty}
  \newgeometry{a4paper,margin=1.7cm,top=1.6cm,bottom=1.4cm}%
  \begin{tikzpicture}[remember picture,overlay]
    \fill[poporange!18] ([xshift=-3.2cm,yshift=-3.6cm]current page.north east) circle (4.6cm);
    \fill[poppurple!14] ([xshift=2.4cm,yshift=7.6cm]current page.south west) circle (3.8cm);
  \end{tikzpicture}%
  {\poplogo\fontsize{30}{30}\selectfont\color{popink}OnBuch\color{poporange}+}\hfill
  \raisebox{0.6ex}{\poppill{Cours signé \textbullet\ Premium}{popink}{popcream}}\par
  \vspace{1pt}\popsquiggle{poppurple}\par
  \vspace{8pt}
  \begin{tcolorbox}[enhanced,colback=poporange,colframe=popink,boxrule=2.8pt,arc=13pt,
      drop shadow=popink,left=18pt,right=18pt,top=12pt,bottom=13pt,after skip=10pt]
    \poppill{#2 \textbullet\ #3}{popink}{popcream}\hspace{5pt}\poppill{#4}{white}{popink}\par\vspace{8pt}
    {\popdisplay\fontsize{14}{14}\selectfont\color{popink}LEÇON #5}\par\vspace{4pt}
    {\raggedright\hyphenpenalty=10000\exhyphenpenalty=10000\popdisplay\fontsize{25}{27}\selectfont\color{popcream}\MakeUppercase{#1}\par}
    \vspace{8pt}
    {\popxbold\fontsize{9.5}{9.5}\selectfont\color{popink}\MakeUppercase{Durée conseillée : #6}}
  \end{tcolorbox}
  \begin{tcolorbox}[enhanced,colback=white,colframe=popink,boxrule=2.2pt,arc=10pt,
      drop shadow=popink,left=16pt,right=16pt,top=10pt,bottom=10pt,after skip=8pt]
    \poppill{Dans cette leçon}{poppurple}{white}\par\vspace{5pt}
    \popsemi\fontsize{9.3}{12.6}\selectfont\color{popink}#7
  \end{tcolorbox}
  \noindent\begin{minipage}[t]{0.487\linewidth}
    \begin{tcolorbox}[enhanced,colback=popgreen,colframe=popink,boxrule=2.2pt,arc=10pt,
        drop shadow=popink,left=11pt,right=11pt,top=10pt,bottom=10pt,height=3.1cm,valign=center]
      \tcbox[colback=white,colframe=popink,boxrule=1.3pt,arc=5pt,boxsep=0pt,nobeforeafter,
        left=3pt,right=3pt,top=3pt,bottom=3pt,valign=center]{\qrcode[height=1.9cm]{https://play.google.com/store/apps/details?id=com.luvvix.onbuch&hl=fr}}%
      \hspace{8pt}\begin{minipage}[c]{0.5\linewidth}
        {\popdisplay\fontsize{11}{12.5}\selectfont\color{white}\MakeUppercase{Télécharge OnBuch}}\par\vspace{4pt}
        {\raggedright\popsemi\fontsize{8.4}{11}\selectfont\color{white}Gratuit \textbullet\ Google Play\\Tuteur IA, annales, concours, résultats\par}
      \end{minipage}
    \end{tcolorbox}
  \end{minipage}\hfill
  \begin{minipage}[t]{0.487\linewidth}
    \begin{tcolorbox}[enhanced,colback=poppurple,colframe=popink,boxrule=2.2pt,arc=10pt,
        drop shadow=popink,left=11pt,right=11pt,top=10pt,bottom=10pt,height=3.1cm,valign=center]
      \tikz[baseline=-0.7ex]\node[circle,fill=white,draw=popink,line width=1.3pt,minimum size=1.6cm,inner sep=0]{\popdisplay\fontsize{24}{24}\selectfont\color{poppurple}+};%
      \hspace{8pt}\begin{minipage}[c]{0.6\linewidth}
        {\popdisplay\fontsize{11}{12.5}\selectfont\color{white}\MakeUppercase{Passe à OnBuch+}}\par\vspace{4pt}
        {\raggedright\popsemi\fontsize{8.4}{11}\selectfont\color{white}Tous les cours signés, Tuteur IA illimité, fiches PDF — onglet OnBuch+ de l'app\par}
      \end{minipage}
    \end{tcolorbox}
  \end{minipage}
  \vspace{8pt}
  \popcontenu
  \vfill
  \signatureonbuch
  \restoregeometry
  \newpage
}

% Rangée « Ce cours contient » (page de garde)
\newcommand{\poptile}[3]{% #1 couleur #2 gros texte #3 légende
  \begin{tcolorbox}[enhanced,colback=#1,colframe=popink,boxrule=2pt,arc=9pt,shadow={2.5pt}{-2.5pt}{0pt}{popink},
      left=6pt,right=6pt,top=9pt,bottom=9pt,halign=center,valign=center,height=1.75cm,width=\linewidth,nobeforeafter]
    {\popdisplay\fontsize{12.5}{13}\selectfont\color{white}\MakeUppercase{#2}}\par\vspace{3pt}
    {\centering\popsemi\fontsize{7.8}{9.5}\selectfont\color{white}#3\par}
  \end{tcolorbox}}
\newcommand{\popcontenu}{%
  \par\noindent{\popxboldls\fontsize{7.5}{7.5}\selectfont\color{popmuted}CE COURS SIGNÉ CONTIENT}\par\vspace{7pt}
  \noindent\begin{minipage}[t]{0.235\linewidth}\poptile{popblue}{Cours}{complet, illustré, pas à pas}\end{minipage}\hfill
  \begin{minipage}[t]{0.235\linewidth}\poptile{poporange}{Méthodes}{et exercices résolus}\end{minipage}\hfill
  \begin{minipage}[t]{0.235\linewidth}\poptile{poppink}{Exercices}{type examen + corrigés}\end{minipage}\hfill
  \begin{minipage}[t]{0.235\linewidth}\poptile{popgreen}{Fiche bilan}{l'essentiel à retenir}\end{minipage}\par}

% Sommaire
\newtcolorbox{sommaire}{enhanced,colback=white,colframe=popink,boxrule=2.2pt,arc=10pt,
  drop shadow=popink,left=18pt,right=18pt,top=13pt,bottom=13pt,before skip=6pt,after skip=20pt,
  before upper={\poppill{Sommaire}{poppurple}{white}\par\vspace{9pt}\popsemi\fontsize{10.6}{14.5}\selectfont\color{popink}}}

% Signature de fin de cours — poussée tout en bas de la dernière page
\newcommand{\findecours}{%
  \par\vfill\nopagebreak
  \begin{center}\popsquiggle{poporange}\end{center}\vspace{4pt}
  \signatureonbuch
}
\AtBeginDocument{\def\llbracket{\mathopen{[\mkern-3.5mu[}}\def\rrbracket{\mathclose{]\mkern-3.5mu]}}} % intervalles d'entiers (glyphe ⟦ absent de la police)
% Glyphes absents de Fira Math : redéfinis à partir de glyphes disponibles
\AtBeginDocument{%
  \def\setminus{\mathbin{\backslash}}\let\smallsetminus\setminus
  \def\vdots{\mathord{\vbox{\baselineskip3.2pt\lineskiplimit0pt\kern4pt\hbox{.}\hbox{.}\hbox{.}}}}%
  \def\ddots{\mathinner{\mkern1mu\raise6.5pt\hbox{.}\mkern2mu\raise3.6pt\hbox{.}\mkern2mu\raise0.7pt\hbox{.}\mkern1mu}}%
  \def\triangle{\mathord{\scalebox{0.9}{$\symup{\Delta}$}}}%
  \def\nabla{\mathord{\raisebox{\depth}{\scalebox{1}[-1]{$\symup{\Delta}$}}}}%
  \def\checkmark{\popcheck{popdark}}%
  \def\star{\mathbin{\ast}}\def\bigstar{\mathord{\ast}}%
  \def\diamond{\mathbin{\scalebox{0.6}{$\Diamond$}}}%
  \def\bigcup{\mathop{\mathchoice{\raisebox{-0.3em}{\scalebox{1.7}{$\cup$}}}{\scalebox{1.25}{$\cup$}}{\cup}{\cup}}\displaylimits}%
  \def\bigcap{\mathop{\mathchoice{\raisebox{-0.3em}{\scalebox{1.7}{$\cap$}}}{\scalebox{1.25}{$\cap$}}{\cap}{\cap}}\displaylimits}%
  \def\lesseqqgtr{\mathrel{\lessgtr}}\def\bigoplus{\mathop{\oplus}\displaylimits}%
}
\providecommand{\ssi}{si et seulement si} % abréviation fréquente
\AtBeginDocument{\providecommand{\wideparen}{\overparen}} % arc de cercle

%% FIGFIX-BEGIN v3 — correctifs globaux des figures (docs/onbuch-generation/tools/figfix)
\usepackage{adjustbox}\usepackage{etoolbox}
% toute figure TikZ/pgfplots plus large que la ligne est réduite à la largeur disponible
\BeforeBeginEnvironment{tikzpicture}{\begin{adjustbox}{max width=\linewidth}}
\AfterEndEnvironment{tikzpicture}{\end{adjustbox}}
% légendes pgfplots lisibles par-dessus les courbes
\pgfplotsset{every axis/.append style={legend style={fill=white,fill opacity=0.92,text opacity=1,draw=black!15,inner sep=3pt}}}
% étiquettes d'arêtes (syntaxe edge["..."]) avec fond blanc
\tikzset{every edge quotes/.append style={fill=white,inner sep=1.2pt}}
% bulles de cartes mentales : texte à la ligne dans la bulle, bulle un peu plus grande
\tikzset{every concept/.append style={text width=2.0cm,align=center,minimum size=2.3cm,inner sep=2pt}}
%% FIGFIX-END
\begin{document}

\pagegarde{Consolidation des acquis de la dissertation philosophique}{Philosophie}{1re Tech}{Philosophie – classe de Première technique}{N02}{2 séances (fiche de progression harmonisée)}{Cette leçon consolide la maîtrise de la méthode de la dissertation philosophique par l'analyse de sujets ancrés dans le réel camerounais et la rédaction guidée de copies types. Elle vise l'autonomie de l'élève face à l'épreuve du Baccalauréat technique.
\vspace{4pt}
\begin{multicols}{2}\small
\begin{itemize}
\item À la fin de cette leçon, je sais analyser un sujet de dissertation en identifiant ses termes clés, son type et ses présupposés.
\item À la fin de cette leçon, je sais formuler une problématique pertinente qui ouvre le débat philosophique sans être hors-sujet.
\item À la fin de cette leçon, je sais construire un plan dialectique, analytique ou thématique cohérent et progressif.
\item À la fin de cette leçon, je sais rédiger une introduction en entonnoir (accroche, définition, problématique, annonce du plan) et une conclusion synthétique et ouverte.
\item À la fin de cette leçon, je sais argumenter chaque sous-partie selon le schéma : thèse -> argument -> exemple concret (Cameroun/Monde) -> lien avec la problématique.
\item À la fin de cette leçon, je sais mobiliser des références philosophiques (auteurs, concepts) et des faits de société camerounais pour étayer ma démonstration.
\end{itemize}\end{multicols}}

\begin{sommaire}
\begin{multicols}{2}
\begin{enumerate}
\item 1. L'analyse rigoureuse du sujet : la clé de voûte
\item 2. L'architecture du développement : Plans et articulation
\item 3. L'art de l'introduction et de la conclusion : les vitrines de la copie
\item 4. L'argumentation philosophique : Preuves, exemples et style
\item 5. Atelier de consolidation : Sujets types Bac \& Grille d'auto-évaluation
\item 6. Transfert : La dissertation hors examen (Vie professionnelle \& Citoyenne)
\item Activité d'intégration
\item Méthodes et exercices résolus
\item Exercices
\item Corrigés des exercices
\item Fiche bilan
\end{enumerate}
\end{multicols}
\end{sommaire}

\begin{objectifs}
\begin{itemize}
\item À la fin de cette leçon, je sais analyser un sujet de dissertation en identifiant ses termes clés, son type et ses présupposés.
\item À la fin de cette leçon, je sais formuler une problématique pertinente qui ouvre le débat philosophique sans être hors-sujet.
\item À la fin de cette leçon, je sais construire un plan dialectique, analytique ou thématique cohérent et progressif.
\item À la fin de cette leçon, je sais rédiger une introduction en entonnoir (accroche, définition, problématique, annonce du plan) et une conclusion synthétique et ouverte.
\item À la fin de cette leçon, je sais argumenter chaque sous-partie selon le schéma : thèse -> argument -> exemple concret (Cameroun/Monde) -> lien avec la problématique.
\item À la fin de cette leçon, je sais mobiliser des références philosophiques (auteurs, concepts) et des faits de société camerounais pour étayer ma démonstration.
\item À la fin de cette leçon, je sais m'auto-évaluer à l'aide de la grille officielle de notation du Baccalauréat technique.
\item À la fin de cette leçon, je sais corriger les erreurs fréquentes : hors-sujet, paraphrase, catalogue d'auteurs, absence de transition.
\end{itemize}
\end{objectifs}

\begin{prerequis}
\begin{itemize}
\item Définition de la philosophie et distinction opinion/savoir (programme 4ème / début 1ère).
\item Structure de base d'un raisonnement : thèse, argument, exemple, conclusion.
\item Maîtrise des connecteurs logiques (donc, cependant, par ailleurs, en outre, en définitive).
\item Connaissance de base de quelques grands auteurs au programme (Socrate, Descartes, Kant, Marx, Eboussi Boulaga, etc.).
\item Technique de la prise de notes et de la lecture de textes philosophiques courts.
\end{itemize}
\end{prerequis}

\newpage

\coursec{L'analyse rigoureuse du sujet : la clé de voûte}

\courssub{Situation d'accroche : du vacarme des réseaux au silence de la pensée}

Ouvrez Facebook ou WhatsApp un soir à Yaoundé ou à Douala : un débat virulent oppose les partisans de la \og modernité \fg{} à ceux de la \og tradition \fg{} pour expliquer le sous-développement du Cameroun. Les uns accusent les coutumes d'être des freins archaïques ; les autres dénoncent une modernité importée qui détruirait nos valeurs. Mais si l'on regarde de près, que trouve-t-on ? Des slogans répétés, des attaques personnelles (\emph{ad hominem}), des généralités sans preuve : \og les Blancs ont tout détruit \fg{}, \og nos ancêtres étaient des sorciers \fg{}. Personne ne définit ses mots, personne ne justifie ses affirmations, personne n'écoute l'adversaire.

Face à ce vacarme, l'élève de Première technique doit-il se taire ? Non. Mais il ne peut intervenir utilement qu'à une condition : transformer son opinion brute en une pensée structurée. C'est exactement le travail de la dissertation philosophique, et ce travail commence toujours au même endroit : \textbf{l'analyse du sujet}. Avant d'écrire une seule ligne, avant même de penser un plan, il faut comprendre \emph{exactement} ce que le sujet demande. C'est la \cle{clé de voûte} de toute la copie : si elle est mal posée, tout l'édifice s'effondre, aussi belle que soit l'écriture.

\textbf{Problématique de la section :} comment passer d'un énoncé de quelques mots, souvent banal en apparence, à une question philosophique précise, tensionnelle et féconde, capable de guider un développement en trois parties ?

\courssub{Identifier la discipline et le type de sujet}

\begin{definition}[Analyse du sujet]
L'\cle{analyse du sujet} est la première étape de la dissertation philosophique. Elle consiste à décomposer l'énoncé pour en identifier la discipline, le type, les termes clés, les présupposés et l'enjeu, afin d'en dégager une problématique. Elle se fait au brouillon et ne doit jamais être négligée : elle conditionne la pertinence de toute la copie.
\end{definition}

\textbf{Premier réflexe : confirmer la discipline.} Au Probatoire comme au Baccalauréat technique, le sujet relève de la \emph{philosophie}, non de l'histoire, de la sociologie ou de la morale religieuse. Cela signifie que la réponse attendue n'est ni un récit de faits, ni une enquête statistique, ni un sermon : c'est un \emph{raisonnement argumenté} qui interroge des concepts (le travail, l'État, la liberté, la technique, la coutume...).

\textbf{Deuxième réflexe : identifier le type de sujet.} On distingue classiquement quatre formes :

\begin{itemize}
\item \textbf{La question directe en \og Est-ce que... ? \fg{}} : \og La corruption est-elle un mal nécessaire ? \fg{} Elle appelle une réponse nuancée (oui / non / sous quelles conditions ?) et invite à un plan dialectique.
\item \textbf{La question normative en \og Faut-il... ? \fg{}} : \og Faut-il obéir à l'État ? \fg{} Elle porte sur un devoir, une valeur, une légitimité. Attention : \og faut-il \fg{} ne demande pas \og est-ce possible \fg{} mais \og est-ce juste, légitime, souhaitable \fg{}.
\item \textbf{Le sujet nominal (un concept seul)} : \og Le travail \fg{}, \og La technique \fg{}. Le plus difficile, car rien n'est posé : c'est au candidat de construire entièrement la question.
\item \textbf{La citation d'auteur} : par exemple une phrase de Kant ou de Nietzsche à discuter. Il faut d'abord \emph{expliquer} la thèse de l'auteur, puis la \emph{discuter} (l'approuver, la critiquer, la dépasser).
\end{itemize}

\begin{attention}[Ne pas confondre les types de sujets]
Une erreur fréquente consiste à traiter un \og Faut-il... ? \fg{} comme un \og Est-ce que... ? \fg{}. \og Faut-il respecter la coutume ? \fg{} n'interroge pas la réalité du respect (les gens respectent-ils ?) mais sa \emph{légitimité} (a-t-on de bonnes raisons de respecter ?). Identifier le type de sujet, c'est déjà identifier la nature de la réponse attendue.
\end{attention}

\courssub{Définir les termes clés : le dictionnaire philosophique, pas le sens courant}

Une fois le type identifié, on souligne les \cle{termes clés} du sujet et on les définit. Mais attention : il ne s'agit pas du sens courant, celui du marché ou de la rue. Il s'agit du sens \emph{philosophique}, précis et réfléchi.

\begin{exemplebox}[Du sens courant au sens philosophique]
\begin{itemize}
\item \textbf{Travail} — Sens courant : \og avoir un job \fg{}. Sens philosophique : activité par laquelle l'homme transforme la nature et se transforme lui-même (Hegel, Marx), source à la fois d'aliénation et d'humanisation.
\item \textbf{Technique} — Sens courant : \og les machines, les téléphones \fg{}. Sens philosophique : ensemble des moyens artificiels par lesquels l'homme étend ses pouvoirs sur la nature ; elle pose la question de savoir si l'outil sert l'homme ou l'asservit.
\item \textbf{État} — Sens courant : \og le gouvernement, les ministres \fg{}. Sens philosophique : institution politique détenant le monopole de la violence légitime (Max Weber) et chargée d'assurer l'ordre et la justice.
\item \textbf{Coutume} — Sens courant : \og les traditions du village \fg{}. Sens philosophique : ensemble de règles non écrites, héritées des ancêtres, qui régissent la vie d'une communauté ; elle pose le problème du rapport entre tradition et raison critique.
\end{itemize}
\end{exemplebox}

Définir, c'est déjà philosopher : en précisant le mot \og nécessaire \fg{} dans \og mal nécessaire \fg{}, on découvre que le sujet suppose que la corruption pourrait être \emph{indispensable} au fonctionnement social — ce qui est loin d'être évident.

\courssub{Dégager les présupposés et l'alternative fondamentale}

\begin{definition}[Présupposé]
Un \cle{présupposé} est ce qu'un énoncé admet implicitement, sans le dire ni le démontrer. Le repérer, c'est découvrir ce que le sujet \emph{prend pour acquis} et qu'il faut interroger.
\end{definition}

Prenons le sujet : \og L'école technique forme-t-elle pour le chômage ? \fg{} Quels présupposés cache-t-il ?

\begin{enumerate}
\item Que l'école a pour mission essentielle de \emph{former pour l'emploi} (présupposé utilitariste).
\item Qu'il existe un lien direct entre formation et insertion professionnelle.
\item Que le chômage des techniciens serait imputable à l'école, et non au marché du travail ou à l'économie nationale.
\end{enumerate}

Chaque présupposé peut être contesté : l'école ne forme-t-elle pas aussi des \emph{citoyens} et des \emph{esprits critiques}, au-delà de l'employabilité ? C'est en secouant les présupposés que naît le débat.

\textbf{L'alternative fondamentale.} Tout bon sujet de dissertation cache une tension entre deux réponses opposées : une \cle{thèse} et une \cle{antithèse}. Pour \og La corruption est-elle un mal nécessaire ? \fg{} :

\begin{itemize}
\item \textbf{Thèse possible :} la corruption serait un mal nécessaire — elle \og lubrifie \fg{} l'administration, accélère les dossiers, contourne des règles trop rigides (certains parlent de \og graissage de rouages \fg{}).
\item \textbf{Antithèse :} la corruption est un mal tout court — elle détruit la confiance, détourne les ressources publiques (hôpitaux sans médicaments, routes inachevées), et rien de ce qui est injuste ne peut être \og nécessaire \fg{}.
\end{itemize}

Le dépassement (synthèse) consistera alors à distinguer : ce qui est parfois \emph{inévitable} n'est jamais pour autant \emph{légitime} ; la nécessité factuelle ne vaut pas justification morale.

\courssub{Formuler la problématique : la question tensionnelle}

\begin{definition}[Problématique]
La \cle{problématique} est la question philosophique dérivée du sujet, qui met en tension au moins deux thèses opposées et guide tout le développement. Elle se formule en une ou deux phrases interrogatives, placées à la fin de l'introduction.
\end{definition}

\begin{methode}[La méthode des 3 questions pour problématiser]
\begin{enumerate}
\item \textbf{Que dit le sujet ?} Reformulez l'énoncé avec vos propres mots, après avoir défini les termes. (Ex. : le sujet affirme que la technique pourrait être à la fois source de liberté et source d'asservissement.)
\item \textbf{Quel est le paradoxe, la tension ?} Identifiez les deux réponses contraires que le sujet autorise. (Ex. : d'un côté la technique soulage l'homme ; de l'autre elle le rend dépendant et peut l'aliéner.)
\item \textbf{Comment formuler la question qui relie les deux ?} Construisez une phrase interrogative qui tient les deux pôles ensemble. (Ex. : \og La technique libère-t-elle l'homme ou l'aliène-t-elle dans le contexte industriel camerounais ? \fg{})
\end{enumerate}
\end{methode}

\begin{attention}[Définition n'est pas problématique]
Erreur fréquente : confondre la définition des termes et la problématique. La \textbf{définition} explique les mots (\og la coutume est un ensemble de règles héritées... \fg{}) ; la \textbf{problématique} lance le débat (\og la coutume est-elle un guide ou un frein pour la raison ? \fg{}). Une copie qui définit bien mais ne problématise pas reste descriptive : elle ne dépasse guère la moyenne.
\end{attention}

\courssub{Vérifier la pertinence de la problématique}

Avant de passer au plan, soumettez votre problématique à deux tests :

\begin{enumerate}
\item \textbf{Test de fidélité :} ma problématique répond-elle \emph{exactement} au sujet posé, sans le déplacer ? Si le sujet porte sur la \emph{légitimité politique} du chef traditionnel, une problématique sur son \emph{rôle économique} est un hors-sujet partiel.
\item \textbf{Test de fécondité :} ma problématique permet-elle un débat en trois parties (thèse / antithèse / dépassement) ? Si la question n'appelle qu'un \og oui \fg{} évident ou un simple récit, elle est mauvaise.
\end{enumerate}

\begin{aretenir}[L'essentiel de la méthode d'analyse]
\begin{enumerate}
\item Identifier la discipline (philosophie) et le type de sujet.
\item Définir les termes clés au sens philosophique.
\item Dégager les présupposés et l'alternative thèse/antithèse.
\item Formuler une problématique tensionnelle.
\item Vérifier : fidélité au sujet + possibilité d'un plan en trois parties.
\end{enumerate}
\end{aretenir}

\begin{popfigure}
\begin{center}
\begin{tikzpicture}[
  grow=down,
  level distance=1.15cm,
  level 1/.style={sibling distance=4.6cm},
  level 2/.style={sibling distance=2.6cm},
  every node/.style={draw=popblue, fill=popblueL, rounded corners=2pt,
    text width=3.1cm, align=center, font=\small, inner sep=3pt},
  edge from parent/.style={draw=popdark, -{Stealth[length=2mm]}, thick}
]
\node[fill=poporangeL, draw=poporange, text width=4.2cm, font=\small\bfseries] {SUJET\\ \og L'école technique forme-t-elle pour le chômage ? \fg{}}
  child { node[align=center]{Termes à définir\\ \emph{école technique, former, chômage}}
    child { node[align=center]{Type de sujet\\ question \og Est-ce que... ? \fg{}} }
    child { node[align=center]{Présupposés\\ école = emploi ; chômage = faute de l'école} }
  }
  child { node[align=center]{Alternative\\ Thèse : oui (formation inadaptée) / Antithèse : non (école libératrice)}
    child { node[fill=popgreenL, draw=popgreen, align=center]{Problématique\\ L'école technique prépare-t-elle à l'emploi ou à l'échec ?} }
    child { node[fill=popgreenL, draw=popgreen, align=center]{Plan possible\\ I. Oui \quad II. Non \quad III. Former l'homme et le technicien} }
  };
\end{tikzpicture}
\end{center}
\legende{Arbre méthodologique de l'analyse d'un sujet de dissertation : du sujet brut (orange) vers la problématique et le plan (vert), en passant par la définition des termes, le repérage des présupposés et l'alternative fondamentale.}
\end{popfigure}

\begin{popfigure}
\begin{center}
\begin{tikzpicture}
\node[draw=popink, fill=poppinkL, rounded corners=3pt, text width=6.6cm, align=left, font=\small, inner sep=6pt] (bad) at (0,0)
{\textbf{Sujet mal analysé (hors-sujet)}\\[2pt]
Sujet : \og Le chef traditionnel a-t-il encore une légitimité politique ? \fg{}\\[2pt]
$\bullet$ Définition absente ou folklorique\\
$\bullet$ Récit : \og mon chef de village est... \fg{}\\
$\bullet$ Problématique déplacée : \og le chef est-il riche ? \fg{}\\
$\bullet$ Plan : liste d'anecdotes, aucun débat\\
$\bullet$ Note constatée : souvent $< 8/20$};
\node[draw=popgreen, fill=popgreenL, rounded corners=3pt, text width=6.6cm, align=left, font=\small, inner sep=6pt, right=0.5cm of bad] (good)
{\textbf{Sujet bien analysé (problématique tensionnelle)}\\[2pt]
$\bullet$ Termes définis : \emph{légitimité} = droit reconnu de commander ; \emph{politique} = relatif au pouvoir sur la cité\\
$\bullet$ Présupposé : la légitimité moderne (État, élections, lois) concurrencerait la tradition\\
$\bullet$ Problématique : \og l'autorité coutumière peut-elle coexister légitimement avec l'État républicain ? \fg{}\\
$\bullet$ Plan dialectique en 3 parties};
\draw[-{Stealth[length=3mm]}, very thick, poporange] (bad.east) -- (good.west)
  node[midway, above, font=\small\bfseries, text=poporange]{méthode};
\end{tikzpicture}
\end{center}
\legende{Tableau comparatif : à gauche, une copie qui fonce sans analyser (hors-sujet) ; à droite, la même copie reconstruite par la méthode d'analyse (définitions, présupposés, problématique tensionnelle).}
\end{popfigure}

\begin{exemplebox}[Un chiffre qui doit faire réfléchir]
D'après les rapports de jury récents (MINESEC), environ \textbf{68\,\%} des copies notées en dessous de $8/20$ au Probatoire de Philosophie technique devaient leur échec à une mauvaise analyse du sujet (hors-sujet, problématique absente ou déplacée), et non à une mauvaise expression écrite. Autrement dit : ce n'est pas le style qui sauve une copie, c'est la justesse de l'analyse initiale.
\end{exemplebox}

\begin{savaistu}[D'où vient le mot \og dissertation \fg{} ?]
Le mot \og dissertation \fg{} vient du latin \emph{dissertare}, qui signifie \og discuter, débattre, examiner en détail \fg{}. Une dissertation n'est donc pas une récitation de cours apprise par cœur : c'est un véritable \emph{procès du sujet}, où l'on entend l'accusation (thèse), la défense (antithèse), avant de rendre un verdict motivé (synthèse). Celui qui récite plaide sans écouter ; celui qui analyse juge en connaissance de cause.
\end{savaistu}

\courssub{Atelier guidé : trois sujets types Probatoire / Baccalauréat technique}

\begin{exoresolu}[Sujet 1 : \og La corruption est-elle un mal nécessaire ? \fg{}]
Appliquez la méthode complète d'analyse à ce sujet : type, définitions, présupposés, alternative, problématique.
\tcblower
\textbf{Solution détaillée.}
\begin{enumerate}
\item \textbf{Type de sujet :} question directe (\og Est-ce que... ? \fg{}), de nature morale et politique.
\item \textbf{Définitions :} \emph{corruption} : usage illégitime d'une fonction publique ou privée pour un avantage personnel (pots-de-vin, détournements, favoritisme). \emph{Mal nécessaire} : ce qui est nuisible en soi mais serait indispensable pour obtenir un bien ou éviter un pire.
\item \textbf{Présupposés :} (a) la corruption serait parfois utile, voire inévitable dans certaines administrations ; (b) on pourrait hiérarchiser les maux et tolérer le moindre ; (c) la fin pourrait justifier le moyen.
\item \textbf{Alternative :} \textbf{Thèse} — la corruption facilite les démarches, \og débloque \fg{} une administration lente : certains la présentent comme un lubrifiant social. \textbf{Antithèse} — la corruption est un mal absolu : elle vole les ressources destinées aux écoles et aux hôpitaux, détruit la confiance et la justice ; un mal toléré devient une habitude qui corrompt toute la société.
\item \textbf{Problématique :} \og Ce qui est socialement inévitable est-il pour autant moralement acceptable ? La corruption peut-elle être justifiée comme un moyen, ou sa tolérance détruit-elle le bien commun qu'elle prétend servir ? \fg{}
\item \textbf{Vérification :} fidélité (on parle bien de corruption et de nécessité) ; fécondité (plan possible : I. Les apparences de nécessité — II. La réfutation : un mal reste un mal — III. Distinguer l'inévitable du légitime et exiger l'intégrité).
\end{enumerate}
\end{exoresolu}

\begin{exoresolu}[Sujets 2 et 3 : analyse comparée]
Analysez brièvement : (a) \og Le chef traditionnel a-t-il encore une légitimité politique ? \fg{} ; (b) \og L'école technique forme-t-elle pour le chômage ? \fg{}
\tcblower
\textbf{(a) Le chef traditionnel.} Type : question directe sur la légitimité. Définitions : \emph{chef traditionnel} : autorité coutumière héréditaire ou élue selon les rites, garante des traditions ; \emph{légitimité politique} : droit reconnu d'exercer le pouvoir. Présupposé : la modernité républicaine (État, élections, lois) aurait rendu la chefferie obsolète. Alternative : la chefferie conserve une légitimité (proximité, sagesse, rôle d'arbitrage et de cohésion dans les villages) / elle est dépassée (pouvoir non démocratique, parfois source de conflits de succession). Problématique : \og l'autorité héritée de la tradition peut-elle coexister légitimement avec la légalité républicaine ? \fg{}

\textbf{(b) L'école technique.} Type : question directe, à présupposé polémique. Définitions : \emph{former} : transmettre des compétences et développer l'esprit ; \emph{chômage} : situation de celui qui cherche un emploi sans en trouver. Présupposés : l'école vise uniquement l'emploi ; le chômage serait causé par la formation. Alternative : la formation technique serait inadaptée au marché (thèse) / elle donne des compétences concrètes et l'esprit d'initiative, et le chômage relève de la conjoncture économique (antithèse). Problématique : \og la mission de l'école technique est-elle de produire des employés ou de former des hommes capables de créer leur propre emploi ? \fg{} — ce qui ouvre un dépassement fécond : former à la fois le technicien compétent et l'entrepreneur autonome.
\end{exoresolu}

\begin{exercice}[À vous de jouer]
Sujet : \og Faut-il toujours dire la vérité ? \fg{} En vous aidant de l'arbre méthodologique de la figure 1 : 1) identifiez le type de sujet ; 2) définissez \emph{vérité} et le sens du \og faut-il \fg{} ; 3) dégagez deux présupposés ; 4) formulez l'alternative thèse/antithèse ; 5) rédigez une problématique tensionnelle et vérifiez sa pertinence.
\end{exercice}

\begin{corrige}[Exercice : \og Faut-il toujours dire la vérité ? \fg{}]
1) \textbf{Type :} question normative (\og faut-il \fg{} = devoir moral). 2) \textbf{Définitions :} \emph{vérité} : accord de la pensée avec la réalité, ou parole conforme à ce qu'on pense ; \og faut-il toujours \fg{} interroge l'existence d'un devoir \emph{sans exception}. 3) \textbf{Présupposés :} (a) dire la vérité serait un devoir ; (b) il pourrait exister des cas où ce devoir entre en conflit avec d'autres devoirs (ne pas nuire, protéger). 4) \textbf{Alternative :} Kant soutient que le mensonge est toujours condamnable, car il détruit la confiance qui fonde toute société (thèse) ; mais dire la vérité à un agresseur sur la cachette d'un innocent serait criminel : le devoir de vérité semble parfois céder devant le devoir d'humanité (antithèse). 5) \textbf{Problématique :} \og le devoir de vérité est-il absolu, ou la sincérité doit-elle être ordonnée au respect de la personne humaine ? \fg{} Vérification : la question répond exactement au \og faut-il toujours \fg{} et permet un plan en trois parties (I. Oui, la vérité fonde la société — II. Non, il existe des limites — III. La vérité comme exigence ordonnée à la justice et à la bienveillance).
\end{corrige}

\begin{aretenir}[Conclusion de la section]
Analyser un sujet, c'est refuser de répondre trop vite. Le philosophe, comme le bon technicien devant une machine en panne, commence par \emph{diagnostiquer} avant de réparer : il identifie le type de sujet, définit les termes au sens philosophique, débusque les présupposés, oppose thèse et antithèse, puis formule une problématique tensionnelle dont il vérifie la fidélité et la fécondité. Cette lenteur initiale est la condition de la rapidité et de la justesse de tout le reste de la copie. La section suivante montrera comment, à partir de cette problématique, construire l'architecture du développement.
\end{aretenir}

\coursec{L'architecture du développement : Plans et articulation}

Dans la section précédente, nous avons appris à \emph{décortiquer} un sujet : définir les termes, dégager la problématique, formuler la question centrale. Mais une problématique, aussi fine soit-elle, reste une intention. Il faut maintenant lui donner un \cle{corps}, une \cle{ossature}. C'est toute la fonction du \cle{plan} : transformer une question en un chemin de pensée. Un élève qui a bien analysé son sujet mais qui rédige sans plan ressemble à un maçon de Yaoundé qui posséderait d'excellents matériaux — ciment, fer, sable — mais aucun plan de construction : il empile, et l'édifice s'écroule. Le plan est l'architecture invisible qui porte toute la copie.

\courssub{Le choix du plan selon le type de sujet}

\textbf{L'intuition d'abord.} On ne construit pas la même maison pour une famille nombreuse de Douala et pour un commerçant de Bafoussam. De même, on ne construit pas le même plan pour toutes les formes de sujets. La \emph{forme grammaticale} du sujet est un signal : elle indique le type de réponse attendu.

\begin{definition}[Les trois grands types de plans]
On distingue classiquement trois architectures de développement :
\begin{itemize}
\item le \cle{plan dialectique} (Thèse -- Antithèse -- Synthèse), adapté aux sujets interrogatifs du type \og Est-ce que... ? \fg{} ou \og Faut-il... ? \fg{} ;
\item le \cle{plan analytique} (Concepts -- Causes -- Conséquences, ou Constat -- Causes -- Solutions), adapté aux sujets nominaux du type \og Le travail \fg{}, \og La liberté \fg{} ;
\item le \cle{plan thématique} (Auteurs, Époques ou Aspects successifs), adapté au commentaire de citation.
\end{itemize}
\end{definition}

\textbf{Le plan dialectique.} Face à une question fermée (\og La technique est-elle libératrice ? \fg{}), l'esprit hésite naturellement entre le oui et le non. Le plan dialectique épouse ce mouvement : la partie I expose la thèse spontanée (le oui), la partie II la conteste (le non), la partie III dépasse l'opposition en montrant que la question était mal posée ou en découvrant un point de vue supérieur.

\textbf{Le plan analytique.} Face à un concept nu (\og Le devoir \fg{}), il n'y a rien à réfuter : il faut \emph{analyser}, c'est-à-dire décomposer. On examine d'abord ce qu'est la chose (définition, concepts), puis d'où elle vient (causes, fondements), enfin ce qu'elle produit (conséquences, enjeux).

\textbf{Le plan thématique.} Face à une citation d'auteur, on suit le fil du texte : on explicite la thèse de l'auteur (I), on la confronte à d'autres doctrines ou à l'expérience (II), on en discute la portée et les limites (III).

\begin{popfigure}
\begin{center}
\begin{tikzpicture}[font=\small, node distance=0.35cm,
plan/.style={draw=popblue, thick, fill=popblueL, rounded corners=2pt, minimum width=3.4cm, minimum height=0.75cm, align=center},
titre/.style={draw=popdark, thick, fill=popdark!10, rounded corners=2pt, minimum width=3.4cm, minimum height=0.9cm, align=center, font=\small\bfseries},
etq/.style={font=\small\bfseries, text=popdark}]
% Dialectique
\node[etq] at (-4.6,3.1) {Plan dialectique};
\node[titre] (d0) at (-4.6,2.4) {\og Faut-il... ? \fg{}};
\node[plan, below=of d0, align=center] (d1){I. Thèse\\ \emph{Oui : la technique libère}};
\node[plan, below=of d1, align=center] (d2){II. Antithèse\\ \emph{Non : elle aliène}};
\node[plan, fill=popgreenL, draw=popgreen, below=of d2, align=center] (d3){III. Synthèse\\ \emph{Tout dépend de l'usage}};
\draw[-{Stealth}, thick, popblue] (d0) -- (d1);
\draw[-{Stealth}, thick, popblue] (d1) -- (d2);
\draw[-{Stealth}, thick, popblue] (d2) -- (d3);
% Analytique
\node[etq] at (0,3.1) {Plan analytique};
\node[titre] (a0) at (0,2.4) {\og Le travail \fg{}};
\node[plan, below=of a0, align=center] (a1){I. Concepts\\ \emph{Qu'est-ce que travailler ?}};
\node[plan, below=of a1, align=center] (a2){II. Causes\\ \emph{Pourquoi travaille-t-on ?}};
\node[plan, fill=popgreenL, draw=popgreen, below=of a2, align=center] (a3){III. Conséquences\\ \emph{Que produit le travail ?}};
\draw[-{Stealth}, thick, popblue] (a0) -- (a1);
\draw[-{Stealth}, thick, popblue] (a1) -- (a2);
\draw[-{Stealth}, thick, popblue] (a2) -- (a3);
% Thématique
\node[etq] at (4.6,3.1) {Plan thématique};
\node[titre] (t0) at (4.6,2.4) {Citation d'auteur};
\node[plan, below=of t0, align=center] (t1){I. Thèse de l'auteur\\ \emph{Explication du texte}};
\node[plan, below=of t1, align=center] (t2){II. Confrontation\\ \emph{Autres doctrines, expérience}};
\node[plan, fill=popgreenL, draw=popgreen, below=of t2, align=center] (t3){III. Discussion\\ \emph{Portée et limites}};
\draw[-{Stealth}, thick, popblue] (t0) -- (t1);
\draw[-{Stealth}, thick, popblue] (t1) -- (t2);
\draw[-{Stealth}, thick, popblue] (t2) -- (t3);
\end{tikzpicture}
\end{center}
\legende{Les trois types de plans au programme de Première technique. La forme du sujet (interrogative, nominale ou citation) commande le choix de l'architecture. La partie III (en vert) est toujours le moment du dépassement.}
\end{popfigure}

\begin{propriete}[Les trois critères d'un bon plan]
Un plan de dissertation philosophique est recevable s'il respecte trois exigences :
\begin{enumerate}
\item la \cle{cohérence} : unité de ton et de point de vue, chaque partie répond à la même problématique ;
\item la \cle{progressivité} : montée en puissance, chaque partie va plus loin que la précédente ;
\item l'\cle{exhaustivité} : tous les aspects de la problématique dégagée à l'analyse du sujet sont traités.
\end{enumerate}
\end{propriete}

\courssub{La construction du plan détaillé}

\textbf{L'intuition.} Un plan annoncé en introduction (\og I..., II..., III... \fg{}) ne suffit pas : il faut un \cle{plan détaillé}, véritable feuille de route de la rédaction. La règle d'or au Probatoire et au Baccalauréat technique : \textbf{trois grandes parties} (I, II, III), chacune divisée en \textbf{deux ou trois sous-parties} (A, B, C), chaque sous-partie portant \emph{une seule idée force}.

\begin{methode}[La « fiche de plan » au brouillon]
Avant de rédiger, construis au brouillon un tableau à cinq colonnes :
\begin{enumerate}
\item \textbf{Colonne 1 — Partie} : I, II, III puis A, B, C ;
\item \textbf{Colonne 2 — Idée force} : la thèse de la sous-partie en une phrase ;
\item \textbf{Colonne 3 — Arguments et exemples} : au moins un exemple concret, de préférence camerounais ;
\item \textbf{Colonne 4 — Auteurs et concepts} : la référence philosophique mobilisée ;
\item \textbf{Colonne 5 — Transition} : le lien logique vers la sous-partie suivante.
\end{enumerate}
Si une case reste vide (pas d'exemple, pas de transition), la sous-partie est faible : il faut la repenser ou la fusionner avec une autre.
\end{methode}

\textbf{La rédaction des titres.} Les titres de parties et de sous-parties ne s'écrivent pas sur la copie finale, mais ils doivent exister dans ta tête et au brouillon. Ils obéissent à trois règles : ce sont des \emph{phrases nominales affirmatives} (pas de questions), ils sont \emph{parallèles} (même construction grammaticale d'une partie à l'autre), et surtout ils \emph{racontent l'histoire de la pensée} : lus à la suite, ils forment un résumé de ta démonstration.

\begin{exemplebox}[Un plan détaillé sur \og La technique est-elle libératrice ? \fg{}]
\begin{itemize}
\item \textbf{I. La technique comme puissance de libération de l'homme}
\begin{itemize}
\item A. La maîtrise des forces naturelles (ex. : les barrages hydroélectriques d'Édéa et de Song Loulou qui éclairent le Cameroun ; Prométhée, le mythe du feu volé aux dieux).
\item B. L'allègement du travail pénible (ex. : la motoculteur remplaçant la houe dans les champs de maïs de l'Ouest ; Marx : la machine peut réduire le temps de travail nécessaire).
\end{itemize}
\item \textbf{II. La technique comme nouvelle servitude}
\begin{itemize}
\item A. L'aliénation de l'ouvrier par la machine (Marx : le travail aliéné, l'ouvrier devient rouage).
\item B. La dépendance technique des sociétés modernes (ex. : la panne de réseau mobile qui paralyse les transactions Mobile Money à Douala ; Heidegger : l'essence de la technique nous enferme dans le \emph{Gestell}, l'arraisonnement).
\end{itemize}
\item \textbf{III. La technique comme responsabilité : l'usage décide de la valeur}
\begin{itemize}
\item A. La neutralité de l'outil et la responsabilité de l'usager (le même téléphone informe l'élève ou le distrait).
\item B. Vers une technique maîtrisée au service du développement endogène (Eboussi Boulaga : penser le développement à partir des besoins réels des communautés camerounaises).
\end{itemize}
\end{itemize}
Lus à la suite, ces titres racontent déjà toute la démonstration.
\end{exemplebox}

\begin{attention}[Le plan « catalogue »]
L'erreur la plus fréquente et la plus sanctionnée : le plan \emph{catalogue} — A. Platon, B. Descartes, C. Kant — qui aligne des auteurs sans lien logique avec la problématique. \textbf{On planifie des idées, pas des auteurs.} L'auteur est un \emph{témoin} convoqué pour étayer une idée ; il n'est pas lui-même l'idée. Un correcteur qui voit un catalogue de noms conclut immédiatement que l'élève récite sans penser.
\end{attention}

\courssub{L'articulation : transitions et références}

\textbf{L'intuition.} Une dissertation n'est pas trois petits textes juxtaposés : c'est \emph{un} texte continu. Les \cle{transitions} sont les soudures entre les parties. Sans elles, le correcteur a l'impression de sauts brutaux ; avec elles, il suit un raisonnement qui avance.

\begin{definition}[La transition]
Une transition est un court passage (2 à 4 lignes) placé entre deux sous-parties ou deux parties, qui comporte trois temps :
\begin{enumerate}
\item un \cle{bilan partiel} : résumé de ce qui vient d'être établi ;
\item un \cle{lien logique} : un connecteur (\emph{mais}, \emph{cependant}, \emph{par ailleurs}, \emph{si... alors}, \emph{en vérité}) qui annonce le rapport entre les deux moments ;
\item une \cle{annonce} : la présentation de l'idée force suivante.
\end{enumerate}
\end{definition}

\begin{exemplebox}[Une transition réussie]
\og Ainsi, la technique apparaît d'abord comme une puissance d'affranchissement qui libère le Camerounais des fatigues ancestrales. \emph{Cependant}, cette libération n'a-t-elle pas un prix ? L'outil qui décharge l'homme de l'effort ne risque-t-il pas de le déposséder de son propre geste ? C'est ce que nous devons maintenant examiner : la technique comme nouvelle servitude. \fg{}
\end{exemplebox}

\textbf{L'intégration des références.} Règle minimale : \textbf{un auteur ou un concept par sous-partie}. La référence doit être \emph{exploitée}, pas simplement citée : on donne la thèse de l'auteur en une phrase, on la relie à l'argument, on l'illustre. Ainsi Marx pour le travail aliéné, Heidegger pour l'essence de la technique, Eboussi Boulaga pour la démocratie et le développement au Cameroun.

\begin{popfigure}
\begin{center}
\begin{tikzpicture}[font=\small, node distance=0.5cm and 0.55cm,
boite/.style={draw=popblue, thick, fill=popblueL, rounded corners=2pt, minimum width=2.5cm, minimum height=0.9cm, align=center},
boiteT/.style={draw=poporange, thick, fill=poporangeL, rounded corners=2pt, minimum width=2.5cm, minimum height=0.9cm, align=center}]
\node[boite] (i1) {Idée force 1};
\node[boite, right=of i1] (a1) {Argument};
\node[boite, right=of a1, align=center] (e1){Exemple\\ Cameroun};
\node[boite, right=of e1, align=center] (au1){Auteur /\\ Concept};
\node[boiteT, below=0.6cm of au1, align=center] (tr){Transition\\ (bilan + lien + annonce)};
\node[boite, below=0.6cm of i1, fill=popgreenL, draw=popgreen] (i2) {Idée force 2};
\draw[-{Stealth}, thick, popblue] (i1) -- (a1);
\draw[-{Stealth}, thick, popblue] (a1) -- (e1);
\draw[-{Stealth}, thick, popblue] (e1) -- (au1);
\draw[-{Stealth}, thick, poporange] (au1) -- (tr);
\draw[-{Stealth}, thick, poporange] (tr.west) -| (i2.north);
\draw[-{Stealth}, thick, popgreen] (i2.east) -- ++(0.7,0) node[right, text=popdark]{\emph{et le cycle recommence...}};
\end{tikzpicture}
\end{center}
\legende{Le cycle d'une sous-partie : de l'idée force à la transition. Chaque sous-partie enchaîne idée, argument, exemple camerounais et référence, puis se soude à la suivante par une transition.}
\end{popfigure}

\textbf{La progressivité de la partie III.} La synthèse n'est pas un compromis mou du type \og la technique est à la fois bonne et mauvaise \fg{}. Elle doit apporter une \emph{lumière nouvelle} : montrer que l'opposition initiale reposait sur une confusion, déplacer la question (de \og la technique est-elle bonne ? \fg{} à \og quel usage en faisons-nous ? \fg{}), ou ouvrir sur une exigence (la responsabilité, l'éducation, la politique).

\begin{exoresolu}[Construire le plan de \og Faut-il obéir à l'État ? \fg{}]
Énoncé : propose un plan dialectique détaillé (titres de parties et sous-parties, avec une référence par sous-partie) pour le sujet : \emph{Faut-il obéir à l'État ?}
\tcblower
\textbf{Solution détaillée.}
\begin{itemize}
\item \textbf{I. L'obéissance à l'État comme condition de la vie civile}
\begin{itemize}
\item A. La sortie de l'état de nature par le contrat (Hobbes : sans pouvoir commun, la vie est \og solitaire, misérable, brutale et brève \fg{}).
\item B. La loi comme protection de tous (ex. : le code de la route qui sécurise l'axe Yaoundé--Douala ; Kant : l'obéissance à la loi comme condition de la liberté).
\end{itemize}
\item \textbf{II. Les limites et les dangers de l'obéissance}
\begin{itemize}
\item A. L'obéissance peut servir l'injustice (Sophocle, \emph{Antigone} : les lois non écrites contre l'édit de Créon).
\item B. Le devoir de désobéissance civile face à la loi injuste (Thoreau : la désobéissance civile ; ex. : les luttes citoyennes pour les libertés au Cameroun).
\end{itemize}
\item \textbf{III. D'une obéissance aveugle à une obéissance éclairée}
\begin{itemize}
\item A. Obéir par raison et non par peur (Kant : \emph{Sapere aude}, ose penser par toi-même).
\item B. Le citoyen actif, fondement de la démocratie (Eboussi Boulaga : la démocratie comme participation effective des communautés à leur destin).
\end{itemize}
\end{itemize}
La partie III dépasse le compromis : elle déplace la question de l'obéissance vers la \emph{qualité} de l'obéissance, éclairée par la raison et la participation citoyenne.
\end{exoresolu}

\begin{exemplebox}[Volume d'un développement type Bac]
Un développement de qualité représente environ \textbf{1,5 à 2 pages par grande partie} (I, II, III), soit \textbf{5 à 6 pages} de développement. Avec une introduction d'une page et une conclusion d'une demi-page, on obtient une copie de 6 à 7 pages, correspondant typiquement à une note de 12 à 14/20, à condition que le plan soit cohérent, progressif et nourri de références.
\end{exemplebox}

\begin{savaistu}[D'où vient le plan dialectique ?]
La structure Thèse -- Antithèse -- Synthèse a été popularisée par les philosophes allemands \textbf{Fichte} et \textbf{Hegel} au début du XIX\textsuperscript{e} siècle, pour qui la pensée et l'histoire avancent par confrontation des contraires. Mais le mouvement est bien plus ancien : \textbf{Socrate}, dans les rues d'Athènes, pratiquait déjà la \cle{maïeutique}, l'art d'\og accoucher les esprits \fg{} : il faisait d'abord formuler à son interlocuteur sa certitude (thèse), la réfutait par des questions (antithèse), pour conduire l'esprit vers une vérité plus solide (synthèse). Quand tu construis un plan dialectique, tu marches donc dans les pas de Socrate.
\end{savaistu}

\begin{aretenir}[L'essentiel de la section]
\begin{itemize}
\item La forme du sujet commande le plan : dialectique pour \og Est-ce que / Faut-il \fg{}, analytique pour un concept, thématique pour une citation.
\item Plan détaillé : 3 parties, 2 à 3 sous-parties chacune, une idée force par sous-partie, des titres nominaux affirmatifs et parallèles.
\item Chaque sous-partie = idée force + argument + exemple camerounais + auteur, soudée à la suivante par une transition (bilan + lien + annonce).
\item Un bon plan est cohérent, progressif et exhaustif ; la partie III dépasse le compromis.
\item On planifie des \emph{idées}, jamais un catalogue d'auteurs.
\end{itemize}
\end{aretenir}

\begin{exercice}[À toi de bâtir]
Pour chacun des sujets suivants, indique le type de plan adapté et rédige les titres des trois parties : a) \og La liberté \fg{} ; b) \og Peut-on se passer de la religion ? \fg{} ; c) \og \emph{La propriété, c'est le vol} (Proudhon) \fg{}. Puis construis la fiche de plan complète (5 colonnes) du sujet b).
\end{exercice}

\begin{corrige}[Exercice 1]
a) Sujet nominal : plan \textbf{analytique} (I. La liberté comme absence de contrainte ; II. Les conditions réelles de la liberté ; III. La liberté comme autonomie et responsabilité). b) Question fermée : plan \textbf{dialectique} (I. La religion semble indispensable ; II. Ses limites et ses dangers ; III. Le sens religieux dépassé en éthique ou en spiritualité libre). c) Citation : plan \textbf{thématique} (I. La thèse de Proudhon ; II. Confrontation avec la défense de la propriété ; III. Discussion : propriété légitime et propriété abusive). La fiche de plan du sujet b) doit comporter, pour chaque sous-partie, une idée force, un exemple (par ex. les églises et mosquées comme écoles de solidarité au Cameroun), une référence (par ex. Kant, \emph{La religion dans les limites de la simple raison}) et une transition rédigée.
\end{corrige}

\coursec{L'art de l'introduction et de la conclusion : les vitrines de la copie}

Après avoir maîtrisé l'architecture du développement — plans dialectique, analytique ou thématique — et l'articulation logique des arguments (section 2), nous abordons désormais les deux « vitrines » de votre copie : l'introduction et la conclusion. Si le développement est le moteur de votre démonstration, l'introduction en est le volant (elle donne la direction) et la conclusion la signature (elle valide le parcours). Au Probatoire comme au Baccalauréat technique, ces deux moments concentrent à eux seuls 4 à 6 points sur 20 : ils sont le premier et le dernier regard du correcteur. Une introduction floue ou une conclusion bâclée plombent irrémédiablement la note, quand bien même le développement serait solide. Cette section vous donne les codes rigoureux, les ancrages camerounais indispensables et l'entraînement chronométré pour en faire des atouts majeurs.

\courssub{L'introduction : l'entonnoir en quatre temps obligatoires}

L'introduction philosophique n'est pas une « mise en bouche » littéraire. C'est un \cle{dispositif logique} qui conduit le lecteur, par un mouvement d'entonnoir, du contexte général à la question précise que vous allez traiter, pour finir par l'annonce de votre itinéraire. Tout manquement à l'un des quatre temps est sanctionné.

\begin{popfigure}
\begin{tikzpicture}[scale=0.9, every node/.style={font=\small}]
% Couleurs
\colorlet{entcol}{popblue}
\colorlet{defcol}{poppurple}
\colorlet{probcol}{poporange}
\colorlet{plancol}{popgreen}
% Formes de l'entonnoir (polygones superposés)
\draw[entcol, line width=2pt, fill=entcol!10] (-5,3) -- (5,3) -- (3,1.5) -- (-3,1.5) -- cycle;
\draw[defcol, line width=2pt, fill=defcol!10] (-3,1.5) -- (3,1.5) -- (1.5,0) -- (-1.5,0) -- cycle;
\draw[probcol, line width=2pt, fill=probcol!10] (-1.5,0) -- (1.5,0) -- (0,-1) -- (0,-1) -- cycle; % Pointe
\draw[plancol, line width=2pt, fill=plancol!10] (0,-1) -- (-1.5,-2.5) -- (1.5,-2.5) -- cycle; % Évasement inversé (flèche vers le bas)
% Étiquettes temps
\node[entcol, align=center, font=\bfseries] at (0, 2.3) {TEMPS 1 : ACCROCHE\\ \textit{(Large : Contexte, Actualité, Citation)}};
\node[defcol, align=center, font=\bfseries] at (0, 0.75) {TEMPS 2 : DÉFINITIONS \& RAPPEL SUJET\\ \textit{(Rétréci : Clarté des termes, Reformulation)}};
\node[probcol, align=center, font=\bfseries] at (0, -0.5) {TEMPS 3 : PROBLÉMATIQUE\\ \textit{(Pointe : Tension, Question unique)}};
\node[plancol, align=center, font=\bfseries] at (0, -1.8) {TEMPS 4 : ANNONCE DU PLAN\\ \textit{(Évasement : I, II, III structurés)}};
% Flèches de guidage
\draw[->, thick, popmuted] (0, 1.5) -- (0, 1.6);
\draw[->, thick, popmuted] (0, 0) -- (0, -0.4);
\draw[->, thick, popmuted] (0, -1) -- (0, -1.4);
\end{tikzpicture}
\legende{Schéma de l'entonnoir de l'introduction : du général (Accroche) au particulier (Problématique), puis ouverture sur la structure (Plan).}
\end{popfigure}

\paragraph{Temps 1 — L'Accroche ancrée (Le « Cameroun » comme preuve de réel).}
Oubliez les « Depuis la nuit des temps... » ou « Depuis Socrate... ». Au Bac technique camerounais, une accroche réussie est \textbf{datée, sourcée, localisée} et \cle{liée explicitement au sujet}. Elle prouve que le sujet n'est pas une abstraction scolaire mais une question vive pour la société camerounaise d'aujourd'hui.

\begin{methode}[Accroche « Cameroun » : 3 étapes obligatoires]
\begin{enumerate}
    \item \textbf{Un fait précis, chiffré ou daté :} Taux de chômage jeunes (INS 2023 : \num{12,8}\% formel, \num{70}\% informel), Promulgation de la loi n°2024/002 sur la décentralisation, Conflit anglophone (crise depuis 2016), Exploitation minière (cobalt, or à Bétaré-Oya), Numérisation de l'administration (e-governance), Crise alimentaire (importation riz/blé).
    \item \textbf{Une source identifiable :} INS (Institut National de la Statistique), Journal Officiel (JO), Rapport MINJEC/MINEFOP, Auteur local (Mongo Beti, Werewere Liking, Fabien Eboussi Boulaga), Article de presse daté (Cameroon Tribune, Le Jour, Mutations).
    \item \textbf{Un lien explicite avec le sujet :} Phrase de bascule : « Ce constat illustre la tension entre... », « Cette réalité pose avec acuité la question de... », « C'est à la lumière de cette situation que l'on comprend l'enjeu du sujet : ... »
\end{enumerate}
\end{methode}

\begin{exemplebox}[Accroches types selon les thèmes du programme]
\textbf{Sujet : « Le travail est-il une aliénation ? »} \\
« Selon l'Enquête Emploi et Secteur Informel (EESI 2023) de l'INS, plus de \num{70}\% de la population active camerounaise évolue dans le secteur informel, souvent sans contrat, ni protection sociale, ni revenu décent. Cette précarité massive, loin de réaliser l'homme par le travail, semble le réduire à une survie quotidienne. Dès lors, le travail, au Cameroun comme ailleurs, est-il nécessairement aliénation ou peut-il demeurer voie d'émancipation ? »

\textbf{Sujet : « La technique libère-t-elle l'homme ? »} \\
« La loi n°2024/002 du 19 avril 2024 relative à la décentralisation confie aux collectivités territoriales décentralisées (CTD) la gestion de services publics numériques (état civil, foncier). Pourtant, dans les zones rurales de l'Extrême-Nord ou de l'Adamaoua, l'absence d'infrastructures électriques et numériques transforme cet outil de libération administrative en facteur d'exclusion. La technique, dans ce contexte, libère-t-elle vraiment ou creuse-t-elle les inégalités ? »

\textbf{Sujet : « L'État a-t-il pour seule fonction d'assurer la sécurité ? »} \\
« Depuis 2016, les régions du Nord-Ouest et du Sud-Ouest traversent une crise sécuritaire majeure. Si la réponse de l'État a été d'abord militaire (opérations BIR, couvre-feux), les populations réclament aussi le retour des services publics : écoles, hôpitaux, routes. Cette situation met en lumière la tension entre la fonction régalienne de sécurité et la fonction sociale de l'État. L'État se réduit-il au monopole de la violence légitime ? »
\end{exemplebox}

\paragraph{Temps 2 — Définition des termes \& Rappel du sujet.}
Ne définissez pas \emph{tous} les mots du dictionnaire. Ciblez \textbf{les termes-clés porteurs d'ambiguïté philosophique} (ex: \textit{aliénation, technique, nature, liberté, justice, bonheur, État}). Donnez une définition \cle{philosophique} (pas seulement lexicale), en précisant le sens que \textbf{vous} retenez pour la dissertation. Reformulez ensuite le sujet pour montrer que vous en avez saisi la portée exacte (ex: « Le sujet ne demande pas si le travail est dur, mais s'il structurellement aliène l'homme »).

\paragraph{Temps 3 — La Problématique : le cœur nucléaire.}
C'est la \cle{question unique, tensionnelle et féconde} qui guide tout le développement. Elle ne doit pas être une simple reformulation interrogative du sujet (« Le travail est-il une aliénation ? » $\to$ mauvaise problématique). Elle doit exhiber une \textbf{tension entre deux thèses opposées plausibles}.
\begin{itemize}
    \item \textbf{Mauvaise :} « Le travail est-il une aliénation ? » (Trop binaire, appelle oui/non).
    \item \textbf{Bonne :} « Dans quelle mesure la contrainte inhérente au travail (nécessité, division, salariat) s'oppose-t-elle à son pouvoir de réalisation humaine (création, socialisation, autonomie), et comment dépasser cette opposition ? »
\end{itemize}

\paragraph{Temps 4 — L'Annonce du Plan : la carte d'identité du devoir.}
Annoncez \textbf{explicitement} les trois grandes parties (I, II, III) avec leurs titres exacts (ou une phrase résumant la thèse de chaque partie). Pas de « Nous verrons d'abord... ensuite... enfin... ». Écrivez : \textbf{I. Le travail comme aliénation nécessaire (contrainte, division, exploitation). II. Le travail comme réalisation de soi (création, maîtrise, reconnaissance). III. Vers un travail émancipé : conditions juridiques, techniques et éthiques.}

\begin{attention}[Erreur fatale : L'introduction rédigée \textit{après} le développement]
Beaucoup d'élèves rédigent l'introduction \textbf{en premier} sur la copie propre. C'est un piège mortel : au fil du développement, la problématique s'affine, le plan change, des définitions s'ajustent. L'introduction recopiée en premier contredit alors le développement.
\begin{center}
\textbf{RÈGLE D'OR :} Au brouillon, vous rédigez l'introduction \textbf{EN DERNIER} (après avoir bouclé le plan détaillé et les arguments). Sur la copie propre, vous la recopiez \textbf{EN PREMIER}.
\end{center}
\end{attention}

\begin{aretenir}[Le poids des vitrines : 4 à 6 points sur 20]
Au Probatoire/Bac Tech, la grille de correction réserve typiquement :
\begin{itemize}
    \item \textbf{Introduction (2 à 3 pts) :} Accroche pertinente (0,5), Définitions claires (0,5), Problématique tensionnelle (1), Plan annoncé et cohérent (1).
    \item \textbf{Conclusion (2 à 3 pts) :} Bilan synthétique répondant à la problématique (1), Ouverture pertinente (1), Chute stylistique (0,5 à 1).
\end{itemize}
\textbf{Conséquence :} Une introduction ratée (pas de problématique, plan flou) plafonne souvent la note globale à \textbf{10/20 maximum}, même avec un bon développement. Les correcteurs « jugent » la copie en lisant l'intro et la conclusion en premier (\textit{voir boîte « Le saviez-vous ? »}).
\end{aretenir}

\courssub{La conclusion : le triptyque Bilan — Ouverture — Chute}

La conclusion n'est pas un résumé. C'est l'\textbf{acte de jugement final}. Elle suit un mouvement inverse de l'entonnoir : elle part de la réponse précise (Bilan) pour s'ouvrir sur une perspective plus large (Ouverture) et se refermer sur une phrase forte (Chute).

\begin{popfigure}
\begin{tikzpicture}[scale=0.95, every node/.style={font=\small}]
% Styles
\tikzset{boxstyle/.style={rounded corners, draw, thick, fill=popcream, inner sep=6pt, align=center, minimum width=5cm}}
\tikzset{arrowstyle/.style={->, thick, popdark, shorten >=2pt, shorten <=2pt}}
% Nœuds
\node[boxstyle, fill=popblue!20, draw=popblue, align=center] (bilan) at (0, 2){\textbf{BILAN SYNTHÉTIQUE}\\(Flèche convergente $\downarrow$)\\\textit{Réponse directe à la problématique}\\\textit{Sans nouveau argument}};
\node[boxstyle, fill=poporange!20, draw=poporange, align=center] (ouverture) at (0, 0){\textbf{OUVERTURE CONTRÔLÉE}\\(Flèche divergente $\nearrow$)\\\textit{Autre dimension, Auteur, Actualité,}\\\textit{Lien programme (ex: Technique $\to$ Travail)}};
\node[boxstyle, fill=popgreen!20, draw=popgreen, align=center] (chute) at (0, -2){\textbf{PHRASE DE CHUTE}\\(Point final brillant $\bullet$)\\\textit{Formule marquante, image forte}\\\textit{« Le travail n'est pas le prix de la vie, mais son œuvre »}};
% Flèches
\draw[arrowstyle] (bilan.south) -- (ouverture.north) node[midway, right, font=\tiny, popmuted] {Élargissement};
\draw[arrowstyle] (ouverture.south) -- (chute.north) node[midway, right, font=\tiny, popmuted] {Scellement};
% Décoration horizon
\draw[popmuted, dashed, thick] (-4, -3) -- (4, -3) node[midway, below, font=\tiny] {Horizon de la pensée};
\end{tikzpicture}
\legende{Modèle visuel de la conclusion idéale : convergence vers la réponse (Bilan), divergence vers l'horizon (Ouverture), point final scintillant (Chute).}
\end{popfigure}

\paragraph{1. Le Bilan Synthétique (La réponse).}
C'est la \textbf{réponse directe à votre problématique}, formulée en une ou deux phrases denses. Elle synthétise la dialectique du développement (Thèse $\to$ Antithèse $\to$ Synthèse). \textbf{Interdiction formelle :} apporter un nouvel argument, une nouvelle citation, un nouvel exemple. Le procès est clos.

\paragraph{2. L'Ouverture Contrôlée (L'élargissement).}
C'est le moment où vous montrez que votre réflexion ne s'arrête pas aux frontières du sujet. Mais attention : ouverture \textbf{ne veut pas dire} « hors-sujet » ni « fourre-tout ».

\begin{definition}[Ouverture = Élargissement contrôlé]
L'ouverture doit être \textbf{logiquement reliée} au bilan. Elle peut prendre quatre formes canoniques au programme de Première Tech :
\begin{enumerate}
    \item \textbf{Perspective historique :} « Cette conception du travail comme émancipation, défendue par Marx, sera remise en cause par Hannah Arendt dans \textit{La Condition de l'homme moderne} (1958), qui distingue \textit{labor} (contrainte) et \textit{work} (fabrication). »
    \item \textbf{Comparaison culturelle / Anthropologique :} « Au-delà du salariat occidental, la notion bantoue de \textit{Ngong} (entraide communautaire) ou le \textit{Tontine} camerounais réinventent le travail comme lien social et non comme contrat. »
    \item \textbf{Enjeu éthique / Technologique actuel :} « L'irruption de l'IA générative (ChatGPT, Midjourney) au Cameroun (startups Silicon Mountain, administration) relance la question : le travail intellectuel est-il encore le propre de l'homme ? »
    \item \textbf{Lien avec une autre notion du programme (Transversalité) :} \textit{Ex: Sujet sur La Technique $\to$ Ouverture sur Le Travail (aliénation/libération) ou La Nature (technique vs environnement) ou La Justice (justice algorithmique).}
\end{enumerate}
\end{definition}

\begin{attention}[Interdits absolus en conclusion]
\begin{itemize}
    \item \textbf{Définir le sujet} (c'est le rôle du Temps 2 de l'intro).
    \item \textbf{Apporter un nouvel argument} ou un nouvel auteur non cité dans le développement.
    \item \textbf{Finir par :} « En conclusion, nous avons vu que... », « Pour conclure, disons que... », « En somme... » (Banal, vide, style « rédaction de collège »).
    \item \textbf{Ouverture hors-sujet :} Sujet sur « La technique » $\to$ Ouverture sur « L'art est-il imitation de la nature ? » (Zéro point). Sujet sur « L'État » $\to$ Ouverture sur « La liberté de la presse » (sans lien explicite avec la fonction étatique).
    \item \textbf{L'ouverture « tiroir » :} « On pourrait aussi étudier... » (Pas une ouverture, une liste de courses).
\end{itemize}
\end{attention}

\paragraph{3. La Phrase de Chute (La signature).}
Une phrase courte, rythmée, imagée, qui reste en mémoire. Elle scelle la réflexion.
\begin{exemplebox}[Chutes types]
\begin{itemize}
    \item \textit{Sur le Travail :} « Le travail n'est pas le prix que l'homme paie pour exister, mais l'œuvre par laquelle il se prouve. »
    \item \textit{Sur la Technique :} « La technique n'est ni bien ni mal en soi ; elle est le miroir de nos choix collectifs. »
    \item \textit{Sur l'État/Justice :} « Un État qui ne protège que les forts n'est pas un État, c'est une forteresse. »
    \item \textit{Sur le Bonheur :} « Le bonheur n'est pas au bout du chemin, il est la manière de le marcher. »
\end{itemize}
\end{exemplebox}

\begin{savaistu}[Le secret des correcteurs : La « lecture en diagonale »]
Au Bac, les correcteurs ont souvent 20 à 30 copies par heure. La très grande majorité lit \textbf{d'abord l'introduction et la conclusion} (parfois seulement celles-ci) pour « juger » le niveau de la copie : y a-t-il une problématique ? Un plan ? Une réponse ? Une ouverture ? Un style ?
\begin{center}
\textbf{Soignez vos vitrines comme une plaidoirie d'avocat :} \\
C'est votre unique garantie d'être lu \textit{en entier} et \textit{bienveillamment}.
\end{center}
\end{savaistu}

\courssub{Atelier d'analyse : Trois introductions types notées et commentées}

\textbf{Sujet imposé :} \textit{« La technique est-elle une menace pour la nature ? »} (Thème : La Technique / La Nature — Programme 1re Tech).

\paragraph{Sujet d'analyse :} Taux de déforestation au Cameroun (Global Forest Watch 2023 : \num{100 000} ha/an), exploitation forestière industrielle vs artisanale, loi forestière 1994/2023, enjeu carbone/REDD+.

\begin{exoresolu}[Introduction A — Note : 16/20 (Excellente)]
\textbf{Énoncé :} 
« En 2023, le Global Forest Watch alertait sur la perte de \num{100 000} hectares de forêt primaire au Cameroun, soit l'équivalent de \num{140 000} terrains de football, principalement due à l'agro-industrie (palmier à huile, hévéa) et à l'exploitation minière (cobalt, or) qui mobilisent des techniques lourdes (bulldozers, produits chimiques). Cette hémorragie verte, dans un pays qui abrite \num{40}\% de la biodiversité du bassin du Congo, pose avec acuité la responsabilité de nos outils. Si la technique, au sens d'ensemble des procédés rationnels pour transformer la matière, a permis d'extraire la richesse, elle semble ici détruire le socle même de la vie. \textbf{La technique, par sa puissance démultiplicatrice, condamne-t-elle irréversiblement la nature à n'être qu'une ressource exploitable, ou peut-elle, réorientée, devenir l'alliée de sa régénération ?} \textbf{I. La technique comme puissance de destruction : artificialisation, pollution, réduction du vivant à ressource. II. La technique comme puissance de connaissance et de protection : surveillance spatiale, biotechnologies, énergies propres. III. Vers une « technique écologique » : éthique de la responsabilité (Jonas), droit de la nature, low-tech camerounaise.} »

\textbf{Grille de commentaire :}
\begin{itemize}
    \item \textbf{Accroche (2/2) :} Chiffre précis (GFW 2023), causes techniques identifiées (agro-industrie, mines), enjeu biodiversité (bassin Congo), lien explicite (« responsabilité de nos outils »).
    \item \textbf{Définitions/Sujet (1,5/2) :} Définition philosophique de la technique (« procédés rationnels... »), reformulation du sujet (nature = ressource vs socle de vie). Manque juste une nuance sur « menace » (risque vs certitude).
    \item \textbf{Problématique (2/2) :} Tension nette : \textit{Condamnation irréversible} (menace) vs \textit{Alliée de régénération} (espoir). Question unique, féconde, orientant le plan.
    \item \textbf{Plan (2/2) :} Trois parties équilibrées, dialectiques (Destruction $\to$ Protection $\to$ Dépassement éthique), titres précis, annonce d'auteurs (Jonas) et d'ancrage local (low-tech).
    \item \textbf{Style/Orthographe (1/1) :} Fluide, vocabulaire précis (« artificialisation », « hémorragie verte »), connecteurs logiques.
\end{itemize}
\tcblower
\textbf{Analyse du correcteur :} « Copie de niveau Bac+/Probatoire haut. L'élève maîtrise les codes. L'ancrage camerounais (bassin Congo, low-tech) prouve la maturité technique. La problématique est le moteur parfait. »
\end{exoresolu}

\begin{exoresolu}[Introduction B — Note : 10/20 (Moyenne — « Le piège de la fausse intro ») ]
\textbf{Énoncé :}
« Depuis l'Antiquité, les hommes utilisent des techniques pour survivre. Au Cameroun, on coupe les arbres pour faire du bois et de l'huile de palme. La nature souffre beaucoup. Le sujet demande si la technique est une menace pour la nature. La technique, c'est l'ensemble des moyens pour atteindre une fin. La nature, c'est tout ce qui n'est pas fait par l'homme. \textbf{Est-ce que la technique détruit la nature ?} \textbf{I. Oui, la technique détruit la nature (pollution, déforestation). II. Non, la technique protège la nature (panneaux solaires, drones). III. Il faut trouver un équilibre.} »

\textbf{Grille de commentaire :}
\begin{itemize}
    \item \textbf{Accroche (0,5/2) :} « Depuis l'Antiquité... » (Banale, hors-sol). « Au Cameroun, on coupe les arbres... » (Vague, pas de chiffre, pas de source, style oral).
    \item \textbf{Définitions (1/2) :} Définitions lexicologiques (Larousse), pas philosophiques. « Nature = non fait par l'homme » (réducteur, exclut l'homme de la nature).
    \item \textbf{Problématique (0,5/2) :} Simple reformulation interrogative du sujet. Binaire (Oui/Non). Pas de tension conceptuelle.
    \item \textbf{Plan (1/2) :} Structure « Pour/Contre/Équilibre » (Plan binaire + synthèse molle). Titres plats, pas d'auteurs, pas d'ancrage.
    \item \textbf{Style (0,5/1) :} Phrases courtes, répétitions, vocabulaire pauvre (« souffre beaucoup »).
\end{itemize}
\tcblower
\textbf{Analyse du correcteur :} « Élève qui a le plan en tête mais rate l'habillage. L'intro ne « tient pas la route » philosophiquement. Note plafonnée à 10/20 car le développement, même correct, sera lu avec suspicion. »
\end{exoresolu}

\begin{exoresolu}[Introduction C — Note : 6/20 (Insuffisante — « Le naufrage ») ]
\textbf{Énoncé :}
« La technique est très importante aujourd'hui. Les téléphones, les voitures, les usines. Au Cameroun, il y a beaucoup de pollution à Douala et Yaoundé. La nature est belle, il y a des animaux. Je vais vous dire si la technique est une menace. \textbf{La technique est-elle une menace pour la nature ?} \textbf{I. Les avantages de la technique. II. Les inconvénients de la technique. III. Mon avis personnel. »}

\textbf{Grille de commentaire :}
\begin{itemize}
    \item \textbf{Accroche (0/2) :} Aucune. Énumération d'objets techniques. Pollution citée sans source ni précision.
    \item \textbf{Définitions (0/2) :} Absentes.
    \item \textbf{Problématique (0/2) :} Copier-coller du sujet. Aucune appropriation.
    \item \textbf{Plan (0/2) :} Plan « Avantages/Inconvénients/Avis » = \textbf{Hors sujet philosophique} (pas de dialectique, pas de concepts). « Mon avis » interdit en philo.
    \item \textbf{Style (0/1) :} Familier, aucune rigueur.
\end{itemize}
\tcblower
\textbf{Analyse du correcteur :} « Copie « technique » (série F, MEC, etc.) qui n'a pas basculé en philosophie. L'intro signe l'échec. Le correcteur anticipe un développement narratif. Note finale souvent < 8/20. »
\end{exoresolu}

\courssub{Entraînement chronométré : 15 minutes intro / 10 minutes conclusion}

La gestion du temps est une compétence technique. Au Bac (4h), vous devez boucler l'intro au brouillon en \textbf{15 min max} (après 1h30 de réflexion/plan) et la conclusion en \textbf{10 min max} (en fin de copie).

\begin{methode}[Protocole d'entraînement « Chrono-Bac » (à faire 1 fois/semaine)]
\begin{enumerate}
    \item \textbf{Choisir un sujet type Bac} (Banque de sujets MINESEC / Annales 2020-2024).
    \item \textbf{Phase 1 — Brouillon (25 min) :} Analyse $\to$ Problématique $\to$ Plan détaillé (arguments, auteurs, exemples camerounais par partie).
    \item \textbf{Phase 2 — Intro Chrono (15 min) :} Minuteur lancé. Rédigez l'intro \textbf{au brouillon} en appliquant l'entonnoir 4 temps. Vérifiez : Accroche sourcée ? Définitions philosophiques ? Problématique tensionnelle ? Plan I/II/III exact ?
    \item \textbf{Phase 3 — Développement (2h15) :} Rédigez le développement sur la copie propre (en recopiant l'intro \textbf{en premier}).
    \item \textbf{Phase 4 — Conclusion Chrono (10 min) :} Minuteur lancé. Au brouillon : Bilan (réponse à la problé.) $\to$ Ouverture (choix 1 des 4 types) $\to$ Chute (phrase travaillée). Recopiez sur la copie.
    \item \textbf{Phase 5 — Relecture (10 min) :} Cohérence Intro/Développement/Conclusion. Orthographe. Noms d'auteurs.
\end{enumerate}
\end{methode}

\begin{exercice}[Entraînement immédiat — Sujet : « L'État doit-il garantir le bonheur des citoyens ? »]
\textbf{Contexte camerounais :} Politique nationale de développement (PND 2020-2030), Objectifs de Développement Durable (ODD), Indice de Bonheur National Brut (Bhoutan) vs PIB, rôle de l'État social (santé, éducation, emploi jeunes).
\begin{enumerate}
    \item \textbf{Au brouillon (15 min) :} Rédigez une introduction complète (4 temps) en appliquant la méthode « Accroche Cameroun ».
    \item \textbf{Au brouillon (10 min) :} Après avoir esquissé un plan (I. État protecteur / II. Limites : liberté, responsabilité individuelle / III. État créateur de conditions / Bonheur partagé), rédigez la conclusion (Bilan, Ouverture, Chute).
    \item \textbf{Auto-évaluation :} Cochez la grille ci-dessous.
\end{enumerate}
\end{exercice}

\begin{center}
\begin{tabularx}{\linewidth}{|l|X|X|}\hline
\textbf{Critères Introduction} & \textbf{Oui / Non / Partiel} & \textbf{Commentaire / Correction} \\ \hline
Accroche : Fait précis + Source (INS, JO, Auteur) + Lien sujet & & \\ \hline
Définitions : Termes-clés (État, Bonheur, Garantir) sens philo & & \\ \hline
Problématique : Tension unique, féconde (pas binaire) & & \\ \hline
Plan : I, II, III titres précis, dialectiques, annoncés & & \\ \hline \hline
\textbf{Critères Conclusion} & \textbf{Oui / Non / Partiel} & \textbf{Commentaire / Correction} \\ \hline
Bilan : Réponse directe à la problématique (synthèse) & & \\ \hline
Ouverture : Type identifié (Historique / Culturel / Éthique / Transversal) + Lien logique & & \\ \hline
Chute : Phrase courte, marquante, sans « En conclusion... » & & \\ \hline
Interdits respectés : Pas de déf, pas de nvx arg, pas de HS & & \\ \hline
\end{tabularx}
\end{center}

\begin{corrige}[Exercice n°1 — Pistes de correction type « Excellence »]
\textbf{Intro (Extrait) :} « Le Rapport National Volontaire (RNV 2023) du Cameroun sur les ODD révèle que seulement \num{38}\% de la population a accès à une protection sociale efficace, tandis que le taux de pauvreté monétaire stagne à \num{37,5}\% (INS 2022). Pourtant, la Constitution de 1996 (Préambule) garantit « le droit à la vie, à la santé, à l'éducation ». Cette fracture entre droit constitutionnel et réalité statistique interroge la portée de l'action étatique. L'État, personne morale de droit public détentrice du monopole de la violence légitime (Weber) et de la puissance publique, « garantir » implique-t-il une obligation de résultat (bonheur fourni) ou de moyen (conditions créées) ? \textbf{Jusqu'où l'obligation étatique de garantir les conditions du bonheur (sécurité, santé, éducation) peut-elle s'étendre sans empiéter sur la liberté individuelle de définir et poursuivre son propre bonheur ?} \textbf{I. L'État comme garant des conditions objectives du bonheur (besoins vitaux, justice, paix). II. L'impossible garantie subjective : le bonheur comme quête personnelle (liberté, responsabilité, diversité des conceptions). III. L'État « facilitateurs » : capabilités (Sen), justice sociale, indicateur de bonheur national (au-delà du PIB).} »

\textbf{Conclusion (Extrait) :} « \textbf{Bilan :} L'État ne peut ni ne doit « fabriquer » le bonheur — ce serait totalitaire — mais il a le devoir absolu de garantir les \textit{capabilités} (Sen) qui le rendent possible : sécurité, santé, éducation, justice, emploi décent. \textbf{Ouverture :} Cette conception rejoindrait l'expérience du « Bonheur National Brut » au Bhoutan, où l'État mesure non le PIB mais le bien-être psychologique, la culture, l'environnement ; une piste pour repenser le PND camerounais ? \textbf{Chute :} « L'État ne donne pas le bonheur, il trace la route ; à chaque citoyen de marcher. » 
\end{corrige}

\begin{aretenir}[Check-list « Vitrines » avant de rendre la copie]
\begin{itemize}
    \item [ ] \textbf{Intro :} Accroche camerounaise sourcée ? Termes définis philo ? Problématique tensionnelle ? Plan I/II/III exact ?
    \item [ ] \textbf{Conclusion :} Bilan = réponse à la problé. ? Ouverture = 1 type canonique + lien ? Chute = phrase forte ?
    \item [ ] \textbf{Cohérence :} Intro et Conclusion disent-elles la même chose (même problé, même plan) ?
    \item [ ] \textbf{Temps :} Intro < 15 min (brouillon) / Conclu < 10 min (brouillon) respectés ?
\end{itemize}
\end{aretenir}

\coursec{L'argumentation philosophique : Preuves, exemples et style}

Après avoir maîtrisé l'architecture générale de la copie — introduction et conclusion — lors de la section précédente, nous pénétrons maintenant dans le « moteur » de la dissertation : le développement. Une introduction brillante et une conclusion percutante ne sont que des vitrines ; c'est dans le corps du texte, paragraphe après paragraphe, que se joue la vérité philosophique. C'est ici que l'élève prouve qu'il ne récite pas un cours, mais qu'il \emph{pense} en temps réel, qu'il articule des concepts, qu'il confronte des thèses et qu'il ancre sa réflexion dans le réel camerounais. L'argumentation n'est pas un ornement : c'est la chair même de la démonstration.

\courssub{La structure standard d'un paragraphe argumenté (niveau sous-partie A, B ou C)}

Chaque grande partie (I, II, III) de votre plan se décompose en sous-parties (A, B, C). Une sous-partie n'est pas un « résumé de cours » : c'est une \cle{unité argumentative autonome} qui défend une \cle{idée force} (ou thèse partielle) en lien direct avec la problématique. Pour obtenir une cohérence implacable, respectez scrupuleusement la séquence en six temps suivante. C'est le « squelette » de tout bon paragraphe philosophique.

\begin{methode}[La méthode des « 6 A » pour chaque sous-partie]
\textbf{Étape 1 — AFFIRMER (Phrase thèse / Idée force).} Ouvrez le paragraphe par une phrase affirmative, claire, qui annonce l'argument principal. Elle répond à la question : « Que démontre cette sous-partie ? ». Exemple : \emph{« Le travail n'est pas seulement une contrainte naturelle, il est le lieu principal de l'humanisation de l'individu. »}

\textbf{Étape 2 — ANALYSER (Explication conceptuelle).} Déployez les concepts clés de la thèse. Définissez, distinguez, montrez la portée. Pourquoi cette affirmation a-t-elle du sens ? Exemple : \emph{« Contrairement à l'animal qui s'adapte au milieu, l'homme, par le travail, transforme la nature et se transforme lui-même (Marx). Le travail objectifie l'intentionnalité humaine. »}

\textbf{Étape 3 — ARGUMENTER (Argument logique).} Apportez la preuve rationnelle. C'est le « parce que ». Utilisez des connecteurs logiques (\emph{donc, dès lors, par conséquent, en effet, toutefois}). Exemple : \emph{« Dès lors, priver l'homme de travail créateur, c'est l'empêcher de réaliser son essence de sujet libre et productif. »}

\textbf{Étape 4 — ANCRER (Exemple concret camerounais + Contre-exemple éventuel).} Ancrez l'abstraction dans le réel observable. Voir la typologie détaillée ci-dessous. Exemple : \emph{« Au Cameroun, le secteur informel (\"buyam-sellam\"), bien que précaire, illustre cette capacité d'invention : le commerçant transforme l'espace urbain en lieu d'échange et de lien social, humanisant ainsi un environnement hostile. »}

\textbf{Étape 5 — AUTORISER (Référence philosophique : Auteur / Concept).} Citez un auteur pour \emph{appuyer} ou \emph{contester} votre thèse, jamais pour la remplacer. Exemple : \emph{« Comme l'analyse Hannah Arendt dans \textit{La Condition de l'homme moderne}, l'\textit{animal laborans} est enfermé dans la répétition biologique, seul l'\textit{homo faber} et l'homme d'action (\textit{zoon politikon}) accèdent à la liberté par l'œuvre et la parole. »}

\textbf{Étape 6 — ARTICULER (Mini-bilan / Transition).} Refermez le paragraphe en montrant ce qu'il a apporté à la problématique et ouvrez sur l'argument suivant. Exemple : \emph{« Ainsi, le travail libère par sa dimension créatrice. Mais cette libération reste-t-elle effective lorsque le travail est aliéné par la division capitaliste ? C'est ce que nous examinerons maintenant. »}
\end{methode}

\begin{popfigure}
\begin{tikzpicture}[
    boxstyle/.style={rectangle, rounded corners=3pt, draw=popdark, thick, minimum height=2.5cm, text width=2.2cm, align=center, font=\small\sffamily},
    title/.style={font=\bfseries\small\sffamily, text=white, minimum height=0.7cm, align=center},
    arrowstyle/.style={->, >=Stealth, thick, popdark, shorten <=2pt, shorten >=2pt}
]
% Couleurs par étape
\colorlet{c1}{popblue!80}
\colorlet{c2}{poppurple!80}
\colorlet{c3}{popgreen!80}
\colorlet{c4}{poporange!80}
\colorlet{c5}{poppink!80}
\colorlet{c6}{popgold!80}
% Matrice pour alignement parfait
\matrix (M) [row sep=0.2cm, column sep=0.4cm] {
% Ligne titres
\node[boxstyle, fill=c1, title, fill=c1, align=center] (t1){1. AFFIRMER\\ \emph{(Idée force)}}; &
\node[boxstyle, fill=c2, title, fill=c2, align=center] (t2){2. ANALYSER\\ \emph{(Concepts)}}; &
\node[boxstyle, fill=c3, title, fill=c3, align=center] (t3){3. ARGUMENTER\\ \emph{(Logique)}}; &
\node[boxstyle, fill=c4, title, fill=c4, align=center] (t4){4. ANCRER\\ \emph{(Exemple CM)}}; &
\node[boxstyle, fill=c5, title, fill=c5, align=center] (t5){5. AUTORISER\\ \emph{(Auteur)}}; &
\node[boxstyle, fill=c6, title, fill=c6, align=center] (t6){6. ARTICULER\\ \emph{(Transition)}}; \\
% Ligne contenus
\node[boxstyle, fill=c1!70!white, align=center] (b1){Phrase thèse claire\\ Réponse partielle au sujet}; &
\node[boxstyle, fill=c2!70!white, align=center] (b2){Définition, distinction,\\ portée conceptuelle}; &
\node[boxstyle, fill=c3!70!white, align=center] (b3){Raisonnement :\\ Pourquoi c'est vrai ?}; &
\node[boxstyle, fill=c4!70!white, align=center] (b4){Fait social, chiffre,\\ droit, littérature CM}; &
\node[boxstyle, fill=c5!70!white, align=center] (b5){Citation pour appuyer\\ ou contester}; &
\node[boxstyle, fill=c6!70!white, align=center] (b6){Mini-bilan +\\ Ouverture sur suite}; \\
};
% Flèches entre les boîtes du bas (contenus)
\draw[arrowstyle] (b1.east) -- (b2.west);
\draw[arrowstyle] (b2.east) -- (b3.west);
\draw[arrowstyle] (b3.east) -- (b4.west);
\draw[arrowstyle] (b4.east) -- (b5.west);
\draw[arrowstyle] (b5.east) -- (b6.west);
% Titre global
\node[font=\bfseries\large\sffamily, text=popdark, anchor=south] at ([yshift=0.5cm]M.north) {Modèle standard d'un paragraphe argumenté (Sous-partie A, B ou C)};
\end{tikzpicture}
\legende{Schéma de la structure en 6 temps (Méthode des 6 A). Chaque boîte correspond à une étape obligatoire. Le respect de cet ordre garantit la rigueur logique et la densité philosophique.}
\end{popfigure}

\courssub{Typologie des exemples au Cameroun : ancrer sa pensée dans le réel}

L'exemple philosophique n'est pas une illustration décorative ni une anecdote personnelle. C'est un \cle{fait social structurant} qui manifeste, dans la réalité observable, la thèse que vous défendez. Au Cameroun, nous disposons d'un réservoir inépuisable de matériaux concrets. Les correcteurs du Probatoire/Baccalauréat technique attendent que l'élève technique mobilise ce contexte avec précision.

\begin{definition}[Exemple philosophique $\neq$ Anecdote]
L'exemple philosophique doit être \textbf{représentatif} d'une structure universelle, pas un cas isolé.
\begin{itemize}
    \item \textbf{Anecdote (Zéro point) :} \emph{« Mon oncle est fonctionnaire à Yaoundé et il arrive en retard. »} $\rightarrow$ Subjectif, non généralisable, pathétique.
    \item \textbf{Fait social structurant (Points pleins) :} \emph{« La fonction publique camerounaise comptait environ 300\,000 agents en 2023 (source MINFOPRA) pour une population de 28 millions d'habitants, posant la question de l'efficacité de l'État-providence. »} $\rightarrow$ Chiffré, sourcé, porteur de sens sociologique.
\end{itemize}
L'exemple doit permettre au correcteur de \emph{voir} le concept à l'œuvre dans la société camerounaise.
\end{definition}

\begin{attention}[L'erreur de la confusion Auteur / Fait (Preuve d'autorité mal placée)]
Ne confondez jamais \textbf{référence d'auteur} et \textbf{exemple de fait}.
\begin{itemize}
    \item \textbf{Faux :} \emph{« Comme le dit Marx, le prolétariat est exploité. »} $\rightarrow$ Marx donne une \emph{interprétation théorique}, pas un chiffre du chômage au Cameroun en 2024.
    \item \textbf{Vrai :} \emph{« L'exploitation structurelle décrite par Marx trouve un écho dans le secteur informel camerounais où 90\% des emplois (INS, 2022) ne bénéficient d'aucune protection sociale (CNPS). »} $\rightarrow$ L'auteur éclaire le fait ; le fait (chiffre INS) prouve la réalité.
\end{itemize}
\textbf{Règle d'or :} Les faits bruts (stats, lois, histoire) viennent de la réalité ; les auteurs viennent pour \emph{conceptualiser} ces faits.
\end{attention}

Voici une « banque d'exemples » classée par notions du programme, à mémoriser et à adapter :

\begin{popfigure}
\begin{tikzpicture}[
    cell/.style={rectangle, draw=popline, minimum height=1.1cm, align=center, font=\small\sffamily, text width=3.5cm},
    head/.style={rectangle, fill=popblue, text=white, font=\bfseries\small\sffamily, minimum height=0.9cm, align=center, text width=3.5cm},
    subhead/.style={rectangle, fill=popblueL, font=\bfseries\small\sffamily, minimum height=0.7cm, align=center, text width=3.5cm},
    notion/.style={rectangle, fill=popgreenL, font=\bfseries\small\sffamily, minimum height=0.7cm, align=center, text width=2.5cm},
    ex/.style={rectangle, fill=popcream, font=\small\sffamily, minimum height=1.1cm, align=left, text width=5.5cm}
]
% En-têtes
\node[head] (h1) at (0,0) {NOTION DU PROGRAMME};
\node[head] (h2) at (3.5,0) {EXEMPLES CAMEROUNAIS TYPES (Faits, Chiffres, Textes, Œuvres)};
\node[head] (h3) at (9.0,0) {CONCEPTS / AUTEURS ASSOCIÉS};
% Ligne 1 : TRAVAIL
\node[notion] (n1) at (0,-0.9) {TRAVAIL};
\node[ex] (e1) at (3.5,-0.9) {\textbf{Secteur informel :} « Buyam-sellam », moto-taxis (bensikin), tontines. \textbf{Chiffres :} 90\% emploi informel (INS). \textbf{Droit :} Code du travail 1992 (art. 1er), Convention C189 OIT (travail décent). \textbf{Littérature :} \textit{Ville cruelle} (Eza Boto / Mongo Beti) — condition ouvrière.};
\node[cell] (c1) at (9.0,-0.9) {Aliénation (Marx), Travail libre (Hegel), Condition de l'homme moderne (Arendt), Droit du travail};
% Ligne 2 : ÉTAT
\node[notion] (n2) at (0,-2.0) {ÉTAT / POLITIQUE};
\node[ex] (e2) at (3.5,-2.0) {\textbf{Institutions :} Décentralisation (CTD 2019), Chefferies traditionnelles (loi 77/01). \textbf{Chiffres :} 360 communes, 10 régions. \textbf{Texte :} Constitution 1996 (Préambule, Art. 1, 55, 62). \textbf{Art :} \textit{La Palabre} (Hountondji) comme modèle démocratique africain.};
\node[cell] (c2) at (9.0,-2.0) {Contrat social (Rousseau, Hobbes), Légitimité (Weber), Souveraineté, Démocratie, Palabre (Hountondji)};
% Ligne 3 : TECHNIQUE
\node[notion] (n3) at (0,-3.1) {TECHNIQUE / PROGRÈS};
\node[ex] (e3) at (3.5,-3.1) {\textbf{Agriculture :} Moderne (SAA, CDC, Plantations du Haut-Penja) vs Traditionnelle (agriculture itinérante sur brûlis, polyculture). \textbf{Nucléaire :} Projet de centrale (MINEPAT). \textbf{Numérique :} « Silicon Mountain » (Buea), fintechs (Djangui, Maviance). \textbf{Auteur :} Francis Bebey (musique électronique africaine).};
\node[cell] (c3) at (9.0,-3.1) {Technique comme culture (Leroi-Gourhan), Progrès (Condorcet), Aliénation technique (Ellul, Heidegger), Maître/Esclave (Hegel)};
% Ligne 4 : MORALE
\node[notion] (n4) at (0,-4.2) {MORALE / SOCIÉTÉ};
\node[ex] (e4) at (3.5,-4.2) {\textbf{Corruption :} « Gombo », « Arrangement », CONAC (rapports annuels). \textbf{Solidarité :} Tontines, associations de développement (ASD), « Njangi ». \textbf{Droit :} Code pénal (art. 134), Stratégie nationale anticorruption. \textbf{Littérature :} \textit{Le Prix de la liberté} (Calixthe Beyala), \textit{Les Boutiques de l'égalité} (Bebey).};
\node[cell] (c4) at (9.0,-4.2) {Devoir (Kant), Utilitarisme (Bentham/Mill), Vertu (Aristote), Corruption/Intégrité, Communautarisme (Mbiti, Hountondji)};
% Ligne 5 : CONNAISSANCE
\node[notion] (n5) at (0,-5.3) {CONNAISSANCE / VÉRITÉ};
\node[ex] (e5) at (3.5,-5.3) {\textbf{Éducation :} Taux d'alphabétisation (INS/MINEDUB), Bilinguisme officiel (Art. 1er Const.), Enseignement technique vs général. \textbf{Savoirs endogènes :} Médecine traditionnelle (OMS/Cameroun), savoirs agricoles locaux. \textbf{Auteur :} Paulin Hountondji (\textit{Sur la philosophie africaine}, « Dépendance épistémologique »).};
\node[cell] (c5) at (9.0,-5.3) {Vérité (Descartes, Nietzsche), Science/Pseudo-science (Popper), Relativisme/Universalisme, Épistémologie africaine (Hountondji, Wiredu)};
% Cadre global
\draw[popdark, thick] (0,0.45) rectangle (12.5,-5.85);
\draw[popline] (0,-0.45) -- (12.5,-0.45);
\draw[popline] (3.5,0.45) -- (3.5,-5.85);
\draw[popline] (9.0,0.45) -- (9.0,-5.85);
\foreach \y in {-0.9,-2.0,-3.1,-4.2,-5.3} \draw[popline] (0,\y-0.45) -- (12.5,\y-0.45);
\end{tikzpicture}
\legende{Banque d'exemples camerounais classés par notions programmatiques. Chaque entrée combine : fait social, donnée chiffrée (INS, MINEDUB, MINEPAT, CONAC), texte juridique (Constitution, Codes), œuvre littéraire/artistique. À croiser avec les concepts/auteurs de la colonne de droite.}
\end{popfigure}

\courssub{L'utilisation des auteurs : dialoguer avec la tradition philosophique}

Citer un auteur n'est pas un exercice de « name-dropping » (étalage de noms) pour faire savant. La référence philosophique a deux fonctions strictes, codifiées dans la grille d'évaluation :

\begin{propriete}[Les deux usages légitimes de la citation d'auteur]
\textbf{1. L'apport (Appui) :} L'auteur formule mieux que vous la thèse que vous défendez. Vous l'utilisez comme \emph{autorité conceptuelle}.
\begin{quote}
\emph{« Comme le souligne Kant dans \textit{Fondements de la métaphysique des mœurs}, « Agis uniquement d'après la maxime qui fait que tu puisses vouloir en même temps qu'elle devienne une loi universelle. » Le devoir moral ne supporte aucune exception empirique. »}
\end{quote}

\textbf{2. Le contrepoint (Contestation) :} Vous exposez une thèse adverse pour la réfuter ou la nuancer. Vous montrez que vous maîtrisez le débat.
\begin{quote}
\emph{« Contrairement à Hobbes pour qui l'état de nature est une « guerre de tous contre tous » (\textit{Léviathan}, ch. 13), Rousseau soutient dans le \textit{Discours sur l'inégalité} que l'homme naturel est bon et que la société le corrompt. C'est cette seconde thèse que nous retenons pour penser l'éducation... »}
\end{quote}

\textbf{Interdit :} \emph{« Kant dit que... Donc c'est vrai. »} (Argument d'autorité fallacieux). \emph{« Descartes, Kant, Hegel, Marx, Sartre pensent que... »} (Liste de courses sans articulation).
\end{propriete}

\begin{exemplebox}[Profil type d'une copie 14/20 au Probatoire/Bac Tech]
Une analyse des copies excellentes (jurys 2022-2024) montre des constantes quantitatives :
\begin{itemize}
    \item \textbf{12 à 15 références explicites} (Auteurs nommés + Concepts précis : \emph{cogito, surhomme, praxis, catégorique, historicité, palabre...}) réparties sur les 3 grandes parties.
    \item \textbf{6 à 8 exemples concrets variés}, dont \textbf{3 à 4 spécifiquement camerounais} (chiffres INS, articles Constitution, faits sociaux « buyam-sellam/Ngondo/tontines », auteurs camerounais/africains : Mongo Beti, Beyala, Hountondji, Mudimbe, Bebey).
    \item \textbf{0 citation « orpheline »} : chaque auteur cité est \emph{explicité} (on explique sa thèse) et \emph{articulé} (on dit pourquoi il est là).
\end{itemize}
La qualité prime sur la quantité : 3 auteurs bien utilisés valent mieux que 10 noms jetés en vrac.
\end{exemplebox}

\courssub{La qualité du style philosophique : rigueur et impersonnalité}

Le style n'est pas l'enveloppe de la pensée, il \emph{est} la pensée en train de se faire. Une phrase mal construite trahit un concept mal maîtrisé. Le jury sanctionne l'approximation lexicale et la subjectivité non assumée.

\begin{aretenir}[Vocabulaire de précision : distinctions exigibles au Bac]
\begin{center}
\begin{tabularx}{\linewidth}{|l|X|X|}
\hline
\textbf{Concept} & \textbf{Sens philosophique rigoureux} & \textbf{Contre-sens fréquent (Zéro)} \\ \hline
Aliénation & Perte de soi dans l'objet produit (Marx), aliénation juridique (transfert de propriété) & Exploitation, simple souffrance au travail \\ \hline
Exploitation & Rapport structurel : appropriation du surtravail (Marx) & Mauvais traitements, bas salaire \\ \hline
Technique & Ensemble de moyens rationnels pour une fin (Leroi-Gourhan) & Machines, outils, « la technologie » \\ \hline
Progrès & Amélioration cumulative, réversible, évaluable (Condorcet) & Évolution, changement, modernité \\ \hline
État & Monopole de la violence physique légitime (Weber), personne morale de droit public & Gouvernement, administration, pays \\ \hline
Démocratie & Régime où le peuple est souverain (étymologie + Rousseau/Schmitt) & Liberté d'expression, élections seules \\ \hline
\end{tabularx}
\end{center}
\end{aretenir}

\begin{methode}[Check-list style « Zéro faute » pour la rédaction]
\begin{enumerate}
    \item \textbf{Impersonnalité :} Bannir « Je pense que », « À mon avis », « Selon moi ». Préférer : \emph{« Il apparaît que... », « Force est de constater que... », « L'analyse montre que... », « Il convient de noter que... »}.
    \item \textbf{Connecteurs logiques variés :} Ne pas abuser de « Donc ». Utiliser : \emph{Par conséquent, dès lors, en conséquence, corollaire, en revanche, toutefois, néanmoins, par ailleurs, de surcroît, qui plus est, en outre, c'est-à-dire, autrement dit}.
    \item \textbf{Phrases complètes :} Sujet + Verbe + Complément. Pas de phrases nominales style télégraphique (sauf titres de parties).
    \item \textbf{Présent de vérité générale :} \emph{« Le travail humanise l'homme »} (pas « humanisait » ni « humanisera »).
    \item \textbf{Guillemets pour citations :} Citation courte ($<3$ lignes) dans le texte entre guillemets français « ... ». Citation longue ($>3$ lignes) en bloc indenté, sans guillemets, police plus petite.
\end{enumerate}
\end{methode}

\courssub{Gestion du temps et auto-correction guidée : les 3 heures du Bac}

L'examen du Probatoire / Baccalauréat technique dure \textbf{3 heures (180 minutes)}. Une gestion rigoureuse est la condition \emph{sine qua non} de la réussite. Ne laissez pas le temps vous gérer.

\begin{methode}[Chrono-type optimal (180 min)]
\begin{enumerate}
    \item \textbf{00:00 -- 00:30 (30 min) : ANALYSE \& PLAN \& BROUILLON INTRO/CONCL.}
    \begin{itemize}
        \item 10 min : Analyse du sujet (termes, type, présupposés, problématique).
        \item 15 min : Plan détaillé (I, II, III + A, B, C + idées forces + exemples + auteurs).
        \item 5 min : Rédaction au brouillon de l'\textbf{Introduction complète} (accroche, définition, problématique, annonce de plan) et de la \textbf{Conclusion complète} (réponse synthétique, ouverture). \emph{Astuce :} Recopiez-les en propre dès le début de la rédaction pour sécuriser les 4-5 points « faciles ».
    \end{itemize}
    \item \textbf{00:30 -- 02:45 (2h15) : RÉDACTION DU DÉVELOPPEMENT.}
    \begin{itemize}
        \item $\approx$ 45 min par grande partie (I, II, III).
        \item Écrivez \emph{directement sur la copie} en suivant votre plan brouillon. Appliquez la structure « 6 A » par sous-partie.
        \item Si vous bloquez sur un argument : passez au suivant, laissez de l'espace, revenez après.
    \end{itemize}
    \item \textbf{02:45 -- 03:00 (15 min) : RELECTURE ACTIVE \& CORRECTION.}
    \begin{itemize}
        \item Relisez \emph{à voix basse} (bouche les mots) pour entendre les phrases bancales.
        \item Appliquez la grille d'auto-correction ci-dessous.
        \item Vérifiez : Numérotation parties/sous-parties, noms d'auteurs/titres (italiques), accents, accords, lisibilité.
    \end{itemize}
\end{enumerate}
\end{methode}

\begin{attention}[Grille d'auto-correction : 4 pièges mortels à traquer en relecture]
\begin{enumerate}
    \item \textbf{La phrase « vide » (remplissage) :} \emph{« Le travail est très important dans la vie de l'homme depuis la nuit des temps. »} $\rightarrow$ \textbf{Supprimez.} Remplacez par une thèse argumentée.
    \item \textbf{L'assertion non prouvée :} \emph{« La technique nous rend libres. »} $\rightarrow$ \textbf{Pourquoi ?} Ajoutez l'argument logique (Étape 3) et l'exemple (Étape 4).
    \item \textbf{La répétition lexicale :} \emph{« Le travail... le travail... le travail... »} $\rightarrow$ Utilisez pronoms relatifs, synonymes (\emph{l'activité productive, l'œuvre, le labeur, cette praxis...}).
    \item \textbf{La faute de français grave :} Accord sujet/verbe (surtout avec inversion), confusion \emph{a/à, et/est, ou/où, son/sont, ce/se, ces/ses}, participe passé (avoir/être). \textbf{1 faute grave = -0.5pt en langue.} 5 fautes = note de langue plafonnée.
\end{enumerate}
\end{attention}

\courssub{Transfert : La dissertation hors examen (Vie professionnelle \& Citoyenne)}

La dissertation philosophique n'est pas un exercice scolaire artificiel. C'est un \emph{entraînement à la pensée critique et à la communication structurée}, compétences transversales indispensables pour un technicien supérieur, un cadre ou un citoyen responsable.

\begin{savaistu}[Paulin Hountondji et la « dépendance épistémologique »]
Le philosophe béninois Paulin Hountondji (\textit{Sur la philosophie africaine}, 1973) dénonce la \textbf{dépendance épistémologique} : l'intellectuel africain qui ne fait qu'appliquer des concepts forgés ailleurs (Marx, Heidegger, Foucault) à des réalités africaines sans les critiquer ni en produire de nouveaux. Pour Hountondji, philosopher en Afrique, c'est \textbf{produire ses propres concepts à partir de nos réalités}.
\begin{itemize}
    \item \textbf{Exemple :} La \emph{Palabre} (ou \emph{Palaver}) n'est pas juste une « discussion ». C'est une \textbf{procédure démocratique africaine} de résolution des conflits par la parole, le consensus et la réconciliation, distincte du vote majoritaire occidental (majorité/minorité). C'est un concept opératoire pour penser la démocratie au Cameroun (chefferies, conseil municipal, dialogue national).
    \item \textbf{Leçon pour vous :} Quand vous mobilisez « Buyam-sellam », « Tontine », « Ngondo », « Chefferie » comme exemples philosophiques, vous faites exactement ce que Hountondji demande : vous \emph{conceptualisez le réel camerounais}. C'est la marque d'une pensée technique \emph{autonome} et \emph{ancrée}.
\end{itemize}
\end{savaistu}

\begin{exoresolu}[Dissection d'un paragraphe type (Sujet : « Le travail est-il une aliénation ? »)]
\textbf{Énoncé :} Rédiger la sous-partie A du I : « Le travail comme aliénation dans le capitalisme (Marx) ».

\textbf{Solution détaillée (application des 6 A) :}
\begin{enumerate}
    \item \textbf{AFFIRMER :} \emph{« Dans le mode de production capitaliste, le travail salarié devient la forme achevée de l'aliénation du producteur. »} (Idée force claire, liée au sujet).
    \item \textbf{ANALYSER :} \emph{« L'aliénation (Entfremdung) désigne ici le processus par lequel l'ouvrier perd la maîtrise de son activité (le travail est imposé, pas choisi), de son produit (qui appartient au capitaliste), de sa nature d'espèce (Gattungswesen : l'homme créateur devient bête de somme) et des autres hommes (concurrence). »} (Concepts marxistes précis : travail imposé, appropriation du produit, déshumanisation, rapport social).
    \item \textbf{ARGUMENTER :} \emph{« Puisque le capitaliste achète la force de travail et en dispose comme d'une marchandise, le travailleur vend son temps de vie. Le surtravail (plus-value) est extorqué. Le travailleur ne se réalise pas dans son travail, il s'y nie ; il ne s'y développe pas, il s'y atrophie. »} (Logique interne du capital : marchandise force de travail $\rightarrow$ plus-value $\rightarrow$ nihilification).
    \item \textbf{ANCRER (Cameroun) :} \emph{« Au Cameroun, les ouvriers des plantations agro-industrielles (CDC, PHP, SOCAPALM) ou les mineurs de Capam (or) subissent souvent des cadences infernales, une pénibilité extrême, pour des salaires proches du SMIG (41 875 FCFA depuis 2023), sans sécurité de l'emploi (sous-traitance). Le produit de leur sueur (banane, huile de palme, or) enrichit des actionnaires souvent étrangers, illustrant l'expropriation du produit. »} (Fait social précis, chiffres, noms d'entreprises, lien avec thèse).
    \item \textbf{AUTORISER :} \emph{« C'est le cœur de la critique de l'économie politique chez Marx (\textit{Manuscrits de 1844}, \textit{Le Capital} Livre I) : « L'ouvrier ne se sent pas chez lui dans son travail... son travail n'est pas volontaire mais forcé, travail forcé. » »} (Citation exacte, référence précise, utilisée pour \emph{nommer} la structure décrite).
    \item \textbf{ARTICULER :} \emph{« Ainsi, le travail salarié capitaliste structurellement aliène le producteur. Mais cette aliénation est-elle une fatalité ontologique du travail comme tel, ou seulement une forme historique dépassable ? C'est la thèse de la réappropriation que nous défendrons en B. »} (Mini-bilan + Transition parfaite vers la thèse inverse).
\end{enumerate}
\textbf{Note :} Ce paragraphe fait $\approx$ 180 mots. Une partie (I) = 3 tels paragraphes (A, B, C) $\approx$ 550 mots. 3 parties = 1650 mots. + Intro/Concl = 1900 mots. Cible idéale : 1800--2200 mots.
\end{exoresolu}

\begin{exercice}[Entraînement : Construire un paragraphe « 6 A »]
\textbf{Sujet :} « La technique libère-t-elle l'homme ? »
\textbf{Consigne :} Rédigez la sous-partie B du II : « La technique comme amplification de la liberté (thèse prométhéenne) ».
\textbf{Contraintes :} Respecter les 6 étapes. Mobiliser \textbf{un auteur} (Prométhée / Platon / Marx / Heidegger / Ellul / Leroi-Gourhan / Simondon) et \textbf{un exemple camerounais concret} (Agriculture moderne, Numérique « Silicon Mountain », Énergie, Transport). \textbf{Durée : 20 min chrono.}
\end{exercice}

\begin{corrige}[Exercice 1]
\textbf{Piste de rédaction (structure) :}
\begin{enumerate}
    \item \textbf{Thèse :} La technique, en maîtrisant les forces naturelles, libère l'homme de la contrainte biologique et ouvre le champ des possibles.
    \item \textbf{Analyse :} Distinction technique/outil. La technique est savoir-faire rationnel (Leroi-Gourhan : « l'extériorisation de la mémoire et du calcul »). Elle transforme le milieu.
    \item \textbf{Argument :} Moins de temps pour la survie (chasse, cueillette, eau, bois) = plus de temps libre pour la culture, la politique, l'art (Arendt : passage du \emph{labor} à l'\emph{work} et à l'\emph{action}).
    \item \textbf{Exemple CM :} L'adoption de l'irrigation goutte-à-goutte et semences améliorées dans le Nord (MINEPAT/PNUD) libère les paysans de la dépendance pluviométrique, augmente le rendement (trois fois), permet diversification (maraîchage) et scolarisation des enfants. Ou : « Silicon Mountain » (Buea) : le numérique permet à de jeunes ingénieurs de créer des startups (fintech, agritech, edtech) exportant des services, libérés de la contrainte géographique.
    \item \textbf{Auteur :} Prométhée (mythe, Platon \textit{Protagoras}) : don du feu/technique pour compenser la nudité humaine. Ou Marx (\textit{Idéologie allemande}) : « La production des moyens de subsistance... est la base de toute histoire. » Ou Simondon : technique comme « concrétisation » libérant l'homme de l'asservissement à la matière.
    \item \textbf{Transition :} Mais cette libération technique est-elle sans ombre ? Le risque d'aliénation technique (Ellul/Heidegger) appelle une éthique de la responsabilité (Jonas) -- voir sous-partie C.
\end{enumerate}
\end{corrige}

\coursec{Atelier de consolidation : Sujets types Bac \& Grille d'auto-évaluation}

Après avoir maîtrisé les outils de l'argumentation — preuves, exemples, style — il est temps de passer à l'entraînement grandeur nature. Cette section transforme la classe en « salle d'examen » et en « jury de délibération ». L'objectif : automatiser les bons réflexes, identifier ses failles grâce à la grille officielle MINESEC, et repartir avec une fiche mémo opérationnelle pour le jour J.

\courssub{Sujet 1 — Simulation Bac complète : « La technique est-elle une menace pour la nature au Cameroun ? »}

Ce sujet d'actualité brûlante mobilise le chapitre \emph{Technique et Nature} et exige un ancrage local rigoureux. Nous le traitons en classe, minute par minute, comme le jour de l'épreuve.

\begin{methode}[Chronologie type de l'épreuve (4h)]
\begin{enumerate}
\item \textbf{0h00–0h45 : Analyse \& Problématisation.} Lecture active, définition des termes, détection du piège, formulation de la problématique, thèse/antithèse/synthèse.
\item \textbf{0h45–1h15 : Construction du plan détaillé.} Choix du plan (dialectique, analytique, thématique), rédaction des titres de parties/sous-parties, sélection des arguments, exemples, auteurs.
\item \textbf{1h15–1h30 : Rédaction de l'introduction.} Accroche, définition, analyse du sujet, problématique, annonce du plan.
\item \textbf{1h30–3h15 : Rédaction du développement.} Environ 35 min par grande partie (I, II, III). Respecter la structure : Thèse $\rightarrow$ Argument $\rightarrow$ Exemple/Auteur $\rightarrow$ Transition.
\item \textbf{3h15–3h30 : Rédaction de la conclusion.} Bilan, réponse à la problématique, ouverture pertinente.
\item \textbf{3h30–4h00 : Relecture active (3 passes).} Voir boîte \textbf{Méthode} ci-dessous.
\end{enumerate}
\end{methode}

\begin{definition}[Analyse rigoureuse du sujet]
\begin{itemize}
\item \textbf{Technique :} Ensemble des procédés rationnels pour transformer la nature selon nos fins (Heidegger : « mise en réserve » ; Ellul : « système technique »).
\item \textbf{Nature :} Ce qui existe par soi, non fait par l'homme ; au Cameroun : forêts du bassin du Congo, mangroves, biodiversité du Parc de la Lobéké, ressources minières.
\item \textbf{Menace :} Danger, risque de destruction irréversible, mais aussi \emph{pression} qui peut susciter une régulation.
\item \textbf{Au Cameroun :} Ancrage obligatoire. Pas de dissertation hors-sol.
\item \textbf{Problématique :} « Dans quelle mesure l'exploitation technique des ressources naturelles camerounaises, si elle dégrade les écosystèmes, peut-elle aussi devenir le levier de leur préservation durable ? »
\end{itemize}
\end{definition}

\begin{propriete}[Plan dialectique retenu (avec synthèse effective)]
\textbf{I. La technique comme puissance de destruction de la nature camerounaise.}
\begin{itemize}
\item A. L'extractivisme minier et forestier : déforestation massive (ex. exploitation bois d'œuvre Est/Cameroun), pollution au mercure (orpaillage Bétaré-Oya).
\item B. L'agro-industrie et les monocultures : palmiers à huile (Socapalm, Safacam) $\rightarrow$ perte biodiversité, sols lessivés.
\item C. Les mégaprojets infrastructuraux : barrage de Lom Pangar, port de Kribi $\rightarrow$ déplacement populations, fragmentation habitats.
\end{itemize}
\textbf{II. La technique comme instrument de connaissance et de protection.}
\begin{itemize}
\item A. Surveillance par satellite/drones : suivi déforestation temps réel (Global Forest Watch), anti-braconnage (Parc de Waza).
\item B. Techniques agroécologiques : agroforesterie cacao-sous-bois, zaï amélioré Nord, biochar.
\item C. Énergies renouvelables : solaire rural (programme PERG), micro-hydroélectricité — réduction pression bois-énergie.
\end{itemize}
\textbf{III. Vers une \emph{technique écologique} encadrée par le droit et l'éthique (Synthèse).}
\begin{itemize}
\item A. Cadre juridique : Loi 96/12 (cadre environnement), Évaluations d'Impact Environnemental (EIE) obligatoires mais souvent contournées.
\item B. Gouvernance locale : Forêts communautaires (ex. Kilum-Ijim), comités de gestion participative.
\item C. Philosophie : Heidegger (\"gelassenheit\" — laisser-être), Jonas (\"principe responsabilité\"), pensée africaine (ubuntu, totémisme) : la technique doit \emph{servir} la vie, non la \emph{maîtriser}.
\end{itemize}
\end{propriete}

\begin{exemplebox}[Introduction type (chronométrée 15 min)]
\textbf{Accroche :} « La forêt du bassin du Congo, deuxième poumon vert de la planète après l'Amazonie, perd chaque année près de 400 000 hectares au Cameroun selon le Global Forest Watch (2023). »
\textbf{Définitions :} Technique (moyens), Nature (milieu), Menace (risque existentiel).
\textbf{Analyse :} Le sujet suppose une opposition ; il faut la nuancer : la technique n'est pas \emph{une} menace monolithique, mais un ensemble de pratiques aux effets ambivalents.
\textbf{Problématique :} (voir ci-dessus).
\textbf{Annonce du plan :} « Nous montrerons d'abord que le modèle technique dominant menace gravement les écosystèmes camerounais (I), avant d'examiner comment des techniques alternatives permettent de protéger la nature (II), pour enfin dégager les conditions politiques et éthiques d'une technique au service du vivant (III). »
\end{exemplebox}

\begin{exoresolu}[Développement partiel — Partie I.A (Extrait)]
\textbf{Énoncé :} Rédiger l'argumentaire de la sous-partie I.A (L'extractivisme forestier) en mobilisant Ellul et un exemple camerounais précis.
\tcblower
\textbf{Thèse :} L'exploitation forestière industrielle, par sa mécanisation lourde (bulldozers, grumiers), fragilise les sols tropicaux minces.
\textbf{Argument :} L'ouverture de pistes d'exploitation (réseau secondaire) facilite l'accès aux braconniers et aux cultivateurs itinérants (effet « route »), amplifiant la fragmentation.
\textbf{Exemple Cameroun :} Dans l'UFA 09-005 (Est), l'exploitation FSC-certifiée a réduit l'impact direct, mais les pistes secondaires non refermées ont entraîné une hausse de 37\% du braconnage (Rapport WWF 2021).
\textbf{Auteur :} Jacques Ellul, \emph{Le Système technicien} : « La technique ne tolère aucun jugement du dehors. » $\rightarrow$ La logique de rentabilité court-circuite la préservation.
\textbf{Transition :} « Mais si la technique industrielle détruit, d'autres techniques — de surveillance, de gestion douce — peuvent inverser la tendance. »
\end{exoresolu}

\begin{attention}[Piège « Sujet à double sens »]
\textbf{Sujet :} « La technique est-elle une menace pour la nature au Cameroun ? » \\
\textbf{Erreur fréquente :} Faire l'histoire de la technique au Cameroun (chemin de fer, barrages) ou lister des pollutions sans problématiser. \\
\textbf{Correctif :} Le mot « menace » appelle une évaluation \emph{critique} et \emph{mesurée} : \emph{quelles} techniques, \emph{où}, \emph{pour qui}, \emph{avec quelles alternatives} ? La problématique doit porter sur l'ambivalence, pas sur l'inventaire.
\end{attention}

\courssub{Sujet 2 — Travail en îlots : « Faut-il obéir aux lois injustes ? » avec dossier documentaire}

Ce sujet (chapitre \emph{Justice et Droit}) se prête au travail collaboratif. Chaque îlot (4–5 élèves) reçoit un dossier et produit une fiche de plan + intro en 45 min, puis présente au tableau.

\begin{experience}[Dossier documentaire (extraits)]
\textbf{Doc 1 — Constitution du Cameroun (1996, Préambule) :} « Tout peuple a droit au développement. [...] La liberté de communication, d'expression, de presse, de réunion, d'association et de syndicat est garantie. »
\textbf{Doc 2 — Nelson Mandela, \emph{Lettre de prison} (1976) :} « Il y a deux types de lois : les lois justes et les lois injustes. [...] On a non seulement une responsabilité légale, mais une responsabilité morale de désobéir aux lois injustes. »
\textbf{Doc 3 — Article \emph{Le Jour} (mars 2008 / octobre 2016) :} Manifestations contre la vie chère / crise anglophone. Répression, arrestations, état d'urgence. Témoignages : « Nous manifestions pour nos droits constitutionnels ; la loi nous criminalisait. »
\textbf{Doc 4 — Extrait Martin Luther King, \emph{Lettre de Birmingham} (1963) :} « Une loi injuste n'est pas une loi au sens plein (lex injusta non est lex). La désobéissance civile non-violente exprime le respect suprême pour la loi. »
\end{experience}

\begin{methode}[Consignes îlots (45 min)]
\begin{enumerate}
\item Lecture croisée des 4 docs (10 min) : repérer arguments pro-obéissance / pro-désobéissance.
\item Formuler la problématique commune (5 min).
\item Construire un plan dialectique à 3 parties avec arguments/docs/auteurs par sous-partie (20 min).
\item Rédiger l'introduction complète sur papier (10 min).
\item Désigner un rapporteur pour la restitution (2 min).
\end{enumerate}
\end{methode}

\begin{propriete}[Plan type attendu (synthèse des îlots)]
\textbf{I. L'obéissance à la loi comme fondement de l'ordre juste (Thèse).}
\begin{itemize}
\item A. Légitimité démocratique : loi = volonté générale (Rousseau, \emph{Contrat social}).
\item B. Sécurité juridique : prévisibilité, protection des faibles (Kelsen, pyramide des normes).
\item C. Exemple Cameroun : respect code de la route, paiement impôts $\rightarrow$ services publics.
\end{itemize}
\textbf{II. La désobéissance comme devoir moral face à l'injustice (Antithèse).}
\begin{itemize}
\item A. Distinction légalité/légitimité : lois ségrégationnistes, lois d'exception (Antigone, Sophocle ; Thoreau, \emph{Désobéissance civile}).
\item B. Désobéissance civile : publique, non-violente, assumée (King, Mandela, Gandhi).
\item C. Exemples Cameroun : grèves 2008 (vie chère), « villes mortes » 2016-17 (revendications identitaires) — illégales mais légitimes aux yeux de nombreux citoyens.
\end{itemize}
\textbf{III. L'État de droit comme cadre de résolution du conflit (Synthèse).}
\begin{itemize}
\item A. Recours juridictionnels : Conseil constitutionnel, Cour suprême, Cour africaine DH.
\item B. Rôle de la société civile, médias, syndicats : contre-pouvoirs démocratiques.
\item C. Condition : la désobéissance doit viser \emph{l'amélioration} de la loi, non sa destruction (Habermas, \emph{Droit et démocratie}).
\end{itemize}
\end{propriete}

\courssub{Exercice « Correction par les pairs » : Grille officielle MINESEC}

Chaque élève échange sa copie brouillon (Sujet 1 ou 2) avec un binôme. Notation sur 20 selon la grille simplifiée ci-dessous. L'enseignant circule, valide/ajuste les notes, fait expliciter les justifications.

\begin{popfigure}
\centering
\begin{tikzpicture}[
  cell/.style={draw=popline, minimum height=1cm, align=center, font=\small},
  header/.style={cell, fill=popblue, text=white, font=\small\bfseries},
  crit/.style={cell, fill=popblueL, font=\small\bfseries, text width=3.2cm},
  ind/.style={cell, text width=5.5cm},
  pts/.style={cell, text width=1.5cm, font=\small\bfseries},
  auto/.style={cell, text width=2.5cm}
]
% En-têtes
\node[header, minimum width=3.2cm] (c1) at (0,0) {Critère};
\node[header, minimum width=5.5cm] (c2) at (3.2,0) {Indicateurs clés (ce que le correcteur cherche)};
\node[header, minimum width=1.5cm] (c3) at (8.7,0) {Pts};
\node[header, minimum width=2.5cm] (c4) at (10.2,0) {Auto-éval. /20};

% Ligne 1 : Analyse du sujet (4 pts)
\node[crit] (a1) at (0,-1) {1. Analyse du sujet (4)};
\node[ind] (a2) at (3.2,-1) {Termes définis, sens du sujet saisi, piège évité, problématique pertinente, thèse/antithèse claires};
\node[pts] (a3) at (8.7,-1) {4};
\node[auto] (a4) at (10.2,-1) {\phantom{/20}};

% Ligne 2 : Plan (4 pts)
\node[crit] (b1) at (0,-2) {2. Plan \& Structure (4)};
\node[ind] (b2) at (3.2,-2) {Plan visible (I, II, III), logique dialectique/analytique, transitions explicites, équilibre parties, annonce respectée};
\node[pts] (b3) at (8.7,-2) {4};
\node[auto] (b4) at (10.2,-2) {\phantom{/20}};

% Ligne 3 : Contenu (8 pts)
\node[crit] (c1) at (0,-3) {3. Contenu \& Argumentation (8)};
\node[ind] (c2) at (3.2,-3) {Arguments philosophiques valides, exemples concrets (Cameroun + classiques), auteurs cités à propos, nuances, synthèse effective en III};
\node[pts] (c3) at (8.7,-3) {8};
\node[auto] (c4) at (10.2,-3) {\phantom{/20}};

% Ligne 4 : Expression (4 pts)
\node[crit] (d1) at (0,-4) {4. Expression \& Langage (4)};
\node[ind] (d2) at (3.2,-4) {Français correct (syntaxe, orthographe, vocabulaire), style philosophique (connecteurs, impersonnel), lisibilité, propreté};
\node[pts] (d3) at (8.7,-4) {4};
\node[auto] (d4) at (10.2,-4) {\phantom{/20}};

% Total
\node[header, minimum width=3.2cm] (t1) at (0,-5) {TOTAL};
\node[ind] (t2) at (3.2,-5) {};
\node[pts] (t3) at (8.7,-5) {20};
\node[auto] (t4) at (10.2,-5) {\phantom{/20}};
\end{tikzpicture}
\legende{Grille d'évaluation officielle MINESEC simplifiée — Dissertation Philosophie Série Technique (coeff. 2 au Probatoire). Chaque critère est noté sur le nombre de points indiqué ; la somme donne la note sur 20.}
\end{popfigure}

\begin{methode}[La relecture active en 3 passes]
\begin{enumerate}
\item \textbf{Pass 1 — Structure (5 min) :} Le plan est-il visible ? (I, II, III dans la marge). Les transitions existent-elles ? Intro/Concl complètes ? Annonce du plan respectée ?
\item \textbf{Pass 2 — Fond (10 min) :} Chaque argument tient-il ? Exemples pertinents et \emph{datés} ? Auteurs justes (pas de name-dropping) ? Synthèse en III (pas simple résumé) ? Ancrage Cameroun effectif ?
\item \textbf{Pass 3 — Forme (5 min) :} Orthographe (accords, homophones), syntaxe (phrases complètes), lisibilité (ratures propres, paragraphes sautés), connecteurs logiques (\emph{donc, toutefois, en effet, par conséquent}).
\end{enumerate}
\end{methode}

\begin{propriete}[La note de dissertation = f(Structure, Pertinence, Profondeur, Expression)]
Un défaut majeur sur l'un effondre la note globale, même si le style est parfait.
\begin{itemize}
\item \textbf{Hors-sujet (HS) :} Note maximale 4/20 (souvent 0–2), quel que soit le style.
\item \textbf{Paraphrase / Pas de problématique :} Plafond 8/20.
\item \textbf{Plan binaire (Pour/Contre sans synthèse) :} Plafond 10/20.
\item \textbf{Absence d'exemples concrets :} -2 à -3 pts sur Contenu.
\item \textbf{Expression défaillante (illisible, fautes graves) :} -1 à -2 pts sur Expression, mais peut faire basculer en dessous de 10 si le fond est limite.
\end{itemize}
\end{propriete}

\courssub{Analyse de copies types anonymisées : Notation collective}

Trois copies (A, B, C) sur le Sujet 1 sont projetées (ou photocopiées). La classe note collectivement, justifie, identifie points forts/faibles. L'enseignant révèle la note officielle et le commentaire de jury.

\begin{exemplebox}[Copie A — Note 14/20 (Bien) — Points forts / Faibles]
\textbf{Forces :} Problématique excellente (ambivalence technique), plan dialectique clair, 3 exemples Cameroun datés (Lom Pangar, Socapalm, drones Waza), auteurs (Ellul, Jonas, Heidegger) utilisés à bon escient, synthèse III juridico-éthique, style soigné.
\textbf{Faibles :} Transition I$\rightarrow$II un peu abrupte, ouverture conclusion un peu générique (« l'avenir nous le dira »), 2 fautes d'accord.
\textbf{Commentaire jury :} « Copie mature, ancrage réel, pensée structurée. L'ouverture gagnerait à être plus précise (ex. : réforme du code forestier 2023). »
\end{exemplebox}

\begin{exemplebox}[Copie B — Note 7/20 (Insuffisant) — Points forts / Faibles]
\textbf{Forces :} Intro correcte, définitions présentes, quelques exemples (bois, pétrole).
\textbf{Faibles :} Plan binaire « Pour : la technique détruit / Contre : la technique protège » sans III de synthèse. Pas d'auteurs. Exemples vagues (« les forêts », « les usines »). Pas de transition. Conclusion = résumé. Style oral (« du coup », « en fait »).
\textbf{Commentaire jury :} « Plan binaire = plafond 10. L'absence de synthèse et d'auteurs fait chuter à 7. Travail sur la structure dialectique urgent. »
\end{exemplebox}

\begin{exemplebox}[Copie C — Note 11/20 (Passable) — Points forts / Faibles]
\textbf{Forces :} Plan en 3 parties visible, exemples Cameroun (Kribi, Nord), conclusion avec ouverture sur « transition écologique ».
\textbf{Faibles :} Problématique faible (« La technique est-elle bonne ou mauvaise ? »). Partie II = liste d'exemples sans argument philosophique. Auteurs cités mais non explicités (name-dropping : « Kant dit que... » sans lien). Quelques ratures illisibles.
\textbf{Commentaire jury :} « Le fond est là mais la profondeur philosophique manque. Le name-dropping ne vaut pas argumentation. Lisibilité à soigner. »
\end{exemplebox}

\begin{popfigure}
\centering
\begin{tikzpicture}[
  scale=0.85,
  every node/.style={font=\small},
  annot/.style={draw=poporange, fill=poporange!15, rounded corners=2pt, inner sep=2pt},
  symbol/.style={draw=popblue, fill=popblue!20, circle, minimum size=6pt, inner sep=0pt},
  line/.style={draw=popline, thick, -{Stealth[scale=0.8]}}
]
% Simulated copy page
\draw[popline, thick] (0,0) rectangle (14,10);
% Title
\node[align=center, font=\bfseries] at (7,9.5) {SUJET : La technique est-elle une menace pour la nature au Cameroun ?};
% Intro
\node[annot, text width=12cm] (intro) at (7,8.7) {\textbf{Intro :} Accroche (chiffre déforestation). Définitions. Problématique : « Dans quelle mesure... » Plan annoncé. \textcolor{popgreen}{$\checkmark$ Bien}};
% Symboles marge gauche
\node[symbol] at (0.5,8.7) {+};
\node[symbol] at (0.5,7.8) {?};
% Partie I
\node[annot, text width=12cm] (I) at (7,7.5) {\textbf{I. La technique détruit} A. Forestier (Ellul) — Ex: Est Cameroun. B. Minier — Ex: Bétaré-Oya. C. Agro-industrie — Ex: Socapalm. \textcolor{popgreen}{$\checkmark$ Arguments + Exemples}};
\node[symbol] at (0.5,7.5) {+};
% Transition
\node[annot, text width=12cm] (trans) at (7,6.3) {\textbf{Transition :} « Cependant, la technique n'est pas que destruction. » \textcolor{poporange}{$\sim$ Faible}};
\node[symbol] at (0.5,6.3) {T};
% Partie II
\node[annot, text width=12cm] (II) at (7,5.8) {\textbf{II. La technique protège} A. Drones/satellites. B. Agroécologie. C. Énergies renouvelables. \textcolor{poporange}{$\sim$ Liste sans thèse directrice}};
\node[symbol] at (0.5,5.8) {-};
% Partie III
\node[annot, text width=12cm] (III) at (7,4.5) {\textbf{III. Vers une technique écologique} A. Loi 96/12. B. Forêts communautaires. C. Jonas/Heidegger. \textcolor{popgreen}{$\checkmark$ Synthèse}};
\node[symbol] at (0.5,4.5) {+};
% Conclusion
\node[annot, text width=12cm] (concl) at (7,3.2) {\textbf{Concl :} Bilan. Réponse. Ouverture : « Il faut une éthique de la responsabilité. » \textcolor{popgreen}{$\checkmark$ Correcte}};
\node[symbol] at (0.5,3.2) {+};
% Symboles légende
\begin{scope}[yshift=-1.5cm]
\node[symbol] (s1) at (1,0) {};
\node[right=0.2cm of s1] {+ : Argument solide / Exemple pertinent};
\node[symbol] (s2) at (1,-0.8) {};
\node[right=0.2cm of s2] {- : Argument faible / Exemple manquant};
\node[symbol] (s3) at (1,-1.6) {};
\node[right=0.2cm of s3] {? : Incompréhensible / Hors-sujet local};
\node[symbol] (s4) at (1,-2.4) {};
\node[right=0.2cm of s4] {T : Transition};
\node[symbol] (s5) at (1,-3.2) {};
\node[right=0.2cm of s5] {A : Auteur bien utilisé};
\node[symbol] (s6) at (1,-4.0) {};
\node[right=0.2cm of s6] {HS : Hors-sujet (rouge)};
\end{scope}
\end{tikzpicture}
\legende{Simulation d'une copie annotée par un correcteur Bac. Symboles standards : + (point fort), - (point faible), ? (obscur/HS), T (transition), A (auteur), HS (hors-sujet). La marge gauche concentre les symboles pour une lecture rapide.}
\end{popfigure}

\courssub{Révision des 10 pièges du Bac — Mémo anti-erreur}

\begin{center}
\begin{tabularx}{\linewidth}{|c|X|X|}\hline
\textbf{N\textsuperscript{o}} & \textbf{Piège} & \textbf{Parade} \\ \hline
1 & \textbf{Hors-sujet (HS)} & \textbf{Analyse 45 min :} Souligner mots-clés, définir, tester : « Ma thèse répond-elle \emph{exactement} à la question ? » \\ \hline
2 & \textbf{Paraphrase du sujet} & \textbf{Problématiser :} Transformer l'affirmation en question philosophique (Pourquoi ? Dans quelle mesure ?) \\ \hline
3 & \textbf{Plan binaire (Pour/Contre)} & \textbf{Exiger 3 parties :} Thèse — Antithèse — \textbf{Synthèse} (dépassement, conditions, nuances) \\ \hline
4 & \textbf{Absence d'exemples concrets} & \textbf{Banque 10 exemples Cameroun} (voir Fiche Mémo) + 5 classiques. 1 exemple par sous-partie minimum. \\ \hline
5 & \textbf{Style oral / familier} & \textbf{Relecture Pass 3 :} Remplacer « du coup / en fait / genre » par « par conséquent / en réalité / par exemple » \\ \hline
6 & \textbf{Gestion du temps} & \textbf{Chrono rigide :} 45' analyse/plan, 15' intro, 3$\times$35' parties, 15' concl, 20' relecture. \textbf{Ne pas dépasser.} \\ \hline
7 & \textbf{Écriture illisible / ratures sales} & \textbf{Brouillon propre :} Écrire tous les 2 lignes, ratures à la règle, paragraphes sautés. Lisibilité = respect du correcteur. \\ \hline
8 & \textbf{Oubli de la conclusion} & \textbf{Alarme 3h45 :} Arrêter le développement, écrire conclusion (bilan + réponse + ouverture). Mieux vaut conclure imparfait que pas du tout. \\ \hline
9 & \textbf{Ouverture « gadget »} & \textbf{Ouverture = prolongement :} Autre domaine (art, science, politique), autre époque, perspective concrète (projet de loi, ONG). Pas de citation décorative. \\ \hline
10 & \textbf{Contradiction Intro/Développement} & \textbf{Vérification Pass 1 :} Les titres I, II, III correspondent-ils \emph{exactement} à l'annonce du plan ? \\ \hline
\end{tabularx}
\end{center}

\begin{attention}[Piège n°3 — Plan binaire : le tueur silencieux]
Un plan « I. Oui, la technique menace / II. Non, elle ne menace pas » plafonne à \textbf{10/20} (souvent 8). Pourquoi ? Parce qu'il ne \emph{pense} pas la tension, il la \emph{constate}. La philosophie exige une \textbf{synthèse} (III) qui articule les conditions : « La technique menace la nature \emph{lorsqu'elle est soumise à la seule rentabilité} ; elle la protège \emph{lorsqu'elle est encadrée par le droit et orientée par une éthique de la responsabilité}. » C'est là que se joue la note 12 $\rightarrow$ 16.
\end{attention}

\courssub{Remise de la « Fiche Mémo Bac Philo Tech » — 1 page recto-verso}

Distribuée en fin de séance, plastifiée si possible. Chaque élève la garde jusqu'à l'épreuve.

\begin{exemplebox}[FICHE MÉMO BAC PHILO TECH — RECTO : CHECK-LIST 3 TEMPS]
\textbf{AVANT (J-30 à J-1) :}
\begin{itemize}
\item[$\square$] 10 sujets types traités en conditions chrono (4h)
\item[$\square$] 20 concepts clés maîtrisés (définition + auteur + exemple)
\item[$\square$] 10 exemples Cameroun « passe-partout » mémorisés (voir verso)
\item[$\square$] Grille MINESEC connue par cœur
\item[$\square$] Matériel prêt : 2 stylos bleu/noir, règle, montre, carte d'identité, convocation
\end{itemize}
\textbf{PENDANT (4h) :}
\begin{itemize}
\item[$\square$] 0h00–0h45 : Analyse $\rightarrow$ Problématique $\rightarrow$ Plan détaillé (brouillon)
\item[$\square$] 0h45–1h00 : Intro complète (brouillon $\rightarrow$ propre)
\item[$\square$] 1h00–3h15 : Développement (I, II, III) — \textbf{1 partie = 35 min max}
\item[$\square$] 3h15–3h30 : Conclusion
\item[$\square$] 3h30–4h00 : \textbf{Relecture 3 passes} (Structure / Fond / Forme)
\item[$\square$] \textbf{Alarme 3h45 :} STOP développement $\rightarrow$ Conclusion obligatoire
\end{itemize}
\textbf{APRÈS (Sortie de salle) :}
\begin{itemize}
\item[$\square$] Ne pas discuter le sujet (stress contagieux)
\item[$\square$] Repos, hydratation, prochaine épreuve
\end{itemize}
\end{exemplebox}

\begin{exemplebox}[FICHE MÉMO BAC PHILO TECH — VERSO : 20 CONCEPTS CLÉS + 10 EXEMPLES CAMEROUN « PASSE-PARTOUT »]
\textbf{20 Concepts clés (Définition 1 ligne + Auteur repère) :}
\begin{enumerate}
\item \textbf{Technique} : Moyen rationnel de transformation de la nature (Aristote, Heidegger).
\item \textbf{Nature} : Ce qui est par soi, non fait par l'homme (Aristote, Rousseau).
\item \textbf{Travail} : Médiation homme-nature par la technique (Marx, Hegel).
\item \textbf{Progrès} : Amélioration mesurable vs progrès moral (Condorcet, Jonas).
\item \textbf{Responsabilité} : Répondre de ses actes vis-à-vis du futur (Jonas, Ricoeur).
\item \textbf{Justice} : Donner à chacun son dû (Aristote, Rawls, Senghor).
\item \textbf{Droit} : Ensemble de normes coercitives (Kelsen, Hart).
\item \textbf{Légitimité} : Acceptabilité morale du pouvoir (Weber, Locke).
\item \textbf{État de droit} : Soumission du pouvoir à la loi (Kelsen, Constitution 1996).
\item \textbf{Désobéissance civile} : Transgression publique, non-violente, assumée (Thoreau, King, Mandela).
\item \textbf{Liberté} : Autonomie / absence de contrainte (Kant, Sartre, Bergson).
\item \textbf{Conscience} : Connaissance réflexive de soi (Descartes, Husserl).
\item \textbf{Sujet} : Instance qui pense, veut, agit (Descartes, Kant, Ricoeur).
\item \textbf{Vérité} : Adéquation intellect/réalité + cohérence + pragmatisme (Aristote, Nietzsche).
\item \textbf{Science} : Savoir rationnel, démontrable, cumulatif (Popper, Kuhn, Bachelard).
\item \textbf{Technoscience} : Fusion science/technique (Heidegger, Ellul, Hottois).
\item \textbf{Aliénation} : Perte de soi dans l'objet/autre (Marx, Hegel, Fanon).
\item \textbf{Idéologie} : Représentation faussée servant des intérêts (Marx, Althusser, Mannheim).
\item \textbf{Histoire} : Devenir humain temporalisé (Hegel, Marx, Benjamin).
\item \textbf{Condition humaine} : Finitude, mortalité, natalité, pluralité (Arendt, Heidegger, Janssens).
\end{enumerate}

\textbf{10 Exemples Cameroun « Passe-partout » (à adapter selon sujet) :}
\begin{enumerate}
\item \textbf{Barrage de Lom Pangar} (Énergie / Déplacement / Environnement) — Technique, Nature, Justice, Progrès.
\item \textbf{Port autonome de Kribi / Zone industrielle} (Infrastructure / Mangrove / Emploi) — Technique, Nature, Travail, Justice.
\item \textbf{Exploitation forestière Est (UFA, certification FSC)} — Technique, Nature, Responsabilité, Droit.
\item \textbf{Orpaillage artisanal Bétaré-Oya / Kette} (Mercure / Revenus / Santé) — Technique, Nature, Justice, Aliénation.
\item \textbf{Agro-industrie palmier/hévéa (Socapalm, Safacam, HEVECAM)} — Technique, Nature, Travail, Progrès.
\item \textbf{Crise anglophone 2016–... : « Villes mortes », manifestations, répression} — Droit, Justice, Légitimité, Désobéissance, Liberté.
\item \textbf{Manifestations vie chère 2008 (Émeutes de la faim)} — Justice, Droit, Légitimité, Histoire.
\item \textbf{Forêts communautaires Kilum-Ijim / Ngoyla-Mintom} (Gouvernance locale / Biodiversité) — Nature, Justice, Droit, Responsabilité.
\item \textbf{Programme PERG (Électrification rurale solaire)} — Technique, Progrès, Justice, Nature.
\item \textbf{Code du travail 1992 / Conventions CCT (Syndicats / Patronat)} — Travail, Justice, Droit, Liberté.
\end{enumerate}
\end{exemplebox}

\begin{aretenir}[Simulation chronométrée : Objectif « Plancher 10/20 »]
\textbf{Contrat minimal pour ne pas échouer :}
\begin{itemize}
\item \textbf{Zéro hors-sujet} (analyse 45 min respectée).
\item \textbf{Plan visible} (I, II, III dans la marge, transitions écrites).
\item \textbf{3 parties rédigées} (même inégales, même imparfaites).
\item \textbf{Intro + Concl complètes}.
\end{itemize}
\textbf{Résultat :} Note \textbf{10–11/20} quasi garantie (structure + pertinence minimale).
\textbf{Le delta 10 $\rightarrow$ 16 se joue sur :}
\begin{itemize}
\item Finesse des exemples (datés, localisés, variés).
\item Pertinence des auteurs (cités \emph{pour} argumenter, pas pour décorer).
\item Qualité de la synthèse (III = vraie articulation, pas résumé).
\item Style philosophique (connecteurs, impersonnel, vocabulaire précis).
\item Lisibilité / propreté (respect du correcteur).
\end{itemize}
\end{aretenir}

\begin{savaistu}[Le jury du Bac Technik valorise l'ancrage local]
Les correcteurs du Baccalauréat Technique (MINESEC) le répètent chaque année en délibération : \textbf{un exemple camerounais précis, daté, contextualisé, vaut mieux que trois références européennes vagues}. Cela prouve deux choses :
\begin{enumerate}
\item \textbf{Maturité citoyenne :} L'élève philosophe \emph{ici et maintenant}, pas dans une tour d'ivoire.
\item \textbf{Capacité de transfert :} Il sait mobiliser ses connaissances (géographie, histoire, actualité) pour éclairer un concept philosophique.
\end{enumerate}
\textbf{Exemple :} Dans une copie sur « La technique libère-t-elle l'homme ? », l'élève qui cite « L'introduction du téléphone mobile (MTN/Orange) dans les zones rurales du Nord a permis aux éleveurs de connaître les cours des marchés à Douala en temps réel, réduisant leur dépendance aux intermédiaires » marque 2 points de plus que celui qui écrit « Le téléphone portable facilite la communication. » \textbf{Ancrage = Preuve de pensée en acte.}
\end{savaistu}

\begin{exercice}[Entraînement individuel — Devoir maison (1 semaine)]
Traitez \textbf{l'un} des deux sujets ci-dessous en conditions réelles (4h chrono, brouillon + copie propre, relecture 3 passes). Remettez la copie + la grille d'auto-évaluation remplie.
\begin{enumerate}
\item « Le travail technique dégrade-t-il l'homme ? » (Chapitre Technique / Travail)
\item « La science peut-elle fonder la morale ? » (Chapitre Science / Vérité / Morale)
\end{enumerate}
\textbf{Conseil :} Utilisez la Fiche Mémo. Visez le « Plancher 10 » d'abord, puis la finesse.
\end{exercice}

\begin{corrige}[Exercice — Grille d'auto-évaluation type (à joindre à la copie)]
\begin{center}
\begin{tabularx}{\linewidth}{|c|c|X|}\hline
\textbf{Critère} & \textbf{Note / Max} & \textbf{Justification (1 phrase)} \\ \hline
Analyse (4) & /4 & \\ \hline
Plan (4) & /4 & \\ \hline
Contenu (8) & /8 & \\ \hline
Expression (4) & /4 & \\ \hline
\textbf{TOTAL} & \textbf{/20} & \\ \hline
\end{tabularx}
\end{center}
\textbf{Engagement :} « Je certifie avoir respecté le chrono 4h et effectué les 3 passes de relecture. » \hfill Signature :
\end{corrige}

\coursec{Transfert : La dissertation hors examen (Vie professionnelle \& Citoyenne)}

Après l'atelier de consolidation sur les sujets types du Baccalauréat technique, une question légitime se pose : à quoi sert tout cela une fois l'examen passé ? La réponse est simple et décisive : la dissertation philosophique n'est pas un exercice scolaire artificiel, c'est un \cle{entraînement de la pensée} dont la structure se retrouve partout dans la vie professionnelle et citoyenne. Le technicien qui rédige un rapport de stage, l'employé qui défend un projet en réunion, le citoyen qui intervient au conseil municipal, le jeune diplômé qui écrit une lettre de motivation : tous mobilisent, sans toujours le savoir, les quatre compétences fondamentales de la dissertation : \cle{analyser}, \cle{structurer}, \cle{argumenter}, \cle{synthétiser}. Cette dernière section montre comment transférer ces acquis hors de la salle d'examen.

\courssub{Les quatre compétences transférables}

Commençons par une intuition. Quand un chef de chantier à Douala doit expliquer à son client pourquoi les travaux accusent un retard, il ne peut pas se contenter de dire « c'est compliqué ». Il doit d'abord \emph{analyser} la situation (quelle est la vraie cause du retard ?), puis \emph{structurer} son explication (dans quel ordre présenter les faits ?), ensuite \emph{argumenter} (quelles preuves avancer : bons de livraison, météo, planning ?), enfin \emph{synthétiser} (que propose-t-on concrètement ?). C'est exactement le parcours d'une dissertation : analyse du sujet, plan, argumentation, conclusion.

\begin{aretenir}[Les quatre compétences mères]
La dissertation philosophique développe quatre compétences directement transférables :
\begin{enumerate}
\item \cle{Analyser} : décomposer une situation complexe, identifier le vrai problème (équivalent de l'analyse du sujet et de la problématique).
\item \cle{Structurer} : organiser les idées selon une progression logique (équivalent du plan).
\item \cle{Argumenter} : justifier ses affirmations par des preuves et des exemples (équivalent du développement).
\item \cle{Synthétiser} : dégager une conclusion claire et actionnable (équivalent de la conclusion).
\end{enumerate}
\end{aretenir}

\begin{popfigure}
\begin{center}
\begin{tikzpicture}[
comp/.style={rectangle, rounded corners=3pt, draw=popblue, fill=popblueL, thick, minimum width=2.6cm, minimum height=0.85cm, align=center, font=\small\bfseries},
ctx/.style={rectangle, rounded corners=3pt, draw=poporange, fill=poporangeL, thick, minimum width=2.9cm, minimum height=0.85cm, align=center, font=\small},
arr/.style={-{Stealth[length=2.5mm]}, thick, popmuted}]
% Compétences (gauche)
\node[comp] (c1) at (0,3.6) {Analyser};
\node[comp] (c2) at (0,2.4) {Structurer};
\node[comp] (c3) at (0,1.2) {Argumenter};
\node[comp] (c4) at (0,0) {Synthétiser};
\node[font=\small\itshape, popblue] at (0,4.5) {Compétences de la dissertation};
% Contextes (droite)
\node[ctx] (x1) at (8.2,4.2) {Rapport de stage};
\node[ctx] (x2) at (8.2,3.15) {Réunion de chantier};
\node[ctx] (x3) at (8.2,2.1) {Débat municipal};
\node[ctx] (x4) at (8.2,1.05) {Lettre de motivation};
\node[ctx] (x5) at (8.2,0) {Veille technologique};
\node[font=\small\itshape, poporange] at (8.2,5.1) {Contextes professionnels et citoyens};
% Flèches
\draw[arr] (c1.east) -- (x1.west);
\draw[arr] (c1.east) -- (x5.west);
\draw[arr] (c2.east) -- (x1.west);
\draw[arr] (c2.east) -- (x4.west);
\draw[arr] (c3.east) -- (x2.west);
\draw[arr] (c3.east) -- (x3.west);
\draw[arr] (c4.east) -- (x2.west);
\draw[arr] (c4.east) -- (x4.west);
\end{tikzpicture}
\end{center}
\legende{Diagramme de transfert : les quatre compétences de la dissertation philosophique irriguent les écrits et les prises de parole de la vie technique, professionnelle et citoyenne.}
\end{popfigure}

\begin{savaistu}[De la dissertation au \emph{Design Thinking}]
La méthode de la dissertation (problématiser $\rightarrow$ analyser $\rightarrow$ synthétiser) est la base du \emph{Design Thinking} et de la \emph{résolution de problèmes complexes} enseignés dans les grandes écoles d'ingénieurs et de management du monde entier. Ce que vous pratiquez sur une copie de philosophie est donc, structurellement, la même démarche que celle des ingénieurs qui conçoivent un produit ou des consultants qui diagnostiquent une entreprise : partir d'une question bien posée, explorer les solutions possibles, puis proposer une réponse justifiée.
\end{savaistu}

\courssub{La note de synthèse et le rapport technique}

En BTS, en HND puis dans la vie active, l'écrit professionnel dominant est la \cle{note de synthèse} ou le \cle{rapport technique} : rapport de stage, rapport d'intervention, note de service, étude de faisabilité. Or sa structure est rigoureusement celle de la dissertation.

\begin{center}
\begin{tabularx}{\linewidth}{|l|X|X|}\hline
\rowcolor{popblueL} \textbf{Étape} & \textbf{Dissertation philosophique} & \textbf{Rapport technique / note de synthèse} \\ \hline
1 & Analyse du sujet, problématique & Définition précise du problème à résoudre \\ \hline
2 & Plan en deux ou trois parties & Analyse structurée : causes, constats, options \\ \hline
3 & Arguments et exemples & Données chiffrées, mesures, témoignages \\ \hline
4 & Conclusion ouverte & Recommandations concrètes et hiérarchisées \\ \hline
\end{tabularx}
\end{center}

\begin{methode}[Méthode universelle de l'écrit professionnel]
Pour tout écrit professionnel (note, rapport, mémoire) :
\begin{enumerate}
\item \textbf{Question centrale} : formuler la problématique en une phrase interrogative (« Comment réduire les pannes de la machine X ? »).
\item \textbf{Plan en trois axes} : \emph{Pourquoi ?} (diagnostic, causes), \emph{Comment ?} (solutions possibles comparées), \emph{Et après ?} (mise en œuvre, suivi).
\item \textbf{Preuves} : chiffres, mesures, témoignages, références — jamais d'affirmation sans justification.
\item \textbf{Synthèse actionnable} : une conclusion qui dit quoi faire, par qui, quand.
\end{enumerate}
\end{methode}

\begin{exemplebox}[Exemple : note de synthèse d'un technicien en maintenance]
Un technicien d'une minoterie de Bafoussam constate des arrêts fréquents d'un moteur. Sa note suit le schéma de la dissertation : \emph{problématique} (« Pourquoi le moteur M3 tombe-t-il en panne trois fois par mois ? ») ; \emph{analyse} (usure des roulements, surchauffe due à une ventilation obstruée, graissage irrégulier) ; \emph{arguments chiffrés} (températures relevées, historique des pannes) ; \emph{recommandations} (remplacement des roulements, planning de graissage hebdomadaire, coût estimé). Un tel document, clair et structuré, est immédiatement pris au sérieux par la hiérarchie : c'est la dissertation appliquée à l'atelier.
\end{exemplebox}

\begin{exemplebox}[Un chiffre qui doit vous convaincre]
Selon l'enquête INSERTION 2022 (MINESEC/MINEFOP), \textbf{82\,\% des employeurs du secteur technique formel} citent la « capacité à structurer un raisonnement écrit » comme une compétence critique \emph{manquante} chez les jeunes diplômés. Maîtriser la dissertation, c'est donc détenir un avantage décisif sur le marché du travail camerounais.
\end{exemplebox}

\courssub{La prise de parole : le « pitch » philosophique de trois minutes}

En réunion de chantier, devant un jury de stage ou lors d'un débat, on dispose souvent de trois minutes pour convaincre. C'est le format de l'\emph{elevator pitch} (discours d'ascenseur). Sa structure est celle du développement philosophique, comprimée :

\begin{enumerate}
\item \textbf{Thèse} (30 secondes) : annoncer clairement sa position. « Je propose que nous changions de fournisseur de ciment. »
\item \textbf{Arguments} (90 secondes) : deux ou trois raisons ordonnées, du plus simple au plus fort.
\item \textbf{Exemples} (45 secondes) : un fait concret, un chiffre, une expérience vécue.
\item \textbf{Conclusion} (15 secondes) : rappeler la thèse et appeler à la décision.
\end{enumerate}

\begin{attention}[Argumenter n'est pas vaincre]
Ne confondez jamais \emph{argumenter} et \emph{convaincre à tout prix}. La dissertation — et toute discussion honnête — vise la \cle{vérité}, ou du moins la meilleure raison disponible, et non la victoire rhétorique. L'\cle{honnêteté intellectuelle} prime : reconnaître la force d'un argument adverse n'est pas une faiblesse, c'est une preuve de maturité. En réunion comme en copie, celui qui triche avec les faits finit toujours par perdre sa crédibilité.
\end{attention}

\courssub{La lettre de motivation : une dissertation déguisée}

La lettre de motivation est, structurellement, une dissertation dont le « sujet » serait : « Pourquoi suis-je le candidat qu'il vous faut ? »

\begin{itemize}
\item \textbf{Accroche} (= amorce de l'introduction) : une phrase qui capte l'attention du recruteur.
\item \textbf{Problématique} (= le besoin de l'employeur) : montrer qu'on a compris le problème que l'entreprise cherche à résoudre en recrutant.
\item \textbf{Plan « vous — moi — nous »} (= le développement) : vos besoins, mes compétences et expériences, notre projet commun.
\item \textbf{Conclusion} (= appel à l'action) : demander explicitement un entretien.
\end{itemize}

Le même schéma vaut pour un \cle{plaidoyer} (défendre un projet associatif, une demande de financement) : on y défend une thèse avec des arguments et des preuves, exactement comme dans la deuxième partie d'un plan dialectique.

\courssub{L'analyse critique des médias : le détecteur de fallaces}

Le citoyen et le technicien sont bombardés d'informations : réseaux sociaux, éditoriaux, rumeurs de quartier, \emph{fake news}. L'entraînement à la dissertation fournit un véritable \cle{détecteur d'arguments fallacieux} (ou sophismes) : il apprend à repérer la thèse implicite d'un texte, à vérifier ses arguments, à dénoncer ses exemples biaisés et à constater l'absence de problématique — un article qui affirme sans jamais se poser de question est toujours suspect.

\begin{popfigure}
\begin{center}
\begin{tikzpicture}[
fal/.style={rectangle, rounded corners=3pt, draw=poppink, fill=poppinkL, thick, minimum width=3.1cm, minimum height=1.5cm, align=center, font=\scriptsize},
hub/.style={circle, draw=popdark, fill=popgoldL, thick, minimum size=2.5cm, align=center, font=\small\bfseries}]
\node[hub, align=center] (h) at (0,0){Détecteur\\de fallaces};
\node[fal, align=center] (f1) at (-4.6,2.6){\textbf{Homme de paille}\\ « Tu critiques le projet,\\ donc tu es contre le\\ développement du pays »};
\node[fal, align=center] (f2) at (0,3.4){\textbf{Ad hominem}\\ « Ne l'écoute pas,\\ il n'a même pas\\ le Bac »};
\node[fal, align=center] (f3) at (4.6,2.6){\textbf{Fausse dichotomie}\\ « Soit tu es avec nous,\\ soit tu es contre nous »};
\node[fal, align=center] (f4) at (-4.6,-2.6){\textbf{Appel à la peur}\\ « Si on change de\\ fournisseur, ce sera\\ la faillite assurée »};
\node[fal, align=center] (f5) at (4.6,-2.6){\textbf{Généralisation hâtive}\\ « Mon voisin a été mal\\ servi : tous les garages\\ de Yaoundé sont des voleurs »};
\draw[-{Stealth[length=2mm]}, thick, popmuted] (h) -- (f1);
\draw[-{Stealth[length=2mm]}, thick, popmuted] (h) -- (f2);
\draw[-{Stealth[length=2mm]}, thick, popmuted] (h) -- (f3);
\draw[-{Stealth[length=2mm]}, thick, popmuted] (h) -- (f4);
\draw[-{Stealth[length=2mm]}, thick, popmuted] (h) -- (f5);
\end{tikzpicture}
\end{center}
\legende{Le « détecteur de fallaces » : cinq arguments fallacieux courants, illustrés par des exemples de la vie quotidienne camerounaise. Les repérer, c'est appliquer hors l'école l'esprit critique formé par la dissertation.}
\end{popfigure}

\begin{definition}[Principe de charité]
Le \cle{principe de charité} (en philosophie) consiste à interpréter l'argument de l'adversaire sous sa forme la plus forte et la plus cohérente avant de le critiquer. C'est le contraire de l'homme de paille : au lieu de caricaturer la position adverse pour la réfuter facilement, on la reconstruit au mieux, puis on la discute. Ce principe est la base du dialogue démocratique et de toute négociation professionnelle honnête.
\end{definition}

\courssub{L'éthique de la discussion}

La dissertation forme enfin une \cle{éthique de la discussion}, indispensable au citoyen et au travailleur :

\begin{itemize}
\item \textbf{Respecter l'adversaire} en appliquant le principe de charité : discuter ses meilleurs arguments, pas ses plus faibles.
\item \textbf{Distinguer la personne et l'idée} : on peut réfuter une idée sans attaquer celui qui la porte (refus de l'\emph{ad hominem}).
\item \textbf{Refuser l'argument d'autorité nu} : « c'est vrai parce que le chef l'a dit » n'est pas un argument ; l'autorité doit être doublée d'une justification rationnelle.
\item \textbf{Accepter de changer d'avis} si l'argument adverse est meilleur : c'est la marque de la rationalité, non de la faiblesse. Comme le soulignait déjà Socrate, le dialogue vise à accoucher la vérité, non à humilier l'interlocuteur.
\end{itemize}

\begin{exoresolu}[Atelier de transfert : du sujet de philo au rapport de stage]
\textbf{Énoncé.} Vous êtes stagiaire dans un atelier de menuiserie métallique à Ngaoundéré. Le chef d'atelier vous demande une note d'une page sur la question : « Faut-il investir dans une nouvelle soudeuse ou réparer l'ancienne ? » Construisez le plan détaillé de cette note en appliquant la méthode de la dissertation.
\tcblower
\textbf{Solution détaillée.}
\begin{enumerate}
\item \textbf{Problématique} (analyse du sujet) : « Quel choix, entre réparation et remplacement, garantit la continuité de la production au moindre coût ? » Les termes sont définis : coût d'achat, coût de réparation, coût des arrêts de production.
\item \textbf{Premier axe — Pourquoi ?} (diagnostic) : l'ancienne soudeuse tombe en panne deux fois par mois ; chaque arrêt coûte environ \SI{50000}{FCFA} de production perdue ; les réparations cumulées atteignent \SI{300000}{FCFA} sur un an.
\item \textbf{Deuxième axe — Comment ?} (comparaison des options) : réparer (coût immédiat faible, \SI{80000}{FCFA}, mais pannes récurrentes) contre remplacer (coût élevé, \SI{900000}{FCFA}, mais garantie de deux ans et productivité accrue de 20\,\%).
\item \textbf{Troisième axe — Et après ?} (recommandation actionnable) : proposer l'achat échelonné ou la location-vente, avec un calendrier et un responsable désigné.
\item \textbf{Conclusion} : la thèse défendue — « le remplacement est rentable en dix-huit mois » — est justifiée par les chiffres, exactement comme une thèse de dissertation est justifiée par des arguments et des exemples.
\end{enumerate}
\end{exoresolu}

\begin{exercice}[Débat municipal]
Le conseil municipal de votre ville débat de la question : « Faut-il interdire les motos-taxis dans le centre-ville ? » Rédigez en dix lignes un « pitch » de trois minutes (thèse, deux arguments, un exemple, conclusion) en respectant le principe de charité : commencez par présenter honnêtement le meilleur argument du camp opposé avant de défendre votre position.
\end{exercice}

\begin{corrige}[Exercice : Débat municipal]
Éléments de correction attendus : (1) une ouverture charitable, par exemple « Ceux qui défendent les motos-taxis ont raison de rappeler qu'ils assurent le transport de milliers d'habitants et font vivre des milliers de jeunes » ; (2) une thèse claire, par exemple « il faut non pas interdire mais réglementer » ; (3) deux arguments ordonnés (sécurité routière, ordre urbain, mais aussi emploi des jeunes) ; (4) un exemple concret (une ville camerounaise ayant instauré des couloirs réservés ou des cartes professionnelles de moto-taximen) ; (5) une conclusion qui appelle à une décision précise. Le correcteur vérifiera surtout que l'élève distingue la personne de l'idée et n'utilise aucune fallace (ni ad hominem, ni fausse dichotomie « interdire ou laisser faire »).
\end{corrige}

\courssub{Projet de filière : les passerelles avec les modules techniques}

Ce transfert n'est pas un supplément d'âme : il est inscrit dans votre cursus. En classe de Première technique, le module de \cle{Communication professionnelle} vous demande de rédiger des comptes rendus, des courriers et de présenter oralement un dossier : c'est l'application directe des quatre compétences (analyser, structurer, argumenter, synthétiser). De même, le module de \cle{Droit du travail et législation sociale} exige d'analyser des situations juridiques concrètes (licenciement, contrat, conflit collectif) : il faut y dégager le problème de droit (la problématique), examiner les arguments des deux parties (thèse et antithèse), puis conclure en appliquant la règle de droit (synthèse). La philosophie n'est donc pas une matière isolée : elle est la \emph{boîte à outils intellectuelle} de toute votre formation technique.

\begin{aretenir}[L'essentiel de la section]
\begin{itemize}
\item La dissertation forme quatre compétences transférables : analyser, structurer, argumenter, synthétiser.
\item Rapports techniques, notes de synthèse, lettres de motivation et prises de parole suivent le même schéma : problématique $\rightarrow$ analyse structurée $\rightarrow$ preuves $\rightarrow$ conclusion actionnable.
\item L'esprit critique permet de débusquer les arguments fallacieux dans les médias et les débats.
\item L'éthique de la discussion (principe de charité, distinction personne/idée, refus de l'autorité nue) fait de la philosophie une école de citoyenneté et de professionnalisme.
\end{itemize}
\end{aretenir}

Ainsi se referme cette leçon de consolidation : la dissertation, loin d'être un rituel scolaire, est une gymnastique de la liberté. Celui qui sait se poser les bonnes questions, ordonner ses idées et justifier ses choix ne sera jamais ni un exécutant docile ni un citoyen manipulable — mais un technicien qui pense et un citoyen qui décide.

\newpage

\coursec{Activité d'intégration}

\coursec{Consolidation des acquis de la dissertation philosophique}

\courssub{Objectifs de l'activité}

À l'issue de cette activité d'intégration, l'élève doit être capable de :
\begin{itemize}
\item analyser rigoureusement un sujet de dissertation (dégager les termes clés, l'enjeu, la problématique) ;
\item construire un plan dialectique cohérent (thèse, antithèse, synthèse) et articulé ;
\item mobiliser des documents variés (texte philosophique, données statistiques, image, texte juridique) comme arguments et exemples ;
\item rédiger, dans les conditions de l'examen, une copie complète comportant introduction, développement et conclusion ;
\item évaluer une copie à l'aide d'une grille critériée et justifier son jugement ;
\item transférer les acquis de la dissertation à une situation citoyenne authentique : le débat sur le développement du Cameroun.
\end{itemize}

\courssub{Situation}

Dans le cadre de la semaine de la philosophie organisée par ton lycée technique de Douala, le club « Les Jeunes Philosophes du Wouri » décide d'organiser \textbf{Le Grand Débat Philosophique} sur le thème : \emph{« Le Développement du Cameroun à l'épreuve de la Technique »}. Le proviseur a accepté que la production écrite du débat prenne la forme d'une dissertation complète, rédigée dans les conditions du Probatoire technique, puis publiée dans le journal du lycée.

Le sujet imposé par le jury est le suivant :

\begin{center}
\textbf{« La maîtrise de la technique est-elle la condition sine qua non du développement durable au Cameroun ? »}
\end{center}

Pour alimenter la réflexion, le jury remet à chaque groupe un \textbf{Dossier Documentaire} de cinq documents :

\begin{itemize}
\item \textbf{Document 1 (texte philosophique)} : extrait de Martin Heidegger, \emph{La question de la technique} (1953) : « L'essence de la technique n'est rien de technique. La technique est un mode de dévoilement ; mais la technique moderne tend à réduire la nature à un "fonds" disponible, une réserve que l'homme exploite à volonté. »
\item \textbf{Document 2 (données statistiques, INS Cameroun)} : part du secteur industriel (y compris bâtiment et travaux publics) dans le PIB du Cameroun : environ 26 \% en 2010, environ 28 \% en 2015, environ 26 \% en 2020, environ 27 \% en 2023 ; le secteur primaire (agriculture, forêt, élevage) représente environ 17 \% du PIB mais emploie près de 50 \% de la population active.
\item \textbf{Document 3 (photographie légendée)} : le chantier du barrage hydroélectrique de Nachtigal sur le fleuve Sanaga (capacité prévue de 420 MW, investissement d'environ 1 200 milliards de FCFA), avec, en arrière-plan, une zone de forêt défrichée ; la légende précise que le Cameroun perd chaque année environ 200 000 hectares de couvert forestier.
\item \textbf{Document 4 (texte juridique)} : article 9 de la loi n\textdegree 96/12 du 5 août 1996 portant loi-cadre relative à la gestion de l'environnement au Cameroun : « Toute étude d'impact sur l'environnement est obligatoire avant la réalisation des grands projets de développement. »
\item \textbf{Document 5 (texte philosophique)} : extrait de Fabien Eboussi Boulaga : « L'urgence de penser : un développement qui ne fait que copier les techniques venues d'ailleurs, sans les repenser à partir de nos besoins, de nos valeurs et de nos communautés, n'est qu'une modernisation de la dépendance. »
\end{itemize}

Le débat final portera sur la question : \emph{que doit être le philosophe-citoyen camerounais face à la technique ?}

\courssub{Consignes}

\textbf{Séance 1 — Préparation collective (groupes de 4 élèves).}

\begin{enumerate}
\item \textbf{Analyse du sujet.} Souligne les termes clés du sujet (« maîtrise de la technique », « condition sine qua non », « développement durable », « au Cameroun »). Définis chacun d'eux en une phrase. Quelle présupposition le sujet contient-il ? Formule l'enjeu puis la problématique sous forme d'une question.
\item \textbf{Exploitation du dossier.} Pour chaque document, indique : (a) la thèse ou l'information principale ; (b) dans quelle partie du développement il peut servir d'argument ou d'exemple.
\item \textbf{Construction du plan.} Rédige sur la fiche de plan un plan dialectique détaillé (thèse / antithèse / synthèse), avec pour chaque partie : l'idée directrice, deux arguments, un auteur ou un document du dossier, un exemple camerounais.
\item \textbf{Répartition des tâches.} Binôme A : introduction ; Binôme B : première partie ; Binôme C : deuxième partie ; Binôme A : troisième partie et conclusion. Chaque binôme rédige un brouillon de sa partie.
\end{enumerate}

\textbf{Séance 2 — Rédaction et évaluation.}

\begin{enumerate}
\setcounter{enumi}{4}
\item \textbf{Mise en commun (15 min)} : lecture croisée des brouillons, vérification de la cohérence des transitions et de l'unité de la problématique.
\item \textbf{Rédaction individuelle (2 h, conditions Bac)} : rédige seul(e) la copie complète (introduction, trois parties, conclusion), en intégrant au moins trois documents du dossier et deux références philosophiques.
\item \textbf{Correction par les pairs} : échange ta copie avec un élève d'un autre groupe et évalue-la sur la grille officielle (analyse du sujet /4, cohérence du plan /4, qualité de l'argumentation /6, mobilisation des références et exemples /3, langue et présentation /3). Justifie chaque note par une observation écrite.
\item \textbf{Débriefing collectif} : réponds par écrit en dix lignes à la question : « Qu'avons-nous appris sur le philosophe-citoyen camerounais ? »
\end{enumerate}

\courssub{Ressources à mobiliser}

\begin{aretenir}[Les outils de la dissertation à réactiver]
\begin{itemize}
\item \textbf{Analyse du sujet} : repérer les termes clés, dégager le présupposé, formuler l'enjeu et la problématique.
\item \textbf{Structure de l'introduction} : accroche (amorce) $\rightarrow$ présentation du sujet et définition des termes $\rightarrow$ enjeu $\rightarrow$ problématique $\rightarrow$ annonce du plan.
\item \textbf{Plan dialectique} : thèse (la technique comme condition du développement), antithèse (ses limites et ses dangers), synthèse (une maîtrise critique et responsable de la technique).
\item \textbf{Paragraphe type} : idée $\rightarrow$ argument $\rightarrow$ exemple $\rightarrow$ référence $\rightarrow$ phrase de transition.
\item \textbf{Conclusion} : réponse nuancée à la problématique, bilan du parcours, ouverture.
\item \textbf{Références utiles} : Heidegger (la technique comme dévoilement et danger), Eboussi Boulaga (penser le développement à partir de l'Afrique), la distinction développement / croissance, la notion de développement durable (satisfaire les besoins présents sans compromettre ceux des générations futures).
\end{itemize}
\end{aretenir}

\courssub{Démarche et corrigé commenté}

\begin{exoresolu}[Corrigé commenté du Grand Débat]
\textbf{Étape 1 — Analyse du sujet.} Les termes clés sont : \cle{maîtrise de la technique} (capacité à produire, utiliser et contrôler les outils et procédés techniques, et non à les subir) ; \cle{condition sine qua non} (condition absolument nécessaire, sans laquelle rien n'est possible) ; \cle{développement durable} (développement économique, social et respectueux de l'environnement, pensé pour les générations futures) ; \cle{au Cameroun} (ancrage obligatoire des exemples). Le sujet présuppose qu'il pourrait exister un lien de nécessité entre technique et développement durable : c'est précisément ce lien qu'il faut interroger. \emph{Enjeu} : faut-il faire de la technique la clé unique du développement camerounais, ou faut-il la subordonner à d'autres exigences (éthiques, écologiques, culturelles) ? \emph{Problématique} : la technique suffit-elle à fonder un développement durable, ou doit-elle être maîtrisée par une pensée critique pour ne pas détruire ce qu'elle prétend développer ?

\tcblower

\textbf{Étape 2 — Exploitation du dossier.} Document 1 (Heidegger) : argument pour l'antithèse — la technique moderne réduit la nature à une réserve exploitable. Document 2 (INS) : exemple chiffré pour la thèse — la stagnation industrielle (environ 26 \% à 27 \% du PIB entre 2010 et 2023) montre un déficit de maîtrise technique ; mais l'importance de l'agriculture (50 \% des actifs) rappelle que le développement n'est pas qu'industriel. Document 3 (Nachtigal) : exemple ambivalent — prouesse technique (420 MW, 1 200 milliards de FCFA) mais coût écologique (déforestation, environ 200 000 ha perdus par an). Document 4 (loi de 1996) : argument pour la synthèse — le droit encadre la technique par l'étude d'impact. Document 5 (Eboussi Boulaga) : argument de synthèse — il faut repenser la technique à partir des besoins et des valeurs camerounais.

\textbf{Étape 3 — Plan détaillé attendu.}
\begin{itemize}
\item \textbf{I. Thèse} : la maîtrise de la technique apparaît comme la condition du développement. Argument 1 : nul développement sans infrastructures, énergie, industrie (exemple : Nachtigal, 420 MW pour électrifier le pays ; données INS). Argument 2 : la technique libère l'homme des contraintes naturelles (référence : Descartes, « nous rendre comme maîtres et possesseurs de la nature »).
\item \textbf{II. Antithèse} : la technique seule ne suffit pas et peut même détruire. Argument 1 : Heidegger — la technique moderne épuise la nature (exemple : déforestation visible sur le document 3). Argument 2 : une technique importée sans réflexion crée la dépendance (Eboussi Boulaga) ; croissance n'est pas développement.
\item \textbf{III. Synthèse} : la maîtrise de la technique est une condition nécessaire mais non suffisante ; elle doit être éclairée par l'éthique, le droit et la pensée critique. Argument 1 : la loi de 1996 impose l'étude d'impact — le droit encadre la technique. Argument 2 : le développement durable exige une technique appropriée aux besoins locaux (exemple : transformation locale du cacao, énergies solaires décentralisées dans les villages du Nord).
\end{itemize}

\textbf{Étape 4 — Éléments de rédaction.} \emph{Accroche possible} : « Chaque matin, à Douala, des millions de Camerounais dépendent de la technique : motos-taxis, téléphones mobiles, électricité. Pourtant, ce monde technique que nous utilisons, le maîtrisons-nous vraiment ? » \emph{Réponse à la problématique (conclusion)} : la maîtrise de la technique est bien une condition nécessaire du développement durable au Cameroun, mais elle n'en est pas la condition suffisante : sans pensée critique, sans cadre juridique et sans souci des générations futures, la technique devient un instrument de destruction et de dépendance. \emph{Ouverture} : « Le philosophe-citoyen camerounais n'est-il pas précisément celui qui apprend à penser la technique avant de la consommer ? »

\textbf{Étape 5 — Auto-évaluation.} Une copie excellente : définit les quatre termes clés, annonce clairement le plan, cite au moins trois documents avec exactitude (chiffres INS, 420 MW, loi de 1996), articule les parties par des transitions explicites, et conclut par une réponse nuancée (« nécessaire mais non suffisante »). Erreurs à sanctionner : hors-sujet (disserter sur « la technique » en général sans le Cameroun ni le développement durable), paraphrase des documents sans argumentation, absence de problématique.
\end{exoresolu}

\begin{attention}[Pièges fréquents dans cette activité]
\begin{itemize}
\item Confondre \emph{croissance économique} (augmentation du PIB) et \emph{développement durable} (progrès économique, social et écologique) : le document 2 sert précisément à les distinguer.
\item Traiter le sujet comme un exposé géographique : il faut \emph{problématiser}, pas décrire le Cameroun.
\item Juxtaposer les documents sans les articuler à un argument : chaque document doit \emph{prouver} une idée.
\item Oublier l'expression « condition sine qua non » : toute la dissertation porte sur la nécessité du lien entre technique et développement.
\end{itemize}
\end{attention}

\courssub{Bilan de l'activité}

Cette activité a permis de mettre en évidence que la dissertation philosophique n'est pas un exercice scolaire artificiel, mais un \cle{outil de citoyenneté} : analyser un sujet, c'est apprendre à ne pas subir les mots ; construire un plan, c'est apprendre à ordonner sa pensée avant d'agir ; argumenter avec des documents, c'est fonder son opinion sur des preuves plutôt que sur des rumeurs. Les élèves ont également découvert que les notions du programme (la technique, le développement, la responsabilité) éclairent des problèmes concrets du Cameroun — énergie, forêts, industrialisation — et que la confrontation des copies par la correction croisée développe l'esprit critique et l'honnêteté intellectuelle. Le philosophe-citoyen camerounais apparaît ainsi comme celui qui maîtrise à la fois l'art de penser et le souci du bien commun : penser la technique, c'est déjà commencer à la maîtriser.

\newpage

\coursec{Méthodes et exercices résolus}

\coursec{Leçon 02 : Consolidation des acquis de la dissertation philosophique}

\courssub{Méthode 1 : Analyser rigoureusement le sujet (travail pratique de consolidation)}

\begin{methode}[Analyser un sujet de dissertation en 5 étapes]
\begin{enumerate}
\item \textbf{Lire le sujet trois fois} : une première lecture globale, une deuxième pour souligner les mots-clés, une troisième pour repérer la forme du sujet (question directe, citation, affirmation).
\item \textbf{Définir chaque terme important} : donner le sens courant, puis le sens philosophique du mot. Exemple : \emph{liberté} = absence de contrainte (sens courant) ; pouvoir de se déterminer soi-même (sens philosophique).
\item \textbf{Dégager la problématique} : transformer le sujet en question directrice. Chercher la tension entre deux thèses opposées (par exemple : la liberté est-elle un fait ou une conquête ?).
\item \textbf{Formuler les enjeux} : se demander \emph{pourquoi cette question est importante} pour l'homme, la société, la cité.
\item \textbf{Choisir le plan adapté} : plan dialectique (thèse / antithèse / synthèse) pour une question fermée ; plan thématique pour une notion ; plan analytique pour une citation.
\end{enumerate}
\textbf{Réflexe à retenir} : jamais de développement sans problématique annoncée dans l'introduction. La problématique est la \cle{clé de voûte} de toute la copie.
\end{methode}

\begin{exoresolu}[Analyse du sujet : « Le travail est-il une malédiction ? »]
\textbf{Énoncé.} On donne le sujet : \emph{« Le travail est-il une malédiction ? »} (sujet type Probatoire technique). Effectue l'analyse complète du sujet : définition des termes, problématique, enjeux, choix du plan.
\tcblower
\textbf{Solution détaillée.}

\textbf{1. Forme du sujet.} Il s'agit d'une \emph{question directe} portant sur la notion de travail. La réponse attendue ne peut être ni un simple « oui » ni un simple « non » : il faut discuter.

\textbf{2. Définition des termes.}
\begin{itemize}
\item \textbf{Travail} : au sens courant, activité rémunérée (le menuisier de Douala, l'enseignant, l'informaticien). Au sens philosophique, activité par laquelle l'homme transforme la nature et se transforme lui-même (Hegel).
\item \textbf{Malédiction} : au sens courant, punition, souffrance subie. Le mot renvoie au récit biblique : « Tu gagneras ton pain à la sueur de ton front ». Le travail y apparaît comme une peine infligée à l'homme.
\end{itemize}

\textbf{3. Problématique.} La tension est la suivante : si le travail fatigue, aliène et contraint, il semble être une malédiction ; mais s'il libère l'homme de la nature, le socialise et le réalise, il est au contraire une chance. D'où la problématique : \emph{le travail est-il seulement une contrainte subie, ou est-il aussi le moyen par lequel l'homme se libère et s'accomplit ?}

\textbf{4. Enjeux.} La question est concrète pour l'élève de Première technique qui se prépare à la vie professionnelle : faut-il envisager son futur métier comme une corvée ou comme un chemin d'épanouissement ? L'enjeu est donc à la fois personnel (sens de la vie) et social (dignité du travailleur).

\textbf{5. Choix du plan.} La question fermée (« est-il... ? ») appelle un \textbf{plan dialectique} :
\begin{itemize}
\item Thèse : le travail est une malédiction (fatigue, exploitation, aliénation — Marx).
\item Antithèse : le travail est libération (maîtrise de la nature, reconnaissance sociale — Hegel).
\item Synthèse : le travail est ambivalent ; tout dépend des conditions dans lesquelles il s'exerce.
\end{itemize}

\textbf{Conclusion de l'analyse.} Le sujet invite à dépasser l'opposition simple pour montrer que le travail n'est une malédiction que lorsqu'il est déshumanisé ; accompli dans la dignité, il devient moyen d'émancipation.
\end{exoresolu}

\courssub{Méthode 2 : Construire l'architecture du développement}

\begin{methode}[Bâtir un plan dialectique solide]
\begin{enumerate}
\item \textbf{Rédiger d'abord le plan détaillé au brouillon} : chaque partie = une idée directrice + deux ou trois arguments + un auteur + un exemple.
\item \textbf{Partie I (thèse)} : exposer la position la plus immédiate, celle du sens commun, avec ses arguments les plus forts.
\item \textbf{Partie II (antithèse)} : renverser la perspective ; montrer les limites de la thèse et défendre la position opposée.
\item \textbf{Partie III (synthèse)} : dépasser l'opposition, répondre personnellement et justifiée à la problématique. La synthèse n'est pas un compromis mou mais une \emph{réponse argumentée}.
\item \textbf{Articuler les parties} : annoncer à la fin de chaque partie ce qui va suivre (« Mais cette conception ne rend-elle pas compte que d'un aspect du travail ? »).
\item \textbf{Équilibrer} : les trois parties doivent avoir une ampleur comparable (environ un tiers de la copie chacune).
\end{enumerate}
\textbf{Formule à retenir} : une partie = \cle{une idée}, un argument = \cle{une preuve}, un exemple = \cle{une illustration} (jamais une preuve à lui seul).
\end{methode}

\begin{exoresolu}[Construction du plan : « La technique nous libère-t-elle ? »]
\textbf{Énoncé.} Construis le plan détaillé (idées directrices, arguments, auteurs, exemples) d'une dissertation sur le sujet : \emph{« La technique nous libère-t-elle ? »}
\tcblower
\textbf{Solution détaillée.}

\textbf{Analyse préalable.} \emph{Technique} = ensemble des outils et procédés inventés par l'homme ; \emph{libérer} = rendre autonome, dégager d'une contrainte. Problématique : \emph{la technique, créée pour servir l'homme, ne risque-t-elle pas de le dominer à son tour ?}

\textbf{Partie I — Thèse : la technique libère l'homme.}
\begin{itemize}
\item Idée directrice : la technique soulage l'homme des contraintes naturelles.
\item Argument 1 : elle économise l'effort et le temps (Aristote : si les outils faisaient eux-mêmes leur ouvrage, les maîtres n'auraient pas besoin d'esclaves).
\item Argument 2 : elle repousse la maladie et la mort (médecine, vaccins).
\item Exemple : la moto-taxi et le téléphone mobile relient les villages camerounais isolés (de Mora à Kribi) aux marchés et aux services.
\end{itemize}

\textbf{Partie II — Antithèse : la technique peut asservir.}
\begin{itemize}
\item Idée directrice : l'outil créé par l'homme se retourne contre lui.
\item Argument 1 : la dépendance (l'usager ne maîtrise plus l'outil qui le sert).
\item Argument 2 : l'aliénation au travail (Marx : l'ouvrier devient un simple rouage de la machine).
\item Exemple : le jeune qui passe ses journées sur les réseaux sociaux au détriment de ses études ; l'employé surveillé en permanence par ses outils numériques.
\end{itemize}

\textbf{Partie III — Synthèse : la technique est un instrument, sa valeur dépend de l'usage.}
\begin{itemize}
\item Idée directrice : la technique est ambivalente ; c'est la finalité que l'homme lui donne qui la rend libératrice ou esclavagiste.
\item Argument : un couteau peut soigner (chirurgien) ou blesser (agresseur) ; la responsabilité est humaine, non technique.
\item Réponse à la problématique : la technique libère à condition que l'homme reste maître de ses fins — d'où la nécessité de l'éducation et de la réflexion éthique.
\end{itemize}

\textbf{Conclusion.} Le plan est dialectique, équilibré, et chaque partie s'articule à la suivante par une question de transition. Il ne reste plus qu'à rédiger.
\end{exoresolu}

\courssub{Méthode 3 : Rédiger l'introduction et la conclusion}

\begin{methode}[L'introduction en 4 temps et la conclusion en 3 temps]
\textbf{Introduction (environ 10 à 15 lignes) :}
\begin{enumerate}
\item \textbf{Amorce (accroche)} : un constat général, un fait d'expérience ou un lieu commun qui mène naturellement au sujet. Jamais de définition de dictionnaire en ouverture.
\item \textbf{Analyse des termes du sujet} : définir brièvement les mots-clés et montrer leur ambiguïté.
\item \textbf{Problématique} : formuler la question directrice, souvent sous forme d'interrogation explicite.
\item \textbf{Annonce du plan} : « Nous verrons d'abord que..., puis que..., enfin que... ».
\end{enumerate}
\textbf{Conclusion (environ 8 à 12 lignes) :}
\begin{enumerate}
\item \textbf{Bilan} : rappeler le chemin parcouru sans recopier l'annonce du plan.
\item \textbf{Réponse nette à la problématique} : trancher et justifier en une ou deux phrases.
\item \textbf{Ouverture} : poser une nouvelle question qui prolonge la réflexion (sans ouvrir un nouveau sujet).
\end{enumerate}
\textbf{Réflexe} : introduction et conclusion doivent se \cle{répondre} : la conclusion répond exactement à la question posée dans l'introduction.
\end{methode}

\begin{exoresolu}[Rédaction d'une introduction complète]
\textbf{Énoncé.} Rédige l'introduction complète de la dissertation : \emph{« Faut-il toujours dire la vérité ? »}
\tcblower
\textbf{Solution détaillée (modèle rédigé).}

« Dans la vie quotidienne, chacun reconnaît volontiers que la vérité est un bien : un ami qui ment perd notre confiance, un commerçant de Marché Central qui trompe sur la qualité perd sa clientèle. Dire la vérité semble donc être un devoir absolu. Pourtant, l'expérience montre aussi des situations où la vérité blesse : faut-il annoncer brutalement à un malade la gravité de son état ? Faut-il révéler un secret qui détruirait une famille ?

La \emph{vérité} désigne la conformité de ce que l'on dit avec ce qui est ; l'expression \emph{toujours} pose la question de l'universalité du devoir : la règle admet-elle des exceptions ?

Le problème est donc le suivant : le devoir de dire la vérité est-il absolu, sans aucune exception, ou doit-il céder devant d'autres exigences, comme la bonté ou la prudence ?

Nous examinerons d'abord que la vérité est le fondement de toute relation humaine et de toute société ; nous montrerons ensuite qu'il existe des situations où dire la vérité cause un mal réel ; nous soutiendrons enfin que le devoir de vérité demeure, mais qu'il doit s'exercer avec tact, c'est-à-dire selon le moment et la manière. »

\textbf{Justification de la démarche.}
\begin{itemize}
\item L'amorce part d'un constat concret et camerounais (le commerçant), sans définition de dictionnaire.
\item Les deux termes-clés (\emph{vérité}, \emph{toujours}) sont définis et leur tension est dégagée.
\item La problématique est formulée par une question explicite.
\item Le plan dialectique est annoncé en trois étapes clairement numérotées par « d'abord / ensuite / enfin ».
\end{itemize}
\textbf{Conclusion attendue (rappel)} : elle devra répondre : oui, il faut dire la vérité, mais non n'importe comment ni n'importe quand — et s'ouvrir, par exemple, sur la question : « Le silence est-il parfois une forme de mensonge ? »
\end{exoresolu}

\courssub{Méthode 4 : Argumenter : preuves, exemples et références}

\begin{methode}[Construire un paragraphe argumentatif (schéma A-R-E-A)]
\begin{enumerate}
\item \textbf{Affirmation} : énoncer clairement l'idée défendue dans le paragraphe (une seule idée).
\item \textbf{Raisonnement} : justifier l'idée par un argument logique, en enchaînant les étapes avec des connecteurs (\emph{car, en effet, par conséquent, or, donc}).
\item \textbf{Exemple} : illustrer par un cas précis, de préférence tiré de la vie courante ou du contexte camerounais, puis \emph{commenter} l'exemple pour montrer qu'il vérifie l'idée.
\item \textbf{Appui d'auteur} : citer brièvement une thèse de philosophe (nom + idée, pas de longue citation apprise par cœur) : « Comme le montre Kant, agir par devoir, c'est... ».
\end{enumerate}
\textbf{Réflexes de style} : écrire « nous » plutôt que « je » ; éviter « selon moi » en début de phrase ; bannir les abréviations et le langage SMS ; soigner l'orthographe des noms propres (\emph{Platon}, \emph{Descartes}, \emph{Nietzsche}).
\textbf{À retenir} : un exemple n'est jamais une preuve : il \cle{illustre}, seul le raisonnement \cle{prouve}.
\end{methode}

\begin{exoresolu}[Rédaction d'un paragraphe argumentatif]
\textbf{Énoncé.} Dans le cadre du sujet \emph{« La liberté est-elle un donné ou une conquête ? »}, rédige un paragraphe complet (schéma A-R-E-A) défendant l'idée suivante : « La liberté n'est pas un donné naturel, elle se conquiert par l'éducation et l'effort. »
\tcblower
\textbf{Solution détaillée (modèle rédigé).}

« La liberté n'est pas un donné naturel : elle se conquiert par l'éducation et l'effort. En effet, l'enfant qui naît est entièrement dépendant : il ne choisit ni sa langue, ni sa famille, ni ses besoins. Or, être libre, c'est pouvoir déterminer soi-même ses actes en connaissance de cause ; cela suppose donc d'avoir appris à connaître, à juger et à se maîtriser. Par conséquent, la liberté n'est pas un point de départ, mais un aboutissement. L'exemple de l'élève camerounais qui apprend un métier au lycée technique l'illustre : au début de sa formation, il exécute les gestes qu'on lui impose ; à la fin, capable de comprendre et de décider, il devient un technicien autonome, c'est-à-dire plus libre. Cet exemple montre que l'apprentissage transforme la dépendance en autonomie. C'est précisément la thèse de Kant : l'illumination, écrit-il, est la sortie de l'homme hors de la minorité dont il est lui-même responsable — autrement dit, on ne naît pas libre, on le devient en osant penser par soi-même. »

\textbf{Justification de la démarche.}
\begin{itemize}
\item \textbf{Affirmation} : la première phrase pose l'idée sans détour.
\item \textbf{Raisonnement} : enchaînement logique marqué par « En effet / Or / Par conséquent » : dépendance de naissance $\Rightarrow$ définition de la liberté $\Rightarrow$ la liberté comme aboutissement.
\item \textbf{Exemple} : concret (formation technique), puis explicitement commenté (« Cet exemple montre que... »).
\item \textbf{Appui d'auteur} : Kant est cité brièvement, avec son idée exacte, et relié à la thèse du paragraphe.
\end{itemize}
\textbf{Conclusion.} Le paragraphe est complet, cohérent et directement utilisable dans la partie « antithèse » du plan.
\end{exoresolu}

\courssub{Méthode 5 : Appliquer la dissertation aux situations concrètes de la vie courante}

\begin{methode}[Traiter une situation concrète en 4 étapes]
\begin{enumerate}
\item \textbf{Décrire la situation} avec précision : qui ? quoi ? où ? quel est le dilemme ou le problème ?
\item \textbf{Dégager la notion philosophique} en jeu (liberté, devoir, justice, vérité, travail, bonheur...).
\item \textbf{Mobiliser les notions de l'unité} : confronter deux thèses d'auteurs ou deux points de vue opposés sur la situation.
\item \textbf{Rédiger et justifier une conclusion} : prendre position de manière argumentée et proposer une conduite concrète.
\end{enumerate}
\textbf{Réflexe} : la situation concrète n'est pas une anecdote à raconter ; c'est un \cle{cas} que la philosophie doit éclairer. Toujours passer du particulier au général, puis revenir au particulier.
\end{methode}

\begin{exoresolu}[Situation concrète : le devoir et l'intérêt]
\textbf{Énoncé.} Situation : \emph{Manga, élève en classe de Première technique, trouve dans la cour du lycée un portefeuille contenant \num{50000} FCFA. Personne ne l'a vu le ramasser. Son camarade lui dit : « Garde l'argent, tu as de la chance. » Manga hésite.} Analyse cette situation à l'aide des notions de la philosophie morale et rédige une conclusion justifiée.
\tcblower
\textbf{Solution détaillée.}

\textbf{1. Description.} Manga est seul face à un choix : restituer l'argent (devoir d'honnêteté) ou le garder (intérêt personnel). L'absence de témoin est essentielle : aucune sanction n'est à craindre.

\textbf{2. Notion en jeu.} Il s'agit du rapport entre le \emph{devoir} et l'\emph{intérêt} : doit-on faire le bien uniquement lorsque l'on est surveillé ?

\textbf{3. Confrontation des thèses.}
\begin{itemize}
\item \textbf{Thèse de l'intérêt} : le camarade raisonne en calculant les conséquences : sans témoin, garder l'argent ne coûte rien et rapporte beaucoup. C'est une morale des résultats.
\item \textbf{Thèse du devoir} : pour Kant, une action n'a de valeur morale que si elle est faite par devoir, et non par intérêt ou par crainte. Manga doit restituer l'argent parce que c'est son devoir, même si personne ne le saura jamais. Le critère : puis-je vouloir que ma conduite devienne une loi universelle ? Or personne ne peut vouloir vivre dans une société où chacun garde ce qu'il trouve.
\item \textbf{Point de vue de la conscience} : même sans témoin extérieur, il reste un témoin intérieur : la conscience morale, qui juge nos actes.
\end{itemize}

\textbf{4. Conclusion rédigée et justifiée.} « Manga doit restituer le portefeuille à l'administration du lycée. Cette conduite se justifie doublement : d'une part, garder l'argent reviendrait à s'approprier le bien d'autrui, ce qui, généralisé, rendrait toute confiance sociale impossible ; d'autre part, agir honnêtement sans témoin est précisément ce qui distingue l'homme moral de l'homme simplement prudent. L'honnêteté véritable ne se mesure pas devant les regards, mais dans le secret. »

\textbf{Conclusion de l'exercice.} La méthode a permis de passer d'un cas particulier (un portefeuille perdu) à une réflexion générale sur le devoir, puis de revenir au cas avec une décision justifiée.
\end{exoresolu}

\courssub{Méthode 6 : S'auto-évaluer et réviser sa copie (grille type Bac)}

\begin{methode}[La grille d'auto-évaluation en 6 points]
Avant de rendre la copie, vérifier successivement :
\begin{enumerate}
\item \textbf{Pertinence} : ai-je répondu au sujet posé (et non à un sujet voisin) ? Relire le sujet et comparer avec la problématique.
\item \textbf{Structure} : introduction en 4 temps ? trois parties équilibrées et articulées ? conclusion en 3 temps ?
\item \textbf{Argumentation} : chaque idée est-elle justifiée par un raisonnement et illustrée par un exemple commenté ?
\item \textbf{Références} : au moins deux ou trois auteurs cités correctement, sans erreur d'attribution ?
\item \textbf{Langue} : orthographe, accords, ponctuation, connecteurs logiques, « nous » de préférence à « je ».
\item \textbf{Présentation} : écriture lisible, alinéas au début de chaque partie, pas de ratures excessives, pas de titres numérotés dans la copie.
\end{enumerate}
\textbf{Réflexe de gestion du temps} : sur 4 heures, compter environ 30 min d'analyse et de plan détaillé, 2 h 45 de rédaction, 30 min de relecture. La relecture n'est \cle{jamais} facultative.
\end{methode}

\begin{exoresolu}[Correction d'une copie défaillante]
\textbf{Énoncé.} Voici l'introduction d'un élève au sujet \emph{« La conscience fait-elle la grandeur de l'homme ? »} : \og La conscience est un mot qui vient du latin \emph{conscientia}. Dans ce devoir, je vais parler de la conscience. D'abord je vais dire ce que c'est, ensuite je vais donner des exemples, et enfin je vais donner mon avis. \fg{} Relève les défauts de cette introduction à l'aide de la grille d'auto-évaluation, puis propose une version corrigée de la problématique et de l'annonce du plan.
\tcblower
\textbf{Solution détaillée.}

\textbf{Relevé des défauts (selon la grille).}
\begin{enumerate}
\item \textbf{Amorce défaillante} : l'étymologie en ouverture est un lieu commun interdit ; elle ne mène pas au problème.
\item \textbf{Pas d'analyse des termes} : ni \emph{conscience} ni \emph{grandeur} ne sont définis ; l'ambiguïté du mot conscience (conscience morale ? conscience de soi ?) n'est pas dégagée.
\item \textbf{Absence de problématique} : aucune question directrice n'est formulée ; le sujet est simplement répété.
\item \textbf{Annonce de plan scolaire et vide} : « je vais dire ce que c'est / des exemples / mon avis » n'annonce aucune idée ; c'est un plan de rédaction, non un plan philosophique.
\item \textbf{Style} : emploi de « je » et de « dans ce devoir », à proscrire.
\end{enumerate}

\textbf{Version corrigée (problématique + annonce du plan).}

Problématique : « La conscience — c'est-à-dire à la fois la connaissance que l'homme a de lui-même et la voix intérieure qui juge ses actes — est-elle ce qui élève l'homme au-dessus de l'animal, ou bien est-elle une source de souffrance dont les autres vivants sont préservés ? »

Annonce du plan : « Nous montrerons d'abord que la conscience distingue l'homme de l'animal et fonde sa dignité ; nous verrons ensuite qu'elle est aussi source d'angoisse, de remords et de malheur ; nous soutiendrons enfin que la grandeur de l'homme réside précisément dans cette capacité à se juger soi-même, qui le rend responsable de ses actes. »

\textbf{Conclusion.} L'introduction corrigée définit les termes, pose un véritable problème (grandeur ou fardeau ?) et annonce un plan dialectique dont chaque étape porte une idée précise : elle satisferait aux six points de la grille.
\end{exoresolu}

\begin{attention}[Les 5 erreurs qui coûtent des points au Bac]
\begin{enumerate}
\item \textbf{Le hors-sujet} : traiter la notion en général (« parler de la liberté ») au lieu de répondre à la question posée. Réflexe : relire le sujet après chaque partie et vérifier que la problématique y répond exactement.
\item \textbf{La dissertation-récitation} : aligner des citations apprises par cœur sans les relier au sujet ni les commenter. Une référence non expliquée ne rapporte rien ; une référence mal attribuée (confondre Kant et Marx) fait perdre des points.
\item \textbf{L'exemple non commenté} : raconter une anecdote (« mon voisin... ») sans montrer ce qu'elle illustre. Tout exemple doit être suivi d'une phrase du type « Cet exemple montre que... ».
\item \textbf{L'annonce de plan creuse} : écrire « dans un premier temps, dans un deuxième temps » sans énoncer les idées. Le correcteur doit connaître le contenu de chaque partie dès l'introduction.
\item \textbf{Négliger la forme} : orthographe des noms d'auteurs, ratures, absence d'alinéas, « je » au lieu de « nous », conclusion bâclée ou absente. Une copie sans conclusion est une copie inachevée : réserver impérativement les 10 dernières minutes pour la rédiger, même brièvement.
\end{enumerate}
\end{attention}

\newpage

\coursec{Exercices}

\courssub{Je vérifie mes connaissances}

\begin{exercice}[QCM méthodologique]
Pour chaque question, une seule réponse est correcte. Entoure la bonne lettre.
\begin{enumerate}
\item L'analyse d'un sujet de dissertation commence par :
\begin{enumerate}
\item[a)] la rédaction de l'introduction ;
\item[b)] la définition des termes du sujet et la recherche du problème ;
\item[c)] le choix du plan ;
\item[d)] la citation d'un philosophe.
\end{enumerate}
\item La problématique est :
\begin{enumerate}
\item[a)] une simple question quelconque ;
\item[b)] la reformulation du sujet ;
\item[c)] la question qui exprime le problème philosophique soulevé par le sujet ;
\item[d)] le titre de la première partie.
\end{enumerate}
\item Dans un plan dialectique, la troisième partie sert à :
\begin{enumerate}
\item[a)] répéter la première partie ;
\item[b)] dépasser l'opposition entre la thèse et l'antithèse ;
\item[c)] donner des exemples ;
\item[d)] présenter la biographie d'un auteur.
\end{enumerate}
\item La transition entre deux parties sert à :
\begin{enumerate}
\item[a)] remplir la copie ;
\item[b)] montrer la limite de la partie précédente et justifier le passage à la suivante ;
\item[c)] annoncer la conclusion ;
\item[d)] définir les termes du sujet.
\end{enumerate}
\item Une bonne conclusion de dissertation doit :
\begin{enumerate}
\item[a)] introduire un exemple nouveau ;
\item[b)] répondre clairement à la problématique et éventuellement ouvrir sur une question nouvelle ;
\item[c)] recopier l'introduction ;
\item[d)] donner l'avis du correcteur.
\end{enumerate}
\end{enumerate}
\end{exercice}

\begin{exercice}[Vrai ou faux justifié]
Réponds par Vrai ou Faux, puis justifie ta réponse en deux ou trois phrases.
\begin{enumerate}
\item Il est possible de rédiger l'introduction au brouillon avant d'avoir construit le plan détaillé.
\item Un exemple concret peut à lui seul tenir lieu d'argument dans une dissertation.
\item La thèse est l'idée défendue ; l'argument est la raison qui la justifie.
\item Le hors-sujet consiste à traiter un problème voisin de celui du sujet sans traiter le sujet lui-même.
\item Dans la conclusion, il est obligatoire de citer un nouveau philosophe.
\end{enumerate}
\end{exercice}

\begin{exercice}[Définitions]
Définis précisément, en deux ou trois lignes chacune, les notions méthodologiques suivantes : \cle{problématique}, \cle{plan dialectique}, \cle{transition}, \cle{amorce} (ou accroche), \cle{argumentation}.
\end{exercice}

\begin{exercice}[Reconnaître les étapes d'une introduction]
Voici, dans le désordre, les cinq éléments d'une introduction réussie au sujet « La technique nous rend-elle libres ? » :
\begin{enumerate}
\item[a)] « La technique nous rend-elle réellement libres, ou nous enchaîne-t-elle à de nouvelles dépendances ? »
\item[b)] « Nous traiterons d'abord de la technique comme libération, puis de la technique comme nouvelle servitude, enfin des conditions d'un usage libre de la technique. »
\item[c)] « Au quotidien, le téléphone portable transforme la vie des jeunes Camerounais, de Yaoundé à Maroua. »
\item[d)] « La technique désigne l'ensemble des procédés et outils par lesquels l'homme transforme la nature ; la liberté est le pouvoir d'agir selon sa propre volonté. »
\item[e)] « Or, si la technique supprime les contraintes naturelles, elle crée aussi des dépendances : d'où la question de savoir si elle libère vraiment. »
\end{enumerate}
Remets ces éléments dans l'ordre logique d'une introduction et nomme chaque étape (amorce, définitions, problème, problématique, annonce du plan).
\end{exercice}

\courssub{Je m'entraîne}

\begin{exercice}[Pièges à éviter : cinq introductions fautives]
Voici cinq débuts d'introduction rédigés par des élèves sur le sujet « Le travail est-il une contrainte ? ». Identifie la faute méthodologique commise dans chacun (hors-sujet, absence de problématique, plan invisible, définitions absentes, style oral) et propose une correction.
\begin{enumerate}
\item « Moi je pense que le travail c'est bien, parce que chez nous à Bafoussam tout le monde dit que l'oisiveté est mère de tous les vices. Je vais vous expliquer ça. »
\item « Le travail est une question très intéressante. Les philosophes en ont beaucoup parlé. Nous allons voir cela dans ce devoir. »
\item « Le bonheur est ce que tous les hommes recherchent. Aristote disait que le bonheur est la fin suprême. Le bonheur dépend-il de nous ? »
\item « Le travail, du latin \emph{tripalium}, instrument de torture, est une activité humaine. La contrainte est une obligation. Le travail est-il une contrainte ? » (le texte s'arrête là et le développement commence sans annonce de plan)
\item « Depuis la nuit des temps, les hommes travaillent. Le travail est-il une contrainte ? Nous verrons que oui, ensuite que non, ensuite les deux. »
\end{enumerate}
\end{exercice}

\begin{exercice}[Banque d'exemples camerounais]
Pour chacune des dix notions du programme de Première technique — \cle{Technique}, \cle{Travail}, \cle{État}, \cle{Nature}, \cle{Culture}, \cle{Justice}, \cle{Liberté}, \cle{Vérité}, \cle{Art}, \cle{Religion} — propose \textbf{trois exemples camerounais précis} (faits, institutions, œuvres, situations de la vie courante) que tu pourrais mobiliser dans une dissertation. Présente tes réponses dans un tableau à deux colonnes : « Notion » et « Exemples ».
\end{exercice}

\begin{exercice}[Quinze minutes chrono]
En 15 minutes maximum par sujet, rédige au brouillon, pour chacun des cinq sujets suivants : (1) la définition des termes, (2) la problématique, (3) l'annonce du plan en trois parties.
\begin{enumerate}
\item La culture nous détache-t-elle de la nature ?
\item La liberté consiste-t-elle à faire ce que l'on veut ?
\item L'art a-t-il pour fonction d'être utile ?
\item La religion est-elle une affaire privée ?
\item La vérité est-elle relative à chacun ?
\end{enumerate}
\end{exercice}

\begin{exercice}[Construire une transition]
Un élève vient d'achever la première partie de sa dissertation sur le sujet « L'État a-t-il le devoir de garantir le bonheur de ses citoyens ? », dans laquelle il a soutenu que l'État doit assurer les conditions matérielles du bien-être (sécurité, écoles, hôpitaux). Il s'apprête à montrer, dans la deuxième partie, que le bonheur relève d'un choix personnel qui échappe à l'État. Rédige la transition complète (quatre à six lignes) entre ces deux parties, en respectant la structure : rappel du résultat acquis, mise en évidence de sa limite, question ouvrant la partie suivante.
\end{exercice}

\courssub{Je me prépare au Bac}

\begin{exercice}[Sujet type Probatoire technique : analyse complète]
Sujet : \textbf{« L'État a-t-il le devoir de garantir le bonheur de ses citoyens ? »}
\begin{enumerate}
\item Définis chacun des termes du sujet : « État », « devoir », « garantir », « bonheur », « citoyen ».
\item Formule le problème philosophique soulevé par ce sujet, puis rédige une problématique claire sous forme d'une question.
\item Construis un plan détaillé en trois parties (plan dialectique), en indiquant pour chaque partie la thèse défendue, deux arguments et un exemple camerounais précis (par exemple : la gratuité de la santé des enfants de moins de cinq ans, les programmes de lutte contre la pauvreté, le rôle des coutumes locales dans le bonheur des communautés).
\item Rédige l'introduction complète de la dissertation (quinze à vingt lignes), avec amorce, définitions, problème, problématique et annonce du plan.
\item Rédige la conclusion (huit à dix lignes) répondant à la problématique et ouvrant sur une question nouvelle.
\end{enumerate}
\end{exercice}

\begin{exercice}[Sujet type Bac technique : dissertation complète en devoir maison]
Sujet : \textbf{« Le travail technique est-il une aliénation ? »}
\par Tu disposes de 3 heures, en conditions d'examen, sans document. Ta copie doit comporter : une introduction complète, un développement en trois parties articulées par des transitions rédigées, et une conclusion.
\begin{enumerate}
\item Avant de rédiger, établis au brouillon l'analyse du sujet et le plan détaillé (durée conseillée : 30 minutes).
\item Rédige la dissertation en mobilisant au moins deux références philosophiques étudiées dans l'unité et deux exemples tirés de la vie technique et professionnelle au Cameroun (atelier, chantier, usine, exploitation agricole, secteur informel).
\item Après la rédaction, auto-évalue ta copie à l'aide de la grille suivante, en t'attribuant une note sur 4 pour chaque critère : respect du sujet ; qualité de la problématique ; solidité des arguments ; pertinence des exemples ; clarté de l'expression.
\end{enumerate}
\end{exercice}

\begin{exercice}[Note de synthèse philosophique pour des décideurs]
Contexte : le conseil municipal d'une ville camerounaise hésite à installer des caméras de surveillance dans les quartiers. On te demande, en tant que citoyen formé à la philosophie, une \textbf{note de synthèse d'une page} destinée aux élus.
\begin{enumerate}
\item Dégages le problème philosophique posé par cette situation (opposition entre sécurité et liberté).
\item Rédige la note en respectant la structure suivante : présentation du problème (définitions des termes), exposé des deux positions possibles avec un argument chacune, position argumentée et recommandation finale.
\item Explique en cinq lignes en quoi cette démarche reprend les étapes de la dissertation philosophique et ce qu'elle y ajoute (destinataire, visée pratique, contrainte de brièveté).
\end{enumerate}
\end{exercice}

\newpage

\coursec{Corrigés des exercices}

\begin{corrige}[Exercice 1 — QCM méthodologique]
\begin{enumerate}
\item \textbf{Réponse b)} : la définition des termes du sujet et la recherche du problème. L'analyse du sujet est la première étape de toute dissertation : on ne peut ni rédiger l'introduction (a), ni choisir un plan (c), ni citer un auteur (d) avant d'avoir compris ce que le sujet demande. Définir les termes permet de dégager la tension ou la difficulté qui fait du sujet un \emph{problème} philosophique.
\item \textbf{Réponse c)} : la problématique est la question qui exprime le problème philosophique soulevé par le sujet. Ce n'est ni une question quelconque (a) — elle doit naître de l'analyse des termes — ni une simple reformulation (b) — elle doit \emph{approfondir} le sujet en explicitant l'enjeu — ni un titre de partie (d).
\item \textbf{Réponse b)} : dans un plan dialectique (thèse — antithèse — synthèse), la troisième partie dépasse l'opposition entre les deux premières en montrant ce que chacune a de vrai et en proposant une position plus riche. Répéter la première partie (a) serait une erreur de logique ; les exemples (c) et les biographies (d) n'ont pas leur place comme fonction d'une partie entière.
\item \textbf{Réponse b)} : la transition fait le bilan de la partie qui s'achève, montre sa \emph{limite} (ce qu'elle n'a pas résolu) et justifie ainsi le passage à la partie suivante. Elle donne à la copie sa cohérence et son mouvement. Remplir (a) est un hors-d'œuvre ; annoncer la conclusion (c) ou définir les termes (d) relève d'autres moments du devoir.
\item \textbf{Réponse b)} : la conclusion répond clairement à la problématique (c'est le bilan du devoir) et peut s'achever par une ouverture, c'est-à-dire une question nouvelle liée au sujet. Elle ne doit introduire \emph{aucun} élément nouveau (a), ni recopier l'introduction (c), ni deviner l'avis du correcteur (d).
\end{enumerate}
\begin{attention}[Points de vigilance]
Au Probatoire comme au Baccalauréat technique, le correcteur vérifie d'abord trois choses : le sujet est-il \emph{analysé} (termes définis) ? La problématique est-elle \emph{formulée} (une question précise) ? Le plan est-il \emph{annoncé} ? Ce QCM porte exactement sur ces trois attendus.
\end{attention}
\end{corrige}

\begin{corrige}[Exercice 2 — Vrai ou faux justifié]
\begin{enumerate}
\item \textbf{Faux.} L'introduction annonce le plan : on ne peut donc pas la rédiger définitivement avant de connaître ce plan. On peut certes noter des idées d'amorce au brouillon, mais la rédaction de l'introduction suppose que l'analyse du sujet et le plan détaillé sont déjà construits. Ordre correct : analyse du sujet $\rightarrow$ plan détaillé $\rightarrow$ rédaction.
\item \textbf{Faux.} Un exemple illustre, rend sensible et vérifie un argument, mais il ne prouve rien à lui seul : un cas particulier ne suffit pas à établir une vérité générale (un seul téléphone en panne ne prouve pas que la technique est mauvaise). L'exemple doit toujours être \emph{commandé} par un argument qui en dégage la portée.
\item \textbf{Vrai.} La thèse est la position défendue (par exemple : « le travail libère l'homme ») ; l'argument est la raison avancée pour la justifier (par exemple : « par le travail, l'homme transforme la nature et se reconnaît dans son œuvre »). Confondre les deux, c'est affirmer sans prouver.
\item \textbf{Vrai.} Le hors-sujet ne consiste pas seulement à parler d'autre chose : il consiste le plus souvent à traiter un problème \emph{voisin} (par exemple « le bonheur » au lieu de « le travail ») sans jamais affronter la question précise posée. C'est la faute la plus grave et la mieux sanctionnée à l'examen.
\item \textbf{Faux.} La conclusion doit répondre à la problématique et faire le bilan du devoir ; elle n'a pas à introduire d'élément nouveau, et encore moins une citation inédite. Une ouverture (question nouvelle) est possible mais non obligatoire.
\end{enumerate}
\end{corrige}

\begin{corrige}[Exercice 3 — Définitions]
\begin{itemize}
\item \textbf{Problématique} : question précise, née de l'analyse des termes du sujet, qui formule le problème philosophique à résoudre. Elle oriente tout le devoir : chaque partie du développement doit contribuer à y répondre.
\item \textbf{Plan dialectique} : organisation du développement en trois moments — une thèse (une position et ses raisons), une antithèse (la position opposée et ses raisons), puis une synthèse qui dépasse l'opposition en montrant ce que chaque moment contient de vrai.
\item \textbf{Transition} : passage rédigé entre deux parties, qui rappelle le résultat acquis, en montre la limite et pose la question qui justifie la partie suivante. Elle assure la continuité logique de la réflexion.
\item \textbf{Amorce (ou accroche)} : première phrase de l'introduction, qui attire l'attention du lecteur en partant d'un fait, d'une situation concrète ou d'un constat général menant naturellement au sujet.
\item \textbf{Argumentation} : ensemble ordonné des raisons (arguments), des références et des exemples par lesquels on justifie une thèse et on convainc le lecteur de sa validité.
\end{itemize}
\end{corrige}

\begin{corrige}[Exercice 4 — Reconnaître les étapes d'une introduction]
Ordre logique : \textbf{c) $\rightarrow$ d) $\rightarrow$ e) $\rightarrow$ a) $\rightarrow$ b)}.
\begin{enumerate}
\item \textbf{c) Amorce} : « Au quotidien, le téléphone portable transforme la vie des jeunes Camerounais, de Yaoundé à Maroua. » — on part d'un fait concret et proche du lecteur.
\item \textbf{d) Définitions des termes} : « La technique désigne l'ensemble des procédés et outils\ldots{} la liberté est le pouvoir d'agir selon sa propre volonté. » — on fixe le sens des deux notions clés du sujet.
\item \textbf{e) Mise en évidence du problème} : « Or, si la technique supprime les contraintes naturelles, elle crée aussi des dépendances\ldots » — le « Or » signale la tension entre les deux termes définis.
\item \textbf{a) Problématique} : « La technique nous rend-elle réellement libres, ou nous enchaîne-t-elle à de nouvelles dépendances ? » — la question qui dirige le devoir.
\item \textbf{b) Annonce du plan} : « Nous traiterons d'abord\ldots{} puis\ldots{} enfin\ldots » — les trois parties sont annoncées dans l'ordre.
\end{enumerate}
\begin{methode}[Test de vérification]
Une introduction est complète si l'on peut y repérer, dans cet ordre : une amorce, les définitions, le problème, la problématique (une question), l'annonce du plan. Si l'un de ces éléments manque, l'introduction est incomplète.
\end{methode}
\end{corrige}

\begin{corrige}[Exercice 5 — Pièges à éviter : cinq introductions fautives]
\begin{enumerate}
\item \textbf{Fautes : style oral et absence de démarche philosophique.} « Moi je pense », « je vais vous expliquer ça » relèvent de la conversation, non de la dissertation ; l'opinion personnelle et le proverbe remplacent l'analyse. \emph{Correction} : « Dans de nombreuses communautés, comme à Bafoussam, l'oisiveté est condamnée au nom de la valeur du travail. Mais qu'est-ce que le travail : une simple obligation sociale ou une activité aliénante ? »
\item \textbf{Fautes : absence de définitions, de problème et de problématique ; annonce vide.} « Question très intéressante », « les philosophes en ont beaucoup parlé » ne disent rien : aucun terme n'est défini, aucun problème n'est dégagé. \emph{Correction} : définir « travail » (activité de transformation de la nature) et « contrainte » (obligation qui limite la liberté), puis montrer la tension : le travail est nécessaire à la survie, mais est-il pour autant une pure contrainte ?
\item \textbf{Faute : hors-sujet.} L'élève traite du \emph{bonheur} et non du \emph{travail} : problème voisin, sujet abandonné. \emph{Correction} : repartir des termes du sujet ; si l'on veut évoquer le bonheur, ce ne peut être que comme enjeu secondaire (« le travail, source d'épanouissement ou de souffrance ? »), jamais comme sujet.
\item \textbf{Faute : plan invisible.} Définitions et problématique sont présentes, mais le développement commence sans annonce du plan : le lecteur ne sait pas où va la copie. \emph{Correction} : ajouter, après la problématique : « Nous verrons d'abord que le travail s'impose comme une nécessité subie, puis qu'il peut être un moyen d'accomplissement, enfin qu'il dépend des conditions dans lesquelles il s'exerce. »
\item \textbf{Fautes : amorce creuse et annonce de plan maladroite.} « Depuis la nuit des temps » est un cliché vide ; « nous verrons que oui, ensuite que non, ensuite les deux » annonce un plan dialectique de façon scolaire et contradictoire. \emph{Correction} : « Nous examinerons d'abord le travail comme contrainte imposée par la nécessité, puis comme activité par laquelle l'homme se libère, enfin les conditions qui font basculer le travail de la servitude à l'épanouissement. »
\end{enumerate}
\begin{attention}[Point de vigilance]
À l'examen, ces cinq fautes sont repérées dès les premières lignes et orientent immédiatement la notation. Relecture obligatoire de l'introduction : ai-je évité le « je » oral, le cliché, le hors-sujet, l'oubli des définitions et l'annonce du plan ?
\end{attention}
\end{corrige}

\begin{corrige}[Exercice 6 — Banque d'exemples camerounais]
Voici une proposition de corrigé (tout exemple précis, daté ou localisé, est acceptable) :
\begin{center}
\begin{tabularx}{\linewidth}{|l|X|}
\hline
\rowcolor{popblueL} \textbf{Notion} & \textbf{Trois exemples camerounais} \\
\hline
Technique & Le barrage hydroélectrique de Song Loulou sur la Sanaga ; le paiement par mobile money (MTN Mobile Money, Orange Money) ; les ateliers de fabrication de marmites en aluminium recyclé à Douala. \\
\hline
Travail & Le secteur informel (bayam-sellam, motos-taxis « benskin ») ; la culture du cacao dans le Centre et le Sud-Ouest ; les chantiers de construction du stade d'Olembe. \\
\hline
État & L'Assemblée nationale et la Constitution de 1996 ; la gratuité des soins aux enfants de moins de cinq ans dans les formations sanitaires publiques ; la décentralisation et les conseils municipaux. \\
\hline
Nature & Le parc national de Waza et la protection des éléphants ; la montée des eaux et les inondations à Douala ; la déforestation du bassin du Congo. \\
\hline
Culture & Les danses traditionnelles (bikutsi, assiko) ; le festival Ngondo des Sawa ; la diversité linguistique (plus de 250 langues nationales). \\
\hline
Justice & Les tribunaux de première instance et la Cour suprême ; la justice traditionnelle rendue par les chefs coutumiers ; la lutte contre la corruption (CONAC). \\
\hline
Liberté & La liberté de la presse et ses limites ; la libre circulation sur l'axe Yaoundé--Douala ; le débat sur la liberté d'expression des jeunes sur les réseaux sociaux. \\
\hline
Vérité & La rumeur contre l'information vérifiée lors des campagnes de vaccination ; les commissions d'enquête ; le proverbe comme parole de vérité dans les traditions orales. \\
\hline
Art & Les statuettes et masques du Grassfield ; la musique de Manu Dibango ; la sculpture sur bois des artisans de Foumban. \\
\hline
Religion & La coexistence des cultes chrétien, musulman et traditionnels ; les pèlerinages marials (sanctuaire de Mvolyé à Yaoundé) ; les rites funéraires traditionnels. \\
\hline
\end{tabularx}
\end{center}
\begin{attention}[Point de vigilance]
Un bon exemple est \emph{précis} (lieu, institution, œuvre identifiable) et \emph{pertinent} (il illustre exactement l'argument qu'il accompagne). Un exemple vague (« au Cameroun, les gens travaillent ») ne vaut rien dans une copie.
\end{attention}
\end{corrige}

\begin{corrige}[Exercice 7 — Quinze minutes chrono]
\textbf{Sujet 1 : La culture nous détache-t-elle de la nature ?}
\emph{Définitions} : la culture est l'ensemble des œuvres, savoirs et coutumes par lesquels l'homme transforme son milieu ; la nature est ce qui existe sans l'homme et obéit à des lois nécessaires. \emph{Problématique} : la culture nous arrache-t-elle à notre condition naturelle, ou restons-nous des êtres naturels jusque dans nos œuvres ? \emph{Plan} : I. La culture oppose l'homme à la nature (technique, agriculture, vie en ville). II. L'homme reste un être naturel (besoins, corps, maladie, mort). III. La culture est une manière humaine d'habiter la nature, non de la quitter.

\textbf{Sujet 2 : La liberté consiste-t-elle à faire ce que l'on veut ?}
\emph{Définitions} : la liberté est le pouvoir d'agir par soi-même ; « faire ce que l'on veut » désigne la satisfaction immédiate de tout désir. \emph{Problématique} : obéir à tous ses désirs est-il être libre, ou est-ce au contraire en devenir l'esclave ? \emph{Plan} : I. La liberté semble être l'absence d'obstacles. II. Mais le désir illimité est une servitude (l'addiction au téléphone, la dépense impulsive). III. La vraie liberté est maîtrise de soi et obéissance à la loi qu'on s'est donnée.

\textbf{Sujet 3 : L'art a-t-il pour fonction d'être utile ?}
\emph{Définitions} : l'art est la production d'œuvres belles ou expressives ; l'utile est ce qui sert un besoin pratique. \emph{Problématique} : la valeur de l'art se mesure-t-elle à son utilité, ou l'art a-t-il une fin propre ? \emph{Plan} : I. L'art peut être utile (masques rituels, chansons engagées, artisanat). II. Mais l'œuvre d'art vaut aussi par sa beauté gratuite (l'art pour l'art). III. L'utilité de l'art est d'un autre ordre : il élève, éduque, fait penser.

\textbf{Sujet 4 : La religion est-elle une affaire privée ?}
\emph{Définitions} : la religion est un ensemble de croyances et de pratiques liées au sacré ; le privé est ce qui relève de la conscience individuelle, par opposition à l'espace public. \emph{Problématique} : la foi relève-t-elle seulement de la conscience, ou a-t-elle une dimension sociale irréductible ? \emph{Plan} : I. La foi est d'abord un acte intérieur et libre (liberté de conscience). II. Mais la religion est aussi une pratique collective et publique (cultes, fêtes, rites funéraires). III. La laïcité et la tolérance organisent la coexistence du privé et du public.

\textbf{Sujet 5 : La vérité est-elle relative à chacun ?}
\emph{Définitions} : la vérité est l'accord de la pensée avec la réalité ; « relative à chacun » signifie qu'elle dépendrait du point de vue de chacun. \emph{Problématique} : chacun a-t-il sa vérité, ou la vérité est-elle une et commune à tous ? \emph{Plan} : I. Les opinions et les cultures diffèrent : la vérité semble relative. II. Mais la science et la logique exigent une vérité universelle (2 + 2 = 4 partout). III. Il faut distinguer l'opinion, relative, de la vérité, qui s'impose à tous.
\end{corrige}

\begin{corrige}[Exercice 8 — Construire une transition]
\emph{Corrigé proposé} :
« Ainsi, l'État apparaît bien comme le garant des conditions matérielles sans lesquelles aucun bonheur n'est possible : en assurant la sécurité, en construisant des écoles et des hôpitaux, il délivre le citoyen de la misère et de la peur. Mais suffit-il de ne plus souffrir pour être heureux ? Le bien-être matériel n'est que le sol sur lequel le bonheur peut pousser ; il n'est pas le bonheur lui-même. Un citoyen logé, soigné et en sécurité peut rester profondément malheureux. Dès lors, ne faut-il pas reconnaître que le bonheur relève d'un choix personnel, d'une manière de vivre et de donner un sens à son existence, qui échappe par nature à toute décision de l'État ? C'est ce qu'il nous faut maintenant examiner. »
\begin{methode}[Structure vérifiée]
La transition respecte les trois temps exigés : (1) rappel du résultat acquis (« l'État garantit les conditions matérielles du bien-être ») ; (2) mise en évidence de sa limite (« le bien-être n'est pas le bonheur ») ; (3) question ouvrant la partie suivante (« le bonheur n'échappe-t-il pas à l'État ? »).
\end{methode}
\end{corrige}

\begin{corrige}[Exercice 9 — Sujet type Probatoire : analyse complète]
\textbf{1. Définitions des termes.}
\emph{État} : organisation politique souveraine qui exerce le pouvoir sur un territoire et une population, par des institutions (gouvernement, lois, tribunaux). \emph{Devoir} : obligation morale ou juridique, ce qu'on est tenu de faire. \emph{Garantir} : assurer de façon certaine, prendre en charge la réalisation d'un bien. \emph{Bonheur} : état de satisfaction durable, accomplissement de l'existence. \emph{Citoyen} : membre d'un État, titulaire de droits et soumis à des devoirs.

\textbf{2. Problème et problématique.}
\emph{Problème} : le bonheur semble relever de la vie intérieure et des choix personnels ; peut-on alors en faire une charge publique, un devoir de l'État ? \emph{Problématique} : « L'État peut-il et doit-il rendre ses citoyens heureux, ou le bonheur échappe-t-il par nature à l'action politique ? »

\textbf{3. Plan détaillé.}
\begin{itemize}
\item \textbf{I. Thèse} : l'État a le devoir de garantir les conditions du bonheur. \emph{Arguments} : (a) sans sécurité ni subsistance, nul ne peut être heureux — or assurer la sécurité et le bien-être est la raison d'être de l'État ; (b) la justice sociale exige que l'État corrige les inégalités qui font obstacle au bonheur de tous. \emph{Exemple} : la gratuité des soins pour les enfants de moins de cinq ans et les programmes de lutte contre la pauvreté au Cameroun.
\item \textbf{II. Antithèse} : l'État ne peut pas garantir le bonheur. \emph{Arguments} : (a) le bonheur est un sentiment subjectif : ce qui rend l'un heureux laisse l'autre indifférent ; (b) un État qui prétendrait fabriquer le bonheur imposerait un modèle unique et menacerait la liberté. \emph{Exemple} : le rôle des coutumes locales et des solidarités communautaires (tontines, associations villageoises) dans le bonheur des personnes, en dehors de toute action de l'État.
\item \textbf{III. Synthèse} : l'État doit garantir les \emph{conditions} du bonheur, non le bonheur lui-même. \emph{Arguments} : (a) il faut distinguer le bien-être matériel, qui peut être organisé, du bonheur, qui se vit ; (b) l'État juste crée un cadre (droits, services publics) à l'intérieur duquel chacun construit librement son bonheur. \emph{Exemple} : l'école publique offre à chaque enfant camerounais les moyens de choisir sa vie, sans lui imposer ce qu'il doit désirer.
\end{itemize}

\textbf{4. Introduction rédigée (modèle).}
« Chaque jour, au Cameroun, les citoyens attendent de l'État des routes praticables, des hôpitaux équipés, des écoles ouvertes : comme si le pouvoir public était responsable de leur bien-être. Mais cette attente est-elle légitime ? Par État, on entend l'organisation politique souveraine qui gouverne un territoire par des lois et des institutions ; le devoir désigne une obligation morale ou juridique ; garantir signifie assurer de façon certaine ; le bonheur est un état durable d'accomplissement de l'existence ; enfin, le citoyen est le membre d'un État, titulaire de droits et de devoirs. Or, une difficulté apparaît aussitôt : si l'État peut construire des écoles, peut-il pour autant fabriquer le bonheur, qui semble relever de la vie intérieure et des choix de chacun ? Dès lors, l'État a-t-il le devoir de garantir le bonheur de ses citoyens, ou cette tâche excède-t-elle par nature sa compétence ? Nous verrons d'abord que l'État doit garantir les conditions matérielles du bien-être, puis que le bonheur, choix personnel, échappe à son action, enfin qu'il faut distinguer le devoir de créer les conditions du bonheur et la prétention impossible de le produire. »

\textbf{5. Conclusion rédigée (modèle).}
« Au terme de cette réflexion, nous pouvons répondre à notre problématique : l'État n'a pas le devoir de garantir le bonheur lui-même, car le bonheur est une manière personnelle de vivre qu'aucune loi ne peut imposer. Mais il a le devoir strict d'en garantir les conditions : la sécurité, la santé, l'éducation, la justice. Un État qui néglige ces conditions trahit sa mission ; un État qui prétendrait imposer le bonheur trahirait la liberté. Reste alors une question : si le bonheur dépend en dernier ressort de chacun, que devons-nous faire, nous-mêmes, pour être heureux ? »

\begin{attention}[Barème indicatif (sur 20)]
Définitions des cinq termes : 4 pts ; problème et problématique : 3 pts ; plan détaillé cohérent (thèses, arguments, exemples) : 5 pts ; introduction complète (5 étapes) : 4 pts ; conclusion (réponse + ouverture) : 2 pts ; clarté de l'expression : 2 pts. \emph{Points de vigilance} : ne pas confondre « garantir le bonheur » et « garantir les conditions du bonheur » — toute la synthèse repose sur cette distinction ; chaque exemple camerounais doit être précis et rattaché à un argument.
\end{attention}
\end{corrige}

\begin{corrige}[Exercice 10 — Sujet type Bac : dissertation complète]
\textbf{1. Analyse du sujet et plan détaillé (brouillon, 30 min).}
\emph{Définitions} : le travail technique est l'activité de transformation de la matière au moyen d'outils, de machines et de procédés organisés (atelier, chantier, usine) ; l'aliénation est la perte de soi, l'état où l'homme devient étranger à son activité, à son produit et à lui-même. \emph{Problème} : la technique est censée soulager l'homme ; pourtant le travail technique, répétitif et commandé, semble parfois détruire celui qui l'accomplit. \emph{Problématique} : « Le travail technique libère-t-il l'homme ou le dépossède-t-il de lui-même ? »
\emph{Plan dialectique} : I. Le travail technique peut être une aliénation (travail parcellaire, rythme imposé, produit qui échappe au travailleur). II. Mais le travail technique est aussi humanisation et maîtrise (le geste savant, la fierté de l'ouvrage, la transformation de la nature). III. L'aliénation ne tient pas à la technique elle-même mais aux conditions du travail : c'est le rapport social au travail qu'il faut transformer.

\textbf{2. Éléments de rédaction attendus.}
\emph{Introduction} : amorce concrète (l'ouvrier d'une usine de Douala répétant le même geste huit heures par jour), définitions, problème, problématique, annonce du plan. \emph{Partie I} : le travail à la chaîne fragmente le geste et prive l'ouvrier de la vue d'ensemble de son œuvre ; le produit ne lui appartient pas ; exemple : les ouvriers des usines de transformation du cacao ou les manœuvres des grands chantiers payés à la tâche. \emph{Partie II} : par le travail, l'homme transforme la nature et se transforme lui-même ; le technicien se reconnaît dans son ouvrage ; exemple : le menuisier de Foumban, le mécanicien du quartier Nkolndongo qui répare et invente, la fierté du maçon devant la maison achevée. \emph{Partie III} : la technique est un instrument neutre ; c'est l'organisation du travail (salaire, rythmes, reconnaissance, formation) qui décide de l'aliénation ou de l'épanouissement ; exemple : les coopératives artisanales où le travailleur maîtrise tout le processus. \emph{Conclusion} : réponse nuancée à la problématique (le travail technique n'est aliénant que dans certaines conditions) et ouverture (« la formation technique ne serait-elle pas la première condition d'un travail libre ? »).

\textbf{3. Grille d'auto-évaluation (sur 20).}
\begin{center}
\begin{tabularx}{\linewidth}{|l|X|c|}
\hline
\rowcolor{popblueL} \textbf{Critère} & \textbf{Attendu} & \textbf{Note /4} \\
\hline
Respect du sujet & Toutes les parties traitent du \emph{travail technique} et de l'\emph{aliénation}, sans dériver vers le travail en général. & \\
\hline
Problématique & Une question claire, née des définitions, annoncée dans l'introduction. & \\
\hline
Arguments & Chaque thèse est justifiée par au moins deux raisons distinctes et ordonnées. & \\
\hline
Exemples & Au moins deux exemples camerounais précis (atelier, chantier, usine, informel), rattachés aux arguments. & \\
\hline
Expression & Français correct, pas de style oral, transitions rédigées. & \\
\hline
\end{tabularx}
\end{center}
\begin{attention}[Points de vigilance]
Conditions d'application à citer : au moins \emph{deux références philosophiques} de l'unité (par exemple la conception du travail comme transformation de soi, et l'analyse de l'aliénation du travailleur) et \emph{deux exemples professionnels camerounais}. Erreur fréquente : traiter « la technique » en général au lieu du \emph{travail} technique — c'est un hors-sujet partiel. Les transitions doivent être rédigées, pas seulement suggérées.
\end{attention}
\end{corrige}

\begin{corrige}[Exercice 11 — Note de synthèse philosophique pour des décideurs]
\textbf{1. Le problème philosophique.}
La situation oppose deux valeurs légitimes : la \emph{sécurité} (protection des personnes et des biens, qui est une mission de la puissance publique) et la \emph{liberté} (droit d'aller et venir sans être observé, vie privée). Installer des caméras, c'est augmenter la sécurité au risque de surveiller les innocents ; les refuser, c'est protéger la liberté au risque de laisser le crime impuni. Le problème est donc : jusqu'où peut-on limiter la liberté au nom de la sécurité sans détruire ce que la sécurité était censée protéger ?

\textbf{2. Note de synthèse (modèle).}
\emph{Objet : installation de caméras de surveillance dans les quartiers — note à l'attention des élus municipaux.}

\textbf{Présentation du problème.} Par sécurité, on entend la protection effective des personnes et des biens ; par liberté, le droit de chacun de circuler et de vivre sans surveillance permanente. La vidéosurveillance met ces deux valeurs en tension : filmer pour protéger, c'est aussi observer tous les passants, y compris les innocents.

\textbf{Position 1 — Installer les caméras.} Argument : la première mission de la collectivité est de protéger ses habitants ; la dissuasion et l'identification des délinquants réduisent la délinquance, et la sécurité est la condition de toutes les autres libertés (on ne jouit d'aucune liberté quand on a peur).

\textbf{Position 2 — Refuser les caméras.} Argument : une société qui surveille tout le monde traite chaque citoyen en suspect ; la liberté, notamment la vie privée, est un bien qui, une fois perdu, ne se récupère pas, et l'habitude de la surveillance peut servir un jour d'autres fins que la sécurité.

\textbf{Position argumentée et recommandation.} La sécurité et la liberté ne s'opposent pas absolument : la sécurité protège l'exercice de la liberté, mais une sécurité totale détruirait la liberté qu'elle prétend défendre. \emph{Recommandation} : installer les caméras de façon limitée et encadrée — uniquement aux points sensibles (marchés, gares), pour une durée de conservation des images courte, avec une information claire des habitants et un contrôle citoyen de leur usage. Ainsi la collectivité remplit son devoir de protection sans renoncer au respect des libertés.

\textbf{3. Lien avec la dissertation.}
Cette note reprend exactement les étapes de la dissertation : définition des termes, mise en évidence du problème, examen d'une thèse puis de la thèse opposée avec un argument chacune, position de synthèse justifiée. Elle y ajoute trois contraintes propres : un \emph{destinataire} déterminé (les élus, non un correcteur), une \emph{visée pratique} (décider, non seulement comprendre), et une \emph{contrainte de brièveté} (une page, un style direct). La philosophie sort ainsi de la copie d'examen pour éclairer une décision citoyenne réelle.
\end{corrige}

\newpage

\coursec{Fiche bilan}

\begin{popfigure}
\begin{tikzpicture}[mindmap, grow cyclic, every node/.style={concept, font=\small},
root concept/.append style={concept color=popblue, font=\bfseries\small, minimum size=2.6cm},
level 1 concept/.append style={level distance=3.6cm, sibling angle=60, font=\scriptsize, minimum size=1.9cm}]
\node[root concept]{La dissertation philosophique}
child[concept color=poporange]{node[align=center]{Analyser le sujet\\ (termes, problématique)}}
child[concept color=popgreen]{node[align=center]{Construire le plan\\ (thèse, antithèse, synthèse)}}
child[concept color=poppurple]{node[align=center]{Introduction\\ (accroche, problème, annonce)}}
child[concept color=poppink]{node[align=center]{Conclusion\\ (bilan, ouverture)}}
child[concept color=popgold]{node[align=center]{Argumenter\\ (preuves, exemples, auteurs)}}
child[concept color=popblueL]{node[align=center]{Transfert\\ (vie citoyenne et professionnelle)}};
\end{tikzpicture}
\legende{Carte mentale de la leçon : les six compétences clés de la dissertation philosophique.}
\end{popfigure}

\begin{bilan}
\textbf{Définitions clés}
\begin{itemize}
\item \cle{Dissertation philosophique} : exercice écrit qui répond à une question par une réflexion argumentée, structurée et personnelle.
\item \cle{Problématique} : question centrale dégagée de l'analyse du sujet ; elle dirige tout le devoir.
\item \cle{Thèse} : position défendue ; \cle{antithèse} : position opposée ; \cle{synthèse} : dépassement qui réconcilie ou tranche.
\item \cle{Argument} : raison qui justifie une thèse ; \cle{exemple} : cas concret qui illustre l'argument.
\end{itemize}

\textbf{Structure à connaître par c\oe ur}
\begin{center}
\begin{tabularx}{\linewidth}{|l|X|}
\hline
\rowcolor{popblueL} \textbf{Partie} & \textbf{Contenu obligatoire} \\
\hline
Introduction & Accroche, définition des termes, problématique, annonce du plan \\
\hline
Développement & 2 ou 3 parties : thèse, antithèse, synthèse (arguments + exemples + transitions) \\
\hline
Conclusion & Bilan de la réflexion, réponse claire à la problématique, ouverture \\
\hline
\end{tabularx}
\end{center}

\textbf{Méthodes express}
\begin{enumerate}
\item Analyser le sujet : souligner les mots-clés, définir chaque terme, formuler la problématique en une question.
\item Rédiger au brouillon : arguments, exemples, références d'auteurs (Platon, Descartes, Kant, Nietzsche...).
\item Soigner les transitions entre les parties et la clarté du style.
\end{enumerate}
\end{bilan}

\begin{aretenir}[Checklist avant le Bac]
Avant de rendre ma copie, je vérifie :
\begin{itemize}
\item[$\square$] J'ai bien défini tous les termes importants du sujet.
\item[$\square$] Ma problématique est formulée clairement dans l'introduction.
\item[$\square$] Mon plan est annoncé à la fin de l'introduction.
\item[$\square$] Chaque partie défend une idée distincte (thèse, antithèse, synthèse).
\item[$\square$] Chaque argument est appuyé par un exemple concret.
\item[$\square$] J'ai cité au moins deux auteurs ou doctrines philosophiques.
\item[$\square$] Mes transitions relient logiquement les parties.
\item[$\square$] Ma conclusion répond directement à la problématique.
\item[$\square$] Je n'ai pas fait de hors-sujet.
\item[$\square$] Mon écriture est lisible et mon orthographe relue.
\end{itemize}
\end{aretenir}

\findecours

\end{document}
